% FOR FASTER COMPILING: just the section I'm editing
% \includeonly{Zchap_ST_00_Probability}

\documentclass[12pt, oneside]{report}
% \documentclass[12pt, oneside, draft]{report}

\input{FlipLectureMacros}       %% Notes formatting, load first
\input{FlipPreamble}			%% \usepackages, formatting
\input{FlipMacros}              %% Flip's standard macros
\input{FlipMacros_Comments}     %% Flip's macros for comments
\input{FlipMacros_Teaching}     %% Flip's macros for course notes
\input{Flip_listings}           %% Styling for code blocks
\input{FlipAdditionalHeader}    %% Modify this for each project
\input{FlipPreambleEnd}         %% has to be loaded last

%%%%%%%%%%%%%%%%%%%%%%%%%%%%
%% LECTURE NOTES SETTINGS %%
%%%%%%%%%%%%%%%%%%%%%%%%%%%%

% \linenumbers                  %% lineo package
\graphicspath{{figures/}}       %% figure folder
\addbibresource{FlipBib.bib}    %% Define BibLaTeX source(s)

\geometry{                      %% large margin for side notes
    paper=letterpaper, 
    hmargin={1cm,7.25cm},       %% 6.25cm space on right
    vmargin={2cm,2cm}, 
    marginparsep=.5cm, 
    marginparwidth=5.75cm
}

%% Def. full width; uses changepage package
%%  6.25cm to match hmargin difference;
\newenvironment{wide}%
{\begin{adjustwidth}{0cm}{-6.25cm}}%
{\end{adjustwidth}}

% Reset the sidenote number each section 
\let\oldsection\section
\def\section{%
  \setcounter{sidenote}{1}%
  \oldsection
}



\begin{document}

\newgeometry{margin=2cm}
\newcommand{\FlipTR}{UCR-TR-F2024-FLIP-P231} 
\thispagestyle{firststyle}
\pagenumbering{gobble}

\begin{center}
    {\large \textsf{UC Riverside Physics 231, Fall 2024} \par}
    {\huge \textbf{Mathematical Methods for Physicists} \par}\vspace{.5em}
    {\large {Linear algebra, complex analysis, Green's functions, distributions} \par}
    \vskip .5cm
\end{center}

\input{FlipAuthors}

\vspace{2em}\noindent
Lecture notes for Physics 231: Methods of Theoretical Physics, a course on linear algebra, complex analysis, Green's functions, and probability in preparation for graduate-level physics coursework at \acro{UC~R}iverside. The first portion of these notes encompass an introduction to linear algebra ``with indices'' for lower-division undergraduates. 

\vspace{5em}
\begin{center}
% \includegraphics[width=.2\textwidth]{figures/Squiggle.pdf}
\includegraphics[width=.4\textwidth]{figures/vevtilt.pdf}
\end{center}  


% \vspace{2em}
\vspace*{\fill}

\noindent
\textsf{Last Compiled: \today}

\noindent
\textsf{Image: vacuum misalignment}

\noindent
\textsf{CC BY-NC-SA 4.0}~\ccbyncsa 

\noindent % Course notes URL
% \url{https://github.com/fliptanedo/P231-2023-Math-Methods}

%% Front page logos
\vspace*{\fill}
\begin{center}
\includegraphics[height=.1\textwidth]{figures/FlipAmbigram.png}
\hspace{5em}
\includegraphics[height=.1\textwidth]{figures/UCRPnA_banner.png}
\end{center}

\newpage

\small
\setcounter{tocdepth}{2}
\tableofcontents
\normalsize
\clearpage
\restoregeometry        %% Return to lecture note geometry 
\pagenumbering{arabic}  %% Turn on regular page numbers


%%%%%%%%%%%%%%%%%%%%%
%%%  THE CONTENT  %%%
%%%%%%%%%%%%%%%%%%%%%



\include{Zchap_IntroPhD}
\include{Zchap_DimAnalysis}

\part{Linear Algebra}

This part is a review of undergraduate-level linear algebra and forms the core of the Physics 17 course. The material is likely to be very familiar for graduate students in physics, but we provide it here to emphasize the use of indices and the `unified picture' of vector spaces and their duals that appear in the study of Green's functions. Physics 231 students can start with Chapter~\eqref{ch:P231:summary} and refer back as necessary.

\include{Zchap_LA_00_Basics}
\include{Zchap_LA_01_Indexology}
\include{Zchap_LA_02_vectors}
\include{Zchap_LA_03_bundle}
\include{Zchap_LA_04_metric}
\include{Zchap_LA_05_relativity}
\chapter{Linear Transformations}

In this chapter we focus on linear transformations in metric spaces. Recall that the term \emph{linear transformation} means a linear function that takes in one type of object and returns the same type of object. The linearity means that a linear combination of inputs returns a linear combination of outputs. In this course, we reserve the word \textbf{matrix}\index{linear transformation} to mean a linear transformation. We observed that a matrix $M\aij{i}{j}\ket{e_i}\!\bra{e_j}$ can equivalently take in a vector or a row vector and output an object of the same class. 


\section{The hierarchy of transformations}
\label{sec:hierarchy:of:transformations}

There is something odd that I do when I think about mathematical objects. I impose upon these objects some subjective feeling of niceness or not-niceness. Nice matrices are easy to use. The nicest matrices are easy to undo. Matrices that are not nice may be complicated---perhaps they are a complicated manifestation of something which is supposed to be simple. Another not-nice feature is if a matrix is not invertible: once you apply it, the object on which you have applied it has lost information. The worst matrices do not fall into these categories. There is one step below `worst': matrices that are not even matrices by our definition, for example linear maps that do not return vectors in the same vector space as the input.\sidenote{There are even worse (worser?) ``matrices'' like the $N\times (N+1)$ objects you write down to do weird shit like ``reduced row echelon form'' in a traditional linear algebra class. That's just solving a system of linear equations---that's just algebra, it is not even \emph{linear algebra}.}

The following hierarchy of matrices is not official. It is not found in any textbook of mathematics or physics. In fact, I would vehemently deny that I even taught this to you.\sidenote{This reminds me of the advice my Ph.D advisor gave me when I asked what I should do if my first seminar talk does not go well. He told me, ``If it goes \emph{really} poorly, then lie about who your advisor is.''} However, it does help me categorize the types of matrices that \emph{do} show up in physics and how to think about them.\sidenote{You may compare this to what psychologist refer to as Maslow's hierarchy of needs.}

\begin{enumerate}
    \item The top-of-the line, S-rank, nicest matrices are \textbf{multiples of the identity}. These are just numbers: not only are they defined by a single proportionality constant, but they add and multiply as if they are just numbers. Unfortunately, these matrices are \emph{too} nice: all they do is rescale the vector or dual vector that they act on. That is a little boring, even for us.
    \item The A-class matrices are \textbf{diagonal} in the basis you are currently using. We also assume that each of the diagonal elements are non-zero.\sidenote{We say the matrix is \emph{non-degenerate}.} In this basis, the matrix acts by rescaling each component of a vector or row vector. Each component is stretched or shrunk independently: that means that the resulting vector could look rather different. Most of the time we specialize to the case where the diagonal elements are positive.
    \item At \emph{almost} the same level as the diagonal matrices are matrices that are \textbf{diagonalizable} and, when diagonalized, have no zero elements along the diagonal. The verb `to diagonalize' means that there is some \emph{isometry} (rotation) where we may change basis and the components of $A$ would be diagonal. Most of our active transformations in physics---the ones where we ``do something''---take this form. Because diagonal matrices are so nice, it then behooves us to take these matrices and rotate into the proper basis where they are diagonal.\sidenote{This is called the eigenbasis or basis of eigenvectors. The word `eigen' means `proper' in German. At least I think that's what it means. The diagonal components of the matrix in the eigenbasis are called eigenvalues.}
    \item Now I pause to (re-)define a different class of matrices, \textbf{isometries}. Isometries are transformations that preserve the form of the metric. In some sense, they should not even be on this list since they are most naturally understood as \emph{passive} transformations---see Section~\ref{sec:active:passive}. They are a special class because we use these isometries to go into a basis where matrices are diagonal.\sidenote{In my brain diagonalizing matrices is ``doing [the] eigenstuff.''} Even among isometries, there are two types depending on whether or not they preserve the orientation of your basis.\sidenote{This is a matter of whether or not you change the \emph{handedness} of your basis.}
    \item There are matrices that are not so `nice' but are often meaningful and useful. These are called \textbf{projections}. These are matrices that are not invertible because they remove information from the vectors that they act on. A good example is a diagonal matrix with at least one zero diagonal component. When that matrix acts on a vector, the information of that component is completely lost---there is no way to recover it based on the output vector. 
    \item The last useful type of matrix are those that are none of the above but are otherwise $N\times N$ matrices that take vectors from one vector space and return vectors in the same vector space. These may be random $N\times N$ matrices, for example. We neither expect them to be diagonalizable by a rotation nor are they obviously meaningful as projections. These matrices may still be diagonalized, but usually the transformation to do this is not a simple isometry. Further, the diagonal matrices may not necessarily be positive upon diagonalization.
    \item All other matrices---for example, $N\times M$ matries with $M\neq N$---are just completely off the table. 
\end{enumerate}

\begin{example}[Active versus passive] In a passive transformation, tensors do not transform. Instead, we are simply changing our basis---physically, we are changing our reference frame. For such a transformation, the components of a tensor change because the basis vectors change, even if the tensor itself remains the same. 

In an active transformation, the components of a tensor change but the basis remains the same. This corresponds to actually changing the tensor and not changing the observer.
\end{example}

\begin{exercise}
Draw a vector---any vector---in $\mathbbm{R}^2$. Act on this vector with the matrix, whose components we write with respect to the standard basis:
\begin{align}
    A = \begin{pmatrix}
        A\aij{1}{1} & A\aij{1}{2}\\
        A\aij{2}{1} & A\aij{2}{2}
    \end{pmatrix}
    =
    \begin{pmatrix}
        3 & 0\\
        0 & \frac{1}{2}
    \end{pmatrix} \ .
\end{align}
Draw the resulting vector. Do the same for the case with the same matrix except that $A\aij{1}{1} = -3$.
\end{exercise}


% Identity
% Diagonal
% Rotated from the diagonal: symmetric
% Projections
% All others
% not included: $N\times M$

\section{Non-invertible matrices}

The nicest matrices are invertible. In fact, most of the linear transformations you encounter in a physics curriculum are invertible simply because so many of the mathematical tasks in physics involve inverting the dynamics encoded in these matrices. It behooves us to say a few words about the types of matrices that are not so nice.

\subsection{A first look at projections}

Projections are the matrices that are non-invertible because they send some non-zero vector to zero. A simple example follows:
\begin{example}
Consider the matrix
\begin{align}
    M = 
    \begin{pmatrix}
        2 & 0 \\
        0 & 0
    \end{pmatrix} \ .
\end{align}
This matrix takes any vector oriented in the $\hat{e_2}$ direction to zero: $M\ket{e_2} = 0$. This means that the action of the matrix on a general vector $\ket{v} = v^i\ket{e_i}$ is
\begin{align}
    M \ket{v} = v^1 M\ket{e_1} = 2 v^1 \ket{e_1} \ .
\end{align}
It does not matter what $v^2$ is, the image of $M\ket{v}$ only depends on $v^1$. In this way, one has \emph{lost} information. If you know that
\begin{align}
    M\ket{v} = \ket{w}
    \label{eq:Mv:eq:w:projection:eg}
\end{align}
for some known vector $w^i\ket{e_i}$, then there is no way to figure out what $\ket{v}$ is uniquely because there is a degeneracy where any vector $\ket{v}$ satisfying \eqref{eq:Mv:eq:w:projection:eg}, any vector of the form $\ket{v} + \alpha\ket{e_2}$ will also satisfy \eqref{eq:Mv:eq:w:projection:eg}. 

In fact, you may notice that if $w^2 \neq 0$, then \eqref{eq:Mv:eq:w:projection:eg} has no solution.
\end{example}

\begin{exercise}
A slightly less trivial example is as follows:
\begin{align}
    M = 
    \begin{pmatrix}
        2 & 1 & 0 \\
        3 & 2 & 0 \\
        0 & 0 & 0
    \end{pmatrix} \ .
\end{align}
Verify that the image of any vector $\vec{v}$ does not contain any component in the $\ket{e_3}$ direction. Further verify that the image of any any vector $\vec{v}$ is independent of $v^3$.
\end{exercise}

Matrices of these types are projections. We discuss these more fully in Section~\ref{sec:projections}. When a matrix causes a non-zero vector to vanish, it is necessarily a projection.\footnote{A narrow definition of a projection is that most components of a vector are preserved while some are lost. An example of this is the matrix $\ket{e_2}\!\bra{e_2}$. However, for our purposes we use projection to mean any transformation that loses information.} A key theme is that one \emph{loses information} about the original vector $\vec{v}$ when acting on it with a projection. There is no way to unambiguously reconstruct $\ket{v}$ given its image $M\ket{v}$.



\subsection{Degenerate image}

You may think that the problem with the matrices in the example above is that there are too many zeros in the matrix. After all, these zeros throw away the dependence on components of the input vector. This is na\"ive: the numerical values of a matrix are basis-dependent. Here's another example:
% 
\begin{example}
The following matrix is degenerate:
\begin{align}
    M = 
    \begin{pmatrix}
        1 & 1 \\
        2 & 2
    \end{pmatrix} \ .
\end{align}
This matrix has no zero elements---at least not in this basis---but there's clearly something fishy about it. If you act on this matrix on any vector, you find that the result is proportional to $\ket{e_1} + 2 \ket{e_2}$. Go ahead, try it. 

One way to see this is that it takes both basis vectors $\ket{e_{1,2}}$ to
\begin{align}
    M\ket{e_{1,2}} = \ket{e_1} + 2 \ket{e_2} \ .
\end{align}
We see that $M$ is again non-invertible. Given some vector of the requisite form, $\ket{w} = \alpha \ket{e_1} + 2\alpha \ket{e_2}$, there is no unambiguous solution for $\ket{v}$ to the equation $M\ket{v} = \ket{w}$. We can show this simply by writing out two solutions that both gives the same $\ket{w}$:
\begin{align}
    \ket{v_I} &= \alpha \ket{e_1}
    &
    \ket{v_{II}} &= \alpha \ket{e_2} \ .
\end{align}
\end{example}
\begin{exercise}
In the previous example, show that the most general solution to $M\ket{v} = \ket{w}$ is
\begin{align}
    \ket{v} &= \beta\ket{e_1} + (\alpha-\beta)\ket{e_2} \ .
\end{align}
\end{exercise}
\begin{exercise}
In the previous example, argue that the matrix $M$ has many zeros in the basis
\begin{align}
    \ket{f_1}&= \frac{1}{\sqrt{2}}
    \left[\ket{e_1}+\ket{e_2}\right]
    &
    \ket{f_2}&= \frac{1}{\sqrt{2}}
    \left[\ket{e_1}-\ket{e_2}\right] \ .
\end{align}
\end{exercise}
The lesson from this section is that projections are not always so obvious by their zero components. 


\section{Action on basis vectors}

A linear transformation is defined by its action on basis vectors. Consider a matrix $M$ that is non-degenerate. A convenient way of stating this is that the action of $M$ onto each of the basis vectors $M\ket{e_i}$ produces a set of vectors that are \emph{linearly independent}.\sidenote{Unless $M$ is an isometry, the vectors $M\ket{e_i}$ are not orthonormal.} This is another way of saying that $M$ is invertible and so does not lose information. 
\begin{exercise}
In your own words, articulate why the linear independence of the $M\ket{e_i}$ image vectors implies that $M$ is invertible? \textsc{Hint}: one way to do this is to argue that a given vector $\ket{w}$ is the image of a unique vector $\ket{v}$ under $M$: $M\ket{v} = \ket{w}$.
\end{exercise}

This leads us to a useful point:
\begin{bigidea}
A linear transformation is defined---and completely described---by its action on the basis vectors of a vector space.
\end{bigidea}
We can see this by writing a vector as a linear combination of basis vectors, $\ket{v} = v^i\ket{e_i}$. Because $M$ is a linear transformation, knowing the action of $M$ on each $\ket{e_i}$ completely specifies the action of $M$ on the linear combination $\ket{v}$. 

What does it mean to know the action of $M$ on the basis vectors? $M$ maps each basis vector into some other vector. In general the image vector will not be normalized nor orthogonal to the images of other basis vectors. How does one describe the image vector? The obvious way is to give the components of $M\ket{e_i}$ with respect to the $\ket{e_i}$ basis. The component of $M\ket{e_i}$ in the $\ket{e^j}$ direction is simply
\begin{align}
    \la e^j \mid M \mid e_i \ra 
    = 
    \la e^j \mid M\aij{k}{\ell}\ket{e_k}\!\bra{e^\ell} \mid e_i \ra 
    = M\aij{j}{i} \ .
\end{align}
This is one of those \emph{aha} moments where you may first feel surprised that the components of $M\aij{j}{i}$ have popped out, but then you realize that this is \emph{precisely} what the components of a matrix \emph{mean}. 

Indeed, the `ket-bra' basis $\ket{e_k}\!\bra{e^\ell}$ is doing the following: \emph{give me the component in the $\ket{e_\ell}$ direction so that I can send it into the $\ket{e_k}$ direction}. The matrix components $M\aij{k}{\ell}$ instruct the machine $\ket{e_k}\!\bra{e^\ell}$ \emph{how much} of the component in the $\ket{e_\ell}$ direction should be sent into the $\ket{e_k}$ direction.

We thus have the following interpretation of the components of a matrix:
\begin{bigidea}[Components of a Matrix]
The $M\aij{i}{j}$ component of a matrix tells you the components of the image of $\ket{e_j}$ under $M$:
\begin{align}
    M\aij{i}{j} = \la e^j \mid M \mid e_i \ra \ .
\end{align}
This means that the $j^\textnormal{th}$ columns of the matrix $M$ are the components of $M\ket{e_j}$ in the $\ket{e_i}$ basis.
\end{bigidea}




\section{Isometries}
\label{sec:isometries}

We gave a formal definition of isometries in Section~\ref{sec:isometry:next:pass:bases} and put them to use in Section~\ref{sec:spacetime:isometry} for special relativity. It is now worth returning to isometries to highlight their special role in a metric space as symmetries.

An \textbf{isometry}\index{isometry} is a linear transformation that leaves the components of the metric unchanged.\sidenote{These are generalizations of rotations, but I often use the words interchangeably. I try to speak plainly when the loss of precision of language does not detract from the content of the message. Otherwise I am reminded of a friend from the midwest who \emph{sounds} like a midwesterner except for the unexpected moments when he is ordering a \emph{kwah-sohn} from dunkin' donuts and I remember that he's actually from Quebec.} When we harp on the virtues of index notation and tensors in physics, it is because the tensor indices tell us how objects transform \emph{with respect to isometries}. 

We rigorously introduced the idea that indices tell us how objects transform in Section~\ref{sec:transformation:under:symmetries}. In that section, we \emph{said} that the transformations $R$ are rotations. Now we must come clean: in that section we only had a vector space. Without a metric, there is no way to define a rotation. In fact, all of the results in Section~\ref{sec:transformation:under:symmetries} for ``rotations'' $R$ are valid for \emph{any} invertible transformation.
\begin{exercise}
Go back to Section~\ref{sec:transformation:under:symmetries} and check that all of the transformation rules are valid for the case where $R$ any invertible transformation. The conceptual picture is clear: our rule was that the basis vectors transformed with $R$ and the components transformed with $R\inv$---or vice versa depending on the height of the index. Because $RR\inv = \one$ for any matrix $R$ with an inverse, all of the results in that section hold.
\end{exercise}

What is different when we have a metric? The inner product allows us to define length and angle. As we see in the previous section, the action of a transformation is completely encoded in its action on the basis vectors. Rotations can thus be defined as those transformations which preserve lengths and angles. In so doing, rotations are those transformations that take orthonormal bases to other orthonormal bases. 

% By the way ,you ca YOUC AN LOOK AT EFFECT OF NON ORTHONORMAL ...


% active vs passive: 
% passive acts on basis, so components compensate
% active leaves basis, but changes components, actual change

\section{Projections}
\label{sec:projections}



When a matrix acts on a vector and produces another vector, you can take the output vector and \emph{infer} the initial vector. This is because the equation
\begin{align}
    M\vec{v} = \vec{w}
\end{align}
can be solved as long as you know how to compute $M\inv$. Solving $M\inv M = \one$ is has $N^2$ equations for $N^2$ unknowns---but those equations do not necessarily have a solution. A simple example is the case of a diagonal matrix with one component zero. We illustrate this in Fig.~\ref{fig:map:M:no:inv}
\begin{figure}[tb]
    \centering
    \includegraphics[width=.8\textwidth]{figures/maps_noninv.pdf}
    \caption{$M$ is a non-invertible matrix. We notice that there's no problem applying $M$ to vectors. The odd thing is that there are multiple vectors---$\vec{v}$, $\vec{w}$, and $\vec{u}$ in the picture---that all turn into the same vector, $M\vec{v} = M\vec{w}=M\vec{u}\equiv \vec{z}$. The problem shows up when we try to ``undo'' the transformation $M$ by applying $M\inv$ to $\vec{z}$. $M\inv\vec{z}$ should be a vector. Instead, we have at least three different valid options: $\vec{v}$, $\vec{w}$, and $\vec{u}$. This means that $M\inv$ is not a well defined linear map: it can send one vector $\vec{z}$ into many possible vectors. }
    \label{fig:map:M:no:inv}
\end{figure}
% 
We see the the problem is one of uniqueness. A linear map should take each vector $\vec{v}$ into a single, well-defined vector $M\vec{v}$. $M$ is suspicious because it sends multiple different vectors to the \emph{same} vector, say $M\vec{v} = M\vec{w}$. That's fine. Each of $\vec{v}$ and $\vec{w}$ are sent to \emph{a} vector. The problem shows up when we try to take the inverse. If we write $\vec{z} = M\vec{v} = M\vec{w}$, we know that
\begin{align}
    (M\inv)\vec{z} = (M\inv)M\vec{v} &= \vec{v}
    &
    (M\inv)\vec{z} = (M\inv)M\vec{w} &= \vec{w} \ .
\end{align}
However, we assumed that $\vec{v} \neq \vec{w}$. This means makes the above line inconsistent. The inverse transformation is not defined. Bummer.

Fortunately, most of the \emph{dynamics} in physics \emph{is} invertible. That does not mean that we can ignore non-invertible matrices, though. Non-invertible matrices can be understood as \textbf{projections} onto a subspace. You can see this in the example in Fig.~\ref{fig:map:M:no:inv}. The map $M$ is essentially only using the information about the second component of its input vector and mapping it to just the second component of the output vector. That means that $M$ is really just a map to a one-dimensional subspace. In this way, the two dimensional vector space $\RR ^2$ is \textbf{projected} onto a one-dimensional subspace of $\RR ^2$. In that example, the one dimensional space is defined by vectors of the form
\begin{align}
    \vec{v} = 
    \begin{pmatrix}
    0\\ v^2    
    \end{pmatrix} \ .
\end{align}

Projections are really useful in physics. Sometimes you want to take a big vector space and only focus on a subspace that satisfy certain conditions. For example, it turns out that in nature you can have two types of massless matter particles: left-handed and right-handed according to their quantum spin in the direction of their motion. One may want to project onto only the space of left-handed particles because that is easier to deal with than the entire space. Because projections necessarily ``throw out'' information, they are non-invertible.\sidenote{I think this was more intuitive to past generations of physicists. When you throw out a piece of paper, it ends up in a recycling plant and it gets torn up and is lost forever. Since the advent of computers, when you ``throw out'' a file, you can usually control/command-z and undo to bring it back.}

\begin{example}
One of the most curious puzzles in fundamental physics over the last half century is the black hole information paradox.\footnote{\url{https://www.quantamagazine.org/the-most-famous-paradox-in-physics-nears-its-end-20201029/}} One way of stating the paradox is the question of whether or not ``falling into a black hole'' is a \emph{projection} that throws out information. 
\end{example}

\subsection{Projections in Bra-Ket Notation}
Recall our resolution of the identity \eqref{eq:resolution:of:the:identity}. Each piece in this sum is a projection. For example, consider $\ket{e_2}\!\bra{e^2}$. This is a matrix that acts on a vector $\ket{v}$, pulls out the second component $v^2$, and then returns a vector in the $\ket{e_2}$ direction with length $v^2$. We can see this mathematically:
\begin{align}
    \ket{e_2}\!\bra{e^2} \; v^i\ket{e_i} \;
    &= 
    v^1\ket{e_2}\!\la{e^2}\mid {e_1}\ra + 
    v^2\ket{e_2}\!\la{e^2}\mid {e_2}\ra +
    v^2\ket{e_2}\!\la{e^2}\mid {e_2}\ra + \cdots
    \\ 
    &= v^2 \ket{e_2} \ ,
    \label{eq:projection:e2:e2}
\end{align}
where we used the defining relation $\la e^i \mid e_j \ra = \delta^i_j$\ . This means we can think of our resolution of the identity, $\ket{e_i}\!\bra{e^i}$ as a sum over projections along each axis. One can call each term in the resolution of the identity, say $\ket{e_2}\!\bra{e^2}$ a projector\sidenote{``\emph{Even though I know someday you're gonna shine on your own / I will be your projector},'' from `Protector' by Beyonc\'e Knowles-Carter in the album Cowboy Carter (2024).} onto a given axis.

\subsection{Projecting One Vector Onto Another}

Suppose you have two vectors, $\ket{v}$ and $\ket{w}$ and you want find the projection of $\ket{w}$ onto $\ket{v}$, as shown in Figure~\ref{fig:gram:projection}. 
\begin{figure}[tb]
    \centering
    \includegraphics[width=.8\textwidth]{figures/gram-projection.jpg}
    \caption{Projection of a vector $\ket{w}$ onto the axis of another vector, $\ket{v}$.}
    \label{fig:gram:projection}
\end{figure}

Our strategy is to separate the ket $\ket{w}$ into a piece that is parallel to $\ket{v}$ and a piece that is perpendicular to $\ket{v}$,
\begin{align}
    \ket{w}= \ket{w_\paral} + \ket{w_\perp} \ .
\end{align}
The parallel piece $\ket{w_\paral}$ is precisely the projection of $\ket{w}$ onto $\ket{v}$.  The defining characteristic of $\ket{w_\paral}$ is that $\la \hat{w}_\paral, \hat{v}\ra = 1$. That is: the angle between the \emph{unit vectors} is zero---they are parallel. This means that 
\begin{align}
    \ket{w_\paral} \propto \ket{\hat{v}} \ .
\end{align}
What is the proportionality? This is simply the length of $\ket{w}$ along the $\ket{\hat{v}}$ axis, $\la w, \hat v \ra$. Then the statement that $\ket{w_\paral}$ is the vector of length $\la w, \hat v \ra$ in the $\ket{\hat v}$ direction is:
\begin{align}
    \ket{w_\paral} = \la w, \hat{v} \ra \ket{\hat{v}} 
    % =  \ket{\hat{v}}\la \hat{v}, w \ra  
    \ .
\end{align}
We recall from the symmetry of the inner product and our definition of kets \eqref{eq:basis:dual:vectors:as:inner:prod} that we may further write
\begin{align}
    \la w, \hat{v} \ra = \la \hat{v}, w \ra = \la \hat{v} \mid w \ra \ .
\end{align}
What we have done here is written the length of $\ket{w_\paral}$ with respect to a \emph{bra} $\bra{\hat{v}}$ acting on $\ket{w}$. Plugging this back into our expression for $\ket{w_\paral}$, we recover a projection matrix of exactly the form in \eqref{eq:projection:e2:e2}:
\begin{align}
    \ket{w_\paral} = \ket{\hat{v}} \la{\hat{v}} | w \ra  =
    \left(\ket{\hat{v}} \bra{\hat{v}} \right)  \, \ket{w} \ .
     \ .
\end{align}
Evidently $\ket{\hat{v}} \bra{\hat{v}}$ is the matrix that takes $\ket{w}$ and returns $\ket{w_\paral}$, the component of $\ket{w}$ along $\ket{\hat{v}}$. In this language, it is somewhat \emph{obvious} that this is what $\ket{\hat{v}} \bra{\hat{v}}$ does. However, we have not written the components of $\ket{\hat{v}} \bra{\hat{v}}$ in terms of our standard basis. To do that, one simply expands each  of $\ket{\hat{v}}$ and $\bra{\hat{v}}$ in the standard basis. One may invoke linearity to read off the components of the projection matrix, $P = P\aij{i}{j}\ket{e_i}\bra{e^j}$, as seen in the following examples. 



% \begin{example}
% First a simple warm up. Consider the vectors
% \begin{align}
%     \ket{v} &= 2.5 \ket{1}
%     &
%     \ket{w} &= 3\ket{1} + 4\ket{2} \ .
% \end{align}
% We want to find $\ket{w_\paral}$. First we identify the unit vector in the $\ket{v}$ direction, which is simply $\ket{\hat{v}} = \ket{1}$. As a matrix, the projection operator $P = \ket{\hat{v}} \bra{\hat{v}}$ is
% \begin{align}
%     \ket{\hat{v}} \bra{\hat{v}} = \ket{1}\bra{1} 
%     &=
%     \begin{pmatrix}
%         1 & \\
%         & 0
%     \end{pmatrix}
%     \ .
% \end{align}
% Int he last step we have matched the components to  $P = P\aij{i}{j}\ket{i}\bra{j}$.
% Applying this to $\ket{w}$ gives
% \begin{align}
%     \ket{\hat{v}} \bra{\hat{v}}w\ra = \ket{1}\bra{1}\left(3\ket{1}+ 4\ket{2}\right)
%     = 3 \ket{1} \ .
% \end{align}
% \end{example} 


\section{Gram-Schmidt}

The metric allows us to determine whether a basis is orthonormal or not. We \emph{like} orthnormal bases. Sometimes life gives us garbage bases that are not orthonormal. What are we to do? The Gram--Schmidt procedure is a way to use the metric to take any basis and convert it into an orthonormal basis. We already met the key ideas in Section~\ref{sec:projections}. Here we focus on the two dimensional case. Refer back to Fig.~\ref{fig:gram:projection}: we have two vectors, $\ket{v}$ and $\ket{w}$. How do we turn these into two orthonormal basis vectors?
%
The general procedure is:
\begin{enumerate}
    \item Grab one vector from the pile of garbage basis vectors. \emph{Normalize} this vector and stick it in your bag of nice basis vectors. 
    \item Grab another vector from the pile and do whatever you need to do to make it fit as a nice basis vector with the other nice basis vectors in your bag. Start by finding what new `direction' the vector identifies relative to the existing basis vectors. That means finding a direction that is perpendicular\sidenote{It helps to think of this as the part that is not parallel to any existing basis vectors.} to the existing basis vectors. Then normalize the vector pointing in that new direction. The resulting vector can now be placed in the bag of nice basis vectors as an additional orthonormal basis vector.
    \item Repeat the previous step until you have a complete basis. 
\end{enumerate}

We can start with $\ket{v}$ and simply normalize it:
\begin{align}
    \ket{{e}_1}\defeq \ket{\hat{v}}
\end{align}
is the first basis vector. The next step is to take the next vecotr, $\ket{w}$ and pull out the perpendicular part of $\ket{w}$. In Section~\ref{sec:projections} we saw how to do this: simply decompose
\begin{align}
    \ket{w} &= \ket{w_\paral} + \ket{w_\perp}
    &
    \ket{w_\paral} &=  \la {\hat v} , w\ra  \ket{\hat v}  \ .
\end{align}
This gives us $\ket{w_\perp}$ in terms of two vectors $\ket{w}$ and $\ket{w_\paral}$ that we know:
\begin{align}
    \ket{w_\perp} = \ket{w} - \la {\hat v} , w \ra  \ket{\hat{v}} \ .
\end{align}

 We can thus normalize this and that gives us the second basis vector: 
\begin{align}
    \ket{e_2} \defeq \ket{\hat w_\perp}
\end{align}
If you have more vectors, then you can take your next vector and follow the above steps. The only difference with the iterative step is that the decomposition of a vector 
\begin{align}
\ket{u} = \ket{u_\paral} + \ket{u_\perp}    
\end{align}
is such that
\begin{align}
    \ket{u_\paral} = \la \hat{v}, u \ra \ket{\hat{v}} + \la \hat{w}, u \ra \ket{\hat{w}} + \cdots \ ,
\end{align}
where the $\cdots$ represents projections onto all other existing basis vectors. You then normalize the ``leftover'' perpedicular part, $\ket{\hat{u}_\perp}$, and add that to your bag of orthonormal basis vectors.


\begin{exercise}
Would you have the same basis vectors if you started with $\ket{w}$ and then took the perpendicular part of $\ket{v}$? Explain why or why not; your explanation must include a picture illustrating the result.
\end{exercise}


\section{Adjoints}
\label{sec:adjoint}

\begin{quote}
Yes, \emph{this} is the transpose! (c.f.~Figure~\ref{fig:is:this:a:transpose})
\end{quote}

In Section~\ref{sec:transpose} we introduced the \emph{transpose} of a matrix. We noted that while the working definition there is easy to explain, it does not properly capture the core mathematical idea. The reason for this was that we needed to appreciate that matrices are not arrays of numbers, but are linear transformations from vectors into other vectors. The definition that captures this is called the adjoint of a matrix. The \textbf{adjoint}\index{adjoint} of a matrix\sidenote{Here we mean a linear transformations that take vectors to vectors.} $M$ is written $M^\dag$ and is defined by the relation
\begin{align}
    \la M^\dag w, v \ra = \la w, M v \ra \ . 
\end{align}
Using the pre-filled inner product definition of dual vectors \eqref{eq:dual:vec:is:pre:filled:inner:product}, we may rewrite this as
\begin{align}
    \la M^\dag w \mid v \ra = \la w \mid M v \ra \ .
\end{align}
Sometimes this is expressed as $\la w | M | v \ra$ where it is understood that $M$ acting to the \emph{left} onto the bra $\la w |$ is the same thing as taking the dual vector of $M|w\ra$. 


We show below that for real vector spaces, $M^\dag$ is precisely the transpose $M^\trans$. However, this definition generalizes to complex vector spaces where $M^\dag$ is called the \textbf{Hermitian conjugate}. We use the $M^\dagger$ notation to emphasize the utility of this more general definition. 

\begin{marginfigure}%[th]
    \includegraphics[width=\textwidth]{figures/MordorIndices.jpg}
    \captionsetup{font={scriptsize,sf}}
    \caption{``One does not simply walk into Mordor'' meme, adapted to the transpose. Generated at \url{https://imgflip.com/i/8opffm}.}
    \label{fig:transpose:mordor}
\end{marginfigure}
\begin{example}[One does not simply swap the indices] See Figure~\ref{fig:transpose:mordor}. What could possibly go wrong with the usual ``swap the indices'' definition of the transpose? The problem is that in our sophisticated upper/lower index notation, we need to deal with the height of the indices. Matrices have their first index upper and their second index lower. If $M=M\aij{i}{j}\ket{i}\bra{j}$, then what is the appropriate index structure of the transpose?
\begin{align}
    M^\trans \stackrel{?}{=} 
    \begin{cases}
    (M^\trans)\aij{j}{i}\ket{j}\bra{i}    & \text{ or }\\
    (M^\trans)_{j}^{\phantom{j}i}\ket{i}\bra{j}
    \end{cases}
    \quad\text{.}
\end{align}
The adjoint gives us a way to make sense of this rigorously. The existence of a metric is critical in this analysis. 
\end{example}

\begin{example}
By the symmetry of the metric we could have equivalently written
\begin{align}
    \la M^\dag \vec{w}, \vec{v} \ra =
    \la \vec{w} ,  M \vec{v} \ra =
    \la \vec{v}, M^\dag \vec{w} \ra =
    \la M \vec{v}, \vec{w} \ra 
    \ .
\end{align}


The point of the metric is \emph{not} that you're moving the matrix $M$ from one slot to the other, it is that you are changing from a transformation $M$ that acts on $\ket{v}$ to a transformation $M^\dag$ that acts on the \emph{other} input vector $\ket{w}$ defined such that the inner product is the same.
\end{example}

\subsection{The indices of the `transpose'}
When $M$ is a real matrix, we usually say $M^\dag \equiv A^\trans$. What does this \emph{mean}? Write out the components of each side:
\begin{align}
    \la M^\dag v, w \ra
    &=
   w^k \, g_{ki}  (M^\dag)\aij{i}{\ell} v^\ell  \,  
    &
    \la  v,M w \ra
    &=
   v^\ell\, g_{\ell j}  M\aij{j}{k} w^k  
   \ .
   \label{eq:adjoint:transpose:int1}
\end{align}
Now set these two inner products equal to each other gives\sidenote{Here we have canceled $w^kv^\ell$ on both sides. The result is correct, but this justification is suspect because the $k$ and $\ell$ indices are summed over. What we have really done is used the fact \eqref{eq:adjoint:transpose:int1} is true for \emph{all} allowed $w^k$ and $v^\ell$ so that the rest of the equality cannot depend on those factors. Indeed, the equality \eqref{eq:adjoint:transpose:int1} must be true for \emph{each} $k$ and $\ell$. This is a subtle argument that lets use do a slick manipulation. It is worth thinking carefully about it.}
\begin{align}
    g_{ki}(M^\dag)\aij{i}{\ell}
    &=
    g_{\ell j} M\aij{j}{k}
    \label{eq:adjoint:transpose:int2}
\end{align}
Contracting both sides with the inverse metric $g^{km}$ gives
\begin{align}
    (M^\dag)\aij{m}{\ell} &= g_{\ell j} M\aij{j}{k} g^{km} \ .
    \label{eq:adjoint:def}
\end{align}
This relation \emph{defines} the components adjoint, $(M^\dag)\aij{i}{j}$. We see that the adjoint has the normal index structure of a linear transformation: first index raised, second index lowered. Further, when the metric is diagonal, it is clear that the first index of $M^\dag$ corresponds to the second index of $M$ and vice  versa. In this sense, the order of the indices have swapped.  In fact, the right-hand side of  \eqref{eq:adjoint:def} can be written with the metric implicit:
\begin{align}
    g_{\ell j} M\aij{j}{k} g^{km} \equiv 
    M_\ell^{\phantom{\ell}m} \ ,
\end{align}
though we shall avoid this notation since it is a bit uncommon.



\begin{exercise}
In your own words, articulate the difference between a matrix $A$, its inverse $A\inv$, and its transpose $A^\trans$. What if $A$ is a rotation?
\end{exercise}

\begin{exercise}
Show that taking the adjoint twice returns you to the original matrix, $(A^\dag)^\dag = A$.
\end{exercise}

% \begin{example}
% The adjoint also gives us another way to understand why isometries act differently on column vectors (ket) versus row vectors (bra). The metric lets us convert a column vectors into row vectors, as we saw in Section~\ref{sec:machine:to:make:row:vectors}. Start with a vector $\ket{v}$ with components $v^i$. Under a rotation, these components become $(v')^i = R\aij{i}{j}v^j$. From the transformed ket $\ket{v'} = (v')^i \ket{e_i}$ we invoke the metric to convert this into a bra:
% \begin{align}
% \bra{v'} = \la v', \smallslot \ra = \la \smallslot, v' \ra \equiv v'_i \bra{i} \ .    
% \end{align}
% In the last line we have implicitly used the metric: we have a definition for $(v')^i$, so the meaning of $(v')_i$
% is understood to be
% \begin{align}
%     (v')_i = g_{ij} R\aij{j}{k}v^k \ .
% \end{align}
% What are the components of $(v')_i$ relative to $v_i$? To do that we can insert $\one$ in a clever way:
% \begin{align}
%     g^{mn}g_{n\ell} = \delta^m_\ell
% \end{align}
% so that 
% \begin{align}
%     (v')_i = g_{ij} R\aij{j}{k} g^{k\ell}g_{\ell m} v^m
%     = 
%     \left(g_{ij} R\aij{j}{k} g^{k\ell}\right)
%     v_\ell
%     = 
%     v_\ell(R^\dag)\aij{\ell}{i} 
%     \ ,
% \end{align}
% where again we have implicitly \emph{defined} $v_\ell = g_{\ell m}v^m$. That's just what metrics do. That's what they \emph{do}.\footnote{\url{https://www.youtube.com/watch?v=Akq0xeu-RHE}} In the last equality we recognize the definition of the adjoint, \eqref{eq:adjoint:def}. In our case of a real metric space, of course, this is simply $R^\dag = R^T$. This confirms for us that the objects with lower indices transform by contracting with $R^\dag$, which we identified in \eqref{eq:rotation:definition:one} to be the inverse of the rotation $R$. So there you have it: objects with lower indices transform with $R\inv = R^\dag$. 

% We can do a sanity check for this. The inner product is a number and should not transform under rotations. We can see what happens:
% \begin{align}
%     \la v , w \ra \to \la Rv, Rw \ra = \la v, R^\dag R w \ra = \la v, w \ra \ .
% \end{align}
% where we used $R^\dag R = \one$. 




% % Starting with a ket $\ket{v}$, we define the bra
% % \begin{align}
% % \bra{v} = \la v, \smallslot \ra = \la \smallslot, v \ra \ .    
% % \end{align}
% % Now act $\bra{v}$ on a ket $\ket{w}$. By linearity, I can also transform the ket by a matrix $A$ and then feed $\ket{Aw}=A\ket{w}$ to $\bra{v}$:
% % \begin{align}
% %     \la{v} | {Aw} \ra = \bra{v}A \ket{w} = \la v, Aw \ra \equiv \la A^\dag v, w\ra = \bra{A^\dag v} w\ra \ .
% % \end{align}
% % Now suppose $A$ is an isometry. Assuming that the space is real\footnote{That is: every time we say `number' we mean a \emph{real} number, $\RR$.} and that we have the Euclidean metric, this is simply the assumption that $A=R$, a rotation. 
% \end{example}

\subsection{Self-Adjoint}

There is a special class of matrices that are called \textbf{self-adjoint}\index{self-adjoint} or, equivalently, \textbf{Hermitian}\index{Hermitian}. For real matrices, this turns out to simply mean that they are \emph{symmetric}. These are matrices $M$ for which 
\begin{align}
    M^\dag = M \ .
\end{align}
We shall come to appreciate that these matrices not only have nice properties, but that they are precisely the types of matrices that play significant roles in physics.

We can see that self-adjoint matrices are symmetric. The definition of self-adjoint is:
\begin{align}
    \la M w, v \ra = \la w, M^\dag v \ra = \la w, M v \ra \ .
\end{align}
By the symmetry of the metric, we also have 
\begin{align}
    \la M w, v \ra = \la v, Mw \ra \ .
\end{align}
Combining these gives, for a self-adjoint matrix $M$,
\begin{align}
    \la w, M v \ra = \la v, Mw \ra \ .
\end{align}
We apply this result to basis vectors, $\ket{w} = \ket{e_i}$ and $\ket{v} = \ket{e_j}$ for some $i$ and $j$. Recall the \emph{pre-loaded inner product} definition of basis dual vectors in \eqref{eq:basis:dual:vectors:as:inner:prod}: $\bra{e^i} = \la e_i, \slot \ra$, we may write our self-adjoint definition as\sidenote{To be clear: the arguments of the inner product are vectors. This means that $Me_i$ as one of the arguments means $M\ket{e_i}$. } 
\begin{align}
    \la e_i,  Me_j \ra &= \la e_j, M e_i \ra \\
    \la e_i \mid M \mid e_j \ra &= \la e_j \mid M \mid e_i\ra\\
    M\aij{i}{j} &= M\aij{j}{i}
    \ .
    \label{eq:self:adjoint:is:symmetric}
\end{align}
Do not be confused by the notation: 
\begin{enumerate}
    \item $\la e_i \mid M \mid e_j \ra$ simply means that $M$ acts on the vector $\ket{e_j}$ to produce a vector $\ket{Me_j}$. Then the dual vector $\bra{e_j}$ acts on $\ket{Me_j}$ to produce a number.
    \item The number that is produced from $\la e_i \mid M \mid e_j \ra$ is the $j$--$i$ component of $M$. You can see this because the components of $\ket{e_i}$ are zero everywhere and one on the $i^\textnormal{th}$ component. (See example below.)
\end{enumerate}

\begin{example}
In matrix multiplication notation, we may write:
\begin{align}
    \begin{pmatrix}
        0 & 1
    \end{pmatrix}
    \begin{pmatrix}
        M\aij{1}{1} & M\aij{1}{2}\\
        M\aij{2}{1} & M\aij{2}{2}
    \end{pmatrix}
    \begin{pmatrix}
        1 \\ 0
    \end{pmatrix}
    &= 
    M\aij{2}{1} \ .
\end{align}
\end{example}

You may feel justifiably uncomfortable that in \eqref{eq:self:adjoint:is:symmetric} we ended up with an equation where the heights of the $i$ and $j$ indices on the left-hand side do not match with those on the right-hand side. This happened because the \emph{pre-loaded inner product} effectively raises the index of a basis vector. It is a bit of a notational quandry and has to do with the metric's role as an index lower-er or raise-er. If you wish to be more comfortable with this, you can treat \eqref{eq:self:adjoint:is:symmetric} as a component-wise identity. Alternatively, you may think about it component-wise using \eqref{eq:adjoint:def}: a self-adjoint operator has $(M^\dag)\aij{i}{j} = M\aij{i}{j}$. There are factors of the metric that account for raising and lowering the indices.


\begin{subappendices}

\section{A Matrix in Another Basis}

Consider a linear transformation from vectors to vectors, $A$. In a given basis, this linear transformation is represented as a matrix.\sidenote{I sometimes write `linear transformation' and `matrix' interchangeably. Formally the matrix is the array of numbers that is indicated by indices, whereas the linear transformation is the index-less operation that does not assume a basis.} Often one would like to describe the components of the matrix in a different basis. Here's how one converts from the components $A\aij{i}{j}$ in the $\ket{e_{i}}$ basis to its components $\tilde A\aij{a}{b}$ in a different basis, $\ket{f_a}$.\sidenote{Please review Section~\ref{sec:basis:of:matrices} to recall why the ket-bra is the natural basis of matrices.}
\begin{align}
    A &= A\aij{i}{j}\ket{e_i}\!\bra{e^j}
    \\&=
    A\aij{i}{j}\quad
    \one
    \ket{e_i}\!\bra{e^j}
    \one
    \\&=
    A\aij{i}{j}\quad
    \ket{f_a}\!\la {f^a} \mid 
     {e_i} \ra\!  \la {e^j} \mid
    {f_b}\ra \!\bra{f^b}
    \\&=
    \la {f^a} \mid 
     {e_i} \ra
    A\aij{i}{j}
    \la {e^j} \mid
    {f_b}\ra
    \quad
    \ket{f_a}\!\bra{f^b}
\end{align}
Now please review Figures~\ref{fig:basis:change:ef} and \ref{fig:basis:change:fe}. We realize that the numbers $\la {f^a} \mid 
     {e_i} \ra$ and $\la {e^j} \mid
    {f_b}\ra$ are simply components of the rotation matrix between the two bases. Picking a convention\sidenote{You could have swapped which one is $R$ and which is $R\inv$. After all, the inverse of a rotation is simply a rotation in the opposite direction.} for which one we call $R$ and which we call $R\inv$ gives:
\begin{align}
    A &= (R\inv)\aij{a}{i}
    A\aij{i}{j}
    R\aij{j}{b}
    \quad
    \ket{f_a}\!\bra{f^b} 
    \\&\equiv
    \tilde A\aij{a}{b}
    \quad
    \ket{f_a}\!\bra{f^b}
    \ .
\end{align}
This means that the matrix components $\hat A\aij{a}{b}$ in the $\ket{f_{1,2}}$ basis are related to the matrix components $A\aij{i}{j}$ in the $\ket{e_{1,2}}$ basis by
\begin{align}
    \tilde A\aij{a}{b} &= (R\inv)\aij{a}{i}
    A\aij{i}{j}
    R\aij{j}{b}
    &
    A\aij{i}{j} = R\aij{i}{a}\tilde A\aij{a}{b}(R\inv)\aij{b}{j} \ .
\end{align}



\section{Discrete isometries}
\label{sec:discrete:isometries}

One of the notable features of rotations is that you can rotate by any continuous amount, $\theta$. You can rotate by $\theta=\pi/2$, or $\theta = \pi/2 + \varepsilon$ for, say, $\varepsilon = .0001$. Thus the rotation matrix parameterizes an infinite number of possible rotations. Because there is a continuous parameter for the amount by which one performs a rotation, we say that this is a continuous isometry. There are, however, often \emph{discrete} isometries where there is simply a handful of matrices under which the metric is invariant. In Euclidean two dimensional space, one such matrix is 
\begin{align}
    P = \begin{pmatrix}
        1 & \\
        & -1 
    \end{pmatrix} \ .
\end{align}
Under $P$, the second component of a vector (in the standard basis) changes sign. Because the metric transforms with two factors of $P$---here we use $P^\dag = P$---it remains unchanged:
\begin{align}
    P^\text{T}gP = 
    \begin{pmatrix}
        1 & \\ 
        & -1 
    \end{pmatrix}
    \begin{pmatrix}
        1 & \\ 
        & 1 
    \end{pmatrix}
    \begin{pmatrix}
        1 & \\ 
        & -1 
    \end{pmatrix}
    = 
    \begin{pmatrix}
        1 & \\ 
        & 1 
    \end{pmatrix} \ .
\end{align}
\begin{exercise}
Suppose that the components of the metric in $\RR^2$ are not diagonal. That is, $g_{12} = g_{21}\neq 0$. Is $P$ an isometry of this metric space?
\end{exercise}

The total set of isometries is every combination of rotations and discrete transformations. We often say that this is a \textbf{group}\index{group}. One may write a generic isometry as a product of rotations and discrete transformations. 
\begin{example}
The following transformation is a combination of a discrete transformation $P$ and a rotation:
\begin{align}
    \tilde R = PR =
    \begin{pmatrix}
        1 & \\ 
        & 1 
    \end{pmatrix}
    \begin{pmatrix}
        c & -s\\ 
        s & \pp c
    \end{pmatrix}
    =
    \begin{pmatrix}
        c & -s\\ 
        -s & - c
    \end{pmatrix} \ .
    \label{eq:discrete:isometry:eg}
\end{align}
This matrix does \emph{not} look like a rotation. The diagonal elements differ by a sign, while the off-diagonal elements are the same. If we did not know about discrete isometries, we would have said that this matrix does not map one orthonormal basis to another. In fact, it does.
\end{example}
\begin{exercise}
Compare the rotation-and-discrete transformation in \eqref{eq:discrete:isometry:eg} to a transformation that is simply a rotation. For concreteness, pick a rotation angle $\theta = \pi/4$ so that $c=s=1/\sqrt{2}$. Apply the rotations $R$ and $\tilde R$ to the standard basis. How do the ``rotated'' bases differ?
\end{exercise}

In three dimensions this notion shows up as \emph{parity}. Here are two ways to represent the parity transformation in three dimensions:
\begin{align}
    P &=
    \begin{pmatrix}
        -1 & & \\
        & -1 & \\
        & & -1
    \end{pmatrix} 
    &
    P =
    \begin{pmatrix}
        1 & & \\
        & 1 & \\
        & & -1
    \end{pmatrix}
    \ .
\end{align}
The reason why there are different ways to represent this discrete symmetry is that the two are related by a rotation. 
\begin{exercise}
Describe---in words or drawing if that is simpler---the rotation that relates the two parity matrices above. It may help to think about what each representation of $P$ does to the standard basis vectors. 
\end{exercise}

You can tell that the two are related by a rotation if you think about all that ``right hand rule'' stuff you learned when you first did angular momentum in mechanics. The right hand rule enforced a certain \emph{orientation} of your basis vectors. The positive $x$, $y$, and $z$ directions are chosen to have a particular orientation: if you curl the fingers of your right hand from the $x$-axis to the $y$-axis, then your thumb is pointing in the $z$ direction. You can rotate your hand in any way, but this orientation stays the same. Similarly, if you are contrarian and decided to define your axes by a ``left hand rule,'' then (1) several of your angular momentum-related results are off by a minus sign, and (2) there is no way to rotate your left hand to make it match the orientation of your right hand. This notion is pervasive in nature in three spatial dimensions: fundamental particles have a spin that defines an orientation that is either left- or right-handed. Even protein can have a \emph{chirality} where identical sequences of amino acids may spiral in one direction or another---and biological life on earth can only use one orientation and not the other. 

By manipulating objects (like your right hand) you have an intuitive sense of rotations in three-dimensions. What may be less obvious is that you also have a simple tool for understanding what a parity transformation does: look in the mirror. If you raise your right hand, then the image of you in the mirror raises their left hand. If you do physics with a right-handed coordinate system, then your mirror image will do physics in a left-handed coordinate system. You can bet that the image of you is composed of proteins with the opposite chirality and subatomic particles spinning in the opposite orientation. 
\begin{exercise}
Explain why your image in a mirror exchanges left and right but does not swap up and down.\footnote{This is a question that was posed in my freshman physics class. It was not the first time that I heard the question, but I was surprised that it is a mathematical question in the realm of physics. I have since come to appreciate the significant role that parity plays in many aspects of physics. I hope that soon you will as well.}
\end{exercise}



\end{subappendices}

\include{Zchap_LA_07_det}
\include{Zchap_LA_08_eigen}
%!TEX root = Physics231.tex

\chapter{Complex Spaces}


It is time to embrace complex numbers. A complex vector space is one where all of the numbers that we previously assumed were real can now be complex: $\RR \to \CC$. This means that a linear combination of vectors $\alpha\vec{v}+ \beta\vec{w}$ may have complex coefficients, $\alpha,\beta\in\CC$. 

\section{Review of complex numbers}

A complex number has a real and imaginary part. The imaginary part is proportional to  $i \defeq \sqrt{-1}$. We typically write complex numbers $z$ in either the Cartesian or polar form:
\begin{align}
    z &= x+iy 
    &
    z &= r \e^{i\theta}
    \ .
\end{align}
Here $x,y,r$, and $\theta$ are real parameters. The radius is greater than zero ($r>0$) and the angle $\theta$ is defined modulo $2\pi$ so that $\e^{i(\theta+2\pi)} = \e^{i\theta}$.

\paragraph{Are complex numbers a vector space?}
In the Cartesian expression, you might be tempted to read this like a vector space with two basis vectors: $\ket{e_1}=1$ and $\ket{e_2} = 2$ and a complex number $z$ is the linear combination of basis vectors $\ket{z} = x \ket{e_1} + y \ket{e_2}$. The addition rule for complex numbers reaffirms this:
\begin{align}
    (a+ib)+ (x+iy) = (a+x) + i (b+y) \ .
\end{align}
So we can see that complex numbers have the same structure as the vector space $\RR^2$. However, it is incorrect to think that the vector space $\RR^2$ is the same as complex numbers, $\CC$. This is because complex numbers have additional structure beyond that of a vector space---in fact, beyond even that of a metric space. 

This additional structure is a rule by which two `vectors' can be multiplied to produce another vector. We have no such rule in a general vector space.\sidenote{And don't even get me started on the cross product. That is only valid in three dimensions due to a coincidence in the dimensionality of the Hodge dual of two vectors, see Section~\ref{sec:cross:product}.} The rule $i^2 = -1$ means that
\begin{align}
    (a+ib)(x+iy) = (ax - by) + i (ax + by) \ .
\end{align}
Thus complex numbers are a vector space \emph{with additional structure}. The nature of this additional structure is outside the scope of this course, but you can learn more about it in a field called \textbf{geometric algebra}.\footnote{One of my favorite overviews of this field is Jack Rusher's talk ``From Geometry to Algebra and Back Again: 4000 Years of Papers'' at the Strange Loop Conference, \url{https://youtu.be/1cRFfYQYGxE?si=x9U0mIfeqXy3QHLd}. Another nice one is YouTuber sudgylacmoe's ``Swift Introduction to Geometric Algebra,'' \url{https://youtu.be/60z_hpEAtD8}.}\autocite{Doran:2007tqa} This topic has some really neat manifestations in rotational dynamics\footnote{\url{https://www.youtube.com/watch?v=BN4coAyF5PE&feature=youtu.be}}\autocite{10.1119/5.0109883}, relativity\autocite{10.1119/1.4734014}, and electrodynamics\footnote{\url{https://youtu.be/WN_4j8v6cXo}}\autocite{Dressel:2014fna}.

\paragraph{The exponential}
The polar form of a complex number is related to the Cartesian form through Euler's identity:
\begin{align}
    \e^{i\theta} = \cos\theta + i\sin\theta
\end{align}
This identification is somewhat mysterious the first time you see it.\sidenote{Perhaps it remains mysterious and most of us simply get used to it.} One way to appreciate this is to connect it to the definition of the exponential as a repetition of many small transformations:
\begin{align}
    \e^x \defeq \lim_{n\to\infty} 
    \prod _n \left(1+\frac{x}{n}\right)^n \ .
\end{align}
The right-hand side is a product of factors $(1+\epsilon)$ where $\epsilon = x/n$ is a \emph{small} shift in the $x$ direction.\sidenote{For real numbers, there is only one direction: the real axis. But there are often cases where $x$ is a more sophisticated object, such as a matrix.} As $n\to\infty$, the factor $(1+\epsilon)$ is thus a small transformation and the exponential is the product of $n\to\infty$ of these small transformations. 

\begin{example}[Exponential representation of rotations.] Consider rotations in two dimensions. An infinitesimal rotation is:
\begin{align}
    R(1+\epsilon) = 
    \begin{pmatrix}
        1 & -\varepsilon \\
        \varepsilon & \pp 1
    \end{pmatrix}
    = 
    \one + \varepsilon 
    \begin{pmatrix}
        0 & - 1 \\
        1 & \pp 0
    \end{pmatrix} \ .
\end{align}
You can check that this is exactly what you get from a first-order Taylor expansion of the rotation matrix:
\begin{align}
    R(\theta) = 
    R(0)
    + \theta 
    \frac{\D{}}{\D{\theta}}  + \mathcal O(\theta^2) \ .
\end{align}
Please explicitly write out the two matrices on the right-hand side if this is not clear. 

The matrix 
\begin{align}
    T =\begin{pmatrix}
        0 & - 1 \\
        1 & \pp 0
    \end{pmatrix} 
\end{align}
has a special name: it is the \textbf{generator} of rotations. We may write a finite rotation matrix as the exponentiation of this generator:
\begin{align}
    R(\theta)  = 
    \lim_{n\to\infty} \frac{1}{n}
    \left(1+\frac{\theta}{n} T\right)^n
    = 
    \e^{i\theta T} \ .
\end{align}
You may also explicitly write out the first few terms of $e^{i\theta T}$ as a power series to see that the even terms show up on the diagonal and the odd terms show up on the off-diagonal. The terms sum to the usual cosine and sine expansions.
\end{example}

The imaginary number $i$ acts as a rotation in the complex plane.\sidenote{This is `obvious' when writing $i=\e^{i\pi/2}$, but here we are motivating the exponential representation and so cannot use the exponential representation to motivate itself.} Taking a complex number $z$ and multiplying it by $(1+i\epsilon )$ returns a complex number $z'$ that is approximately $z$ but rotated by a small angle $\epsilon$. The length of $z'$ is the same as the length of $z$ up to factors of $\mathcal O(\epsilon^2)$.
\begin{exercise}
The length of a complex number $z=a+ib$ is $|z| = \sqrt{a^2 + b^2}$. Show that the lengths of $z$ and $(1+i \epsilon )z$ have the same length up to corrections of order $\mathcal O(\epsilon^2)$.
\end{exercise}
Thus the polar form $z=r\e^{i\theta}$ is (1) rescaling by a real number $r$, and (2) a rotation by an angle $\theta$ in the complex plane. 


\section{Phasors}

This section is outside the main scope of this course.

One way that both the Cartesian and the polar pictures of a complex number come together is something called a \textbf{phasor diagram}\index{phasor}, shown in Figure~\ref{fig:figure:phasor}.
\begin{marginfigure}%[th]
    \includegraphics[width=.5\textwidth]{figures/Phasor.pdf}
    \captionsetup{font={scriptsize,sf}}
    \caption{Example of a phasor diagram showing the sum, in blue, of many complex numbers of unit length. Each complex number is $\hat z = e^{iS[q]}$, where $q$ represents a path. The function $S$ is an action. The middle of the diagram is the region where the action is minimized so that it does not change appreciably over nearby paths. This region of \emph{least action} dominates in the classical limit, where the spirals become tighter.}
    \label{fig:figure:phasor}
\end{marginfigure}
At its heart, a phasor is simply a complex number of unit length. That means it is of the form $\hat z = e^{i\theta}$. 
Phasors are a nice way to think about an object at the heart of quantum mechanics: the path integral. I strongly recommend reading Feynman's popular science book, \emph{QED: The Strange Theory of Light and Matter}\autocite{Feynman:1986er}. The material is based on Feynman's Douglas Robb memorial lectures.\sidenote{You may still be able to find recordings of these lectures online.}
% 
Richard Feynman developed the path integral formulation of quantum mechanics in his Ph.D thesis\autocite{Feynman:1942us}. The proposal is that the probability amplitude\sidenote{The probability amplitude $\psi$ is a complex number whose magnitude squared $|\psi|^2 = \psi^*\psi$ is the probability.} to go from state $A$ to state $B$ is given by a sum of phasors up to a normalization\sidenote{The perceptive reader will note that this is an important point that I'm ignoring. This normalization assures that there are no probabilities greater than one.} that we do not belabor here. Each phasor represents the probability amplitude for a particular path between $A$ and $B$. 
\begin{example}
Consider the double slit experiment. There are two paths, one from either slit. Each path is associated with a complex number of unit modulus. The probability amplitude is the sum of the phasors for each path. 
\end{example}
The probability amplitude is thus
\begin{align}
    \Psi(A\to B) \propto \sum_\textnormal{path}
    \e^{iS[q_\textnormal{path}]/\hbar} \ ,
\end{align}
where $q_\textnormal{path}(t)$ is some trajectory that goes from $A$ to $B$ as a function of time $t$. The function $S[q]$ is the usual action that you met in mechanics: it is the time integral over the Lagrangian, $S=\int\D{t}\, L[q]$, which is itself the difference between kinetic and potential energy. The number $\hbar$ is Planck's constant; it is required by dimensional analysis and with respect to human scales it is quite small.

Figure~\ref{fig:figure:phasor} shows that when the action changes quickly as a function of $q$, the phase in $e^{iS[q]}$ changes quickly and the amplitude contribution from nearby paths cancel each other out. However, the paths near the minimum of the action $S'[q]=0$ have a nearly constant action and so those phasors add constructively. The contribution to the amplitude is dominated by these paths that minimize the action. The classical limit of quantum mechanics is the case where $\hbar \to 0$ so that the action varies so rapidly that only the path of minimal action contributes. One can then approximate the contribution in the sum over all paths by only the path that extremizes\sidenote{The path may \emph{either} minimize \emph{or} maximize $S[q]$. In principle it may minimize in some directions while maximizing in other directions: all that matters is that the action is \emph{stationary} for small variations of the path.} the action. This is called the \textbf{saddle point approximation}\index{saddle point approximation} or the \textbf{method of steepest descent}. Classical physics is the limit where the approximation is exact.\sidenote{But please understand: what we really means is that classical physics is an approximation of the underlying quantum theory for dynamics where $\hbar$ effects can be ignored.} This is the quantum mechanical `derivation' of the so-called \emph{principle of least action}. For a fantastic discussion, see ``The Emergence of Classical Physics'' in Padmanabhan\autocite{Padmanabhan:2015boa}.

Phasors are not used much in modern physics, but the intuition of the phasor picture of the path integral is at the heart of quantum field theory. By a curious relation between quantum and statistical physics, this is closely related to statistical field theory by a factor of $i$.



\section{Complex vector space}

A complex vector space is one where all `numbers' are now understood to be \emph{complex numbers}. This means that the components of vectors, matrices, and all tensors are complex. It means that you may take linear combinations of objects with complex coefficients. 

Observe that because complex numbers basically have the same arithmetic as real numbers---up to this funny $i^2=-1$ business---nothing really changes from the \emph{vector space} perspective. We still write vectors as
\begin{align}
    \ket{v} = v^i\ket{e_i}
\end{align}
with the understanding that $v^i$ is now a number of the form $v^i = a^i + i b^i$.\sidenote{Please do not confused the index $i$ with the imaginary number $i$. Some people use Greek letter iota $\iota$ instead of $i$, but I find that this looks like a typo where someone lost the dot.}
\begin{example}
Suppose you are in a three-dimensional  complex vector space with a vector $\ket{v}$ whose components are
\begin{align}
    v^1 &= 2+i
    &
    v^2 &= 3i
    &
    v^3 &= -2 \ .
\end{align}
Suppose you have a dual vector $\bra{w}$ whose components are
\begin{align}
    w_1 &= -i
    &
    w_2 &= 1+i
    &
    w_3 &= 2 \ .
\end{align}
Then the action of $\bra{w}$ on $\ket{v}$ is
\begin{align}
    \la w \mid v \ra &=
    (2+i)(-i)+(3i)(1+i)+(-2)2 
    =
    -7 + i
    \ . 
\end{align}
\end{example}


In the above example, observe that we never invoked a metric. Everything so far has been for a complex \emph{vector space} with no mention of the existence of a metric. 

Secondly, note that we have nothing to say about basis vectors, $\ket{e_i}$. Of course we still have basis vectors, $\ket{e_i}$, and an accompanying basis of dual vectors $\bra{e^i}$. But there is nothing different about them. We do not even have to `complexify' them! One of the big deals of linear algebra has been that we can completely abstract away what the basis vectors \emph{mean}. Whether the vector space is real or complex is purely a statement about the \emph{coefficients}---the \emph{numbers}. The abstract `vector-ness' or `tensor-ness' of different types of objects remain as abstract as before.\sidenote{To remind you: there is a basis of vectors that encodes some abstraction. It may encode velocities, differential operators, functions defined on the sphere, Fibonacci sequences, etc. Then we define linear maps on these abstract objects: this gives us dual vectors, matrices, and all other form of tensor.}



\section{Inner product on complex spaces}

Things become a bit more interesting when we complexify a metric space. We have to update our rules a bit. Previously we assumed that the metric is symmetric, $\langle v, w\rangle = \langle w, v\rangle$. This turns out to only be true for real spaces. When we have a complex vector space, the metric must satisfy symmetry with a complex conjugation:
\begin{align}
    \la v, w \ra = \la w, v \ra^* \ .
    \label{eq:complex:metric:star}
\end{align}
That is a little unusual. The general inner product is \emph{not} symmetric!

\subsection{Anti-Linearity}

One way of expressing \eqref{eq:complex:metric:star} is to say that the inner product is \textbf{antilinear}\index{antilinear} in its first component while it is linear in the second component. The word `antilinear' means that somewhere there is a complex conjugation.\sidenote{The most famous anti-linear operation in physics is time reversal. This is not something we can do in an experiment, but theoretically we can flip the direction of time in all of our equations.} This means that for complex numbers $a$ and $b$,
\begin{align}
    \la av + bu, w \ra = 
    a^* \la v, w\ra
    +
    b^* \la v, w\ra
\end{align}
whereas 
\begin{align}
    \la v, aw + bu \ra = 
    a \la v, w\ra
    +
    b \la v, u\ra \ .
\end{align}


\subsection{Norms are Real}
The norm of a complex vector is real. 
\begin{align}
    |v|^2 = \la v, v \ra = \la v, v \ra^* \ .
\end{align}
In my opinion this helps justify the notation $|v|$ for the magnitude of a vector.\sidenote{What really annoys me is that sometimes we write $|M|$ to mean the determinant of a matrix. This almost makes sense because the determinant is some intrinsic number associated with a matrix. However, the determinant can be positive or negative. This gets confusing when one needs to take the absolute value of a determinant and you do not know if $|M|$ means $\det M$ or $|\det M|$.}


\subsection{Metric}
In real space we defined the metric as the symmetric (0,2) tensor that encodes the linear transformation $V\times V \to \mathbbm{R}$ that is the inner product:
\begin{align}
    \la v , w \ra = g_{ij}v^i w^j 
    &
    g_{ij} = g_{ji}
    \ .
\end{align}
In complex space we sometimes\sidenote{You can opt not to do this. In fact, one does not often have an occasion to write out a complex space metric explicitly for most undergraduate physics courses.} write the metric with funny indices $\bar{\imath}$ and $\bar{\jmath}$ instead of $i$ and $j$ to differentiate the first and second components of $g_{\bar{\imath}j}$. This is necessary now that the inner product is no longer symmetric, but is rather symmetric up to a complex conjugation. We write:
\begin{align}
    \la v, w \ra &= g_{\bar{\imath}j}(v^*)^{\bar{\imath}}w^j \ .
\end{align}
Here $(v^*)^{\bar{\imath}}$ is simply $(v^i)^*$.

The inner product rule \eqref{eq:complex:metric:star} then tells us that
\begin{align}
    g_{\bar{\imath}j}(v^*)^{\bar{\imath}}w^j \ .
    &= 
    \left(
        g_{\bar{\jmath}i}
        (w^*)^{\bar{\jmath}}
        v^i
        \right)^*
    \\
    &= 
    g_{j\bar{\imath}}^* w^j(v^*)^{\bar{\imath}}
     \ .
\end{align}
From this we find that the metric satisfies
\begin{align}
g_{j\bar{\imath}}^* =
    g_{\bar{\imath}j} \ ,
\end{align}
where we write $g_{j\bar{\imath}}^* = (g_{\bar\imath j})^*$.  This result is not particularly surprising: the metric has the same symmetry property as the inner product. 


\subsection{Adjoint}

% What about the adjoint of a transformation? For a real vector space, the adjoint is simply the transpose. For a complex vector space, the adjoint is the transpose \emph{and} a complex conjugation on each component.
% :
% \begin{align}
%     (A^\dag)\aij{i}{j} = g^{ik}(A\aij{\ell}{k})^*g_{\ell j} \ .
% \end{align}

Recall that the adjoint of a matrix $A$ is defined by
\begin{align}
    \la A v , w \ra \defeq 
    \la v, A^\dag w \ra \ .
\end{align}
We would like to relate the components of $A^\dag$ to $A$ in a complex space. The components of $A$ in a basis $\ket{e_i}$ are 
\begin{align}
    A\aij{i}{j} = \la e^i \mid A \mid e_j \ra 
    = \la e_i, A e_j \ra \ .
\end{align}
Using \eqref{eq:complex:metric:star} we may then write:
\begin{align}
    (A^\dag)\aij{i}{j}
    = 
    \la e_i, A^\dag e_j \ra
    =
    \la A e_i, e_j \ra
    =
    \la e_j, A e_i \ra^*
    =
    (A\aij{j}{i})^* \ .
\end{align}
This means, for example, in matrix form:
\begin{align}
    \begin{pmatrix}
        A\aij 11 & A\aij 12 \\
        A\aij 21 & A\aij 22 
    \end{pmatrix}^\dag &= 
    \begin{pmatrix}
        (A\aij 11)^* & (A\aij 21)^* \\
        (A\aij 12)^* & (A\aij 12 )^*
    \end{pmatrix} \ .
\end{align}
In a complex space, the adjoint of a matrix is a combination of the transpose and the complex conjugate. This is sometimes called the \textbf{Hermitian conjugate}\index{Hermitian conjugate}.\sidenote{The Hermitian conjugate and the adjoint or a matrix are the same.} For real vector spaces, the adjoint reduces to the transpose. 


\begin{example}
Here's the adjoint of an explicit matrix:
\begin{align}
\begin{pmatrix}
        3.5 +2i & 2.3+7i \\
        8.3 + 5i & 5-2i
    \end{pmatrix}^\dag &= 
    \begin{pmatrix}
        3.5 -2i &  8.3 - 5i\\
        2.3 -7i & 5 +2i
    \end{pmatrix} \ .
\end{align}
\end{example}

It is also useful to consider the adjoint of a product of matrices, $AB$. We can see that
\begin{align}
    \la AB \vec{v}, \vec{w} \ra
    =
    \la B \vec{v},  A^\dag \vec{w} \ra
    =
    \la \vec{v}, B^\dag A^\dag \vec{w} \ra \ .
\end{align}
This tells us that
\begin{align}
    (AB)^\dag = B^\dag A^\dag \ .
    \label{eq:adjoint:of:product}
\end{align}



\begin{example} What do isometries look like in a complex vector space? Let us suppose that we have the Euclidean metric (even though the space is now complex):
\begin{align}
    g_{\bar\imath j} = \text{diag}(1,1,\cdots, 1) = \mathbbm{1} \ .
\end{align}
Let $U$ be an isometry in this space; this is a conventional name whose meaning will be clear soon. The condition for an isometry in this space is that a transformation $\ket v \to U\ket{v}$ preserves the metric. Observe that because the metric is not symmetric, but rather symmetric-up-to-a-conjugate, it is useful to write the inner product using brackets rather than simply index contraction. We have:
\begin{align}
    \la U\vec{v}, U\vec{w} \ra = \la \vec{v}, U^\dag U \vec{w} \ra \stackrel{?}{=} \la \vec{v}, \vec{w} \ra \ .
\end{align}
This tells us that the condition for $U$ to be an isometry is that $U^\dag U = \mathbbm{1}$. Matrices that satisfy this are called \textbf{unitary} matrices, thus justifying the name $U$. 
\end{example}

\begin{example}
The usual rotation matrices (obviously!) satisfy the unitary condition. Thus rotations are a simple example of a unitary matrix. There are more complicated unitary matrices, for example in two complex dimensions:
\begin{align}
    U(\theta) = 
    \begin{pmatrix}
        e^{i\theta} & 0\\
        0 & e^{-i\theta}
    \end{pmatrix}
    \label{eg:unitary:matrix:diagonal}
\end{align}
is a transformation that rephases the first $v^1 \to e^{i\theta}v^1$  and second $v^2 \to e^{-i\theta}v^2$ components of a vector by opposite phases. It should be obvious that $U^\dag U = \one$. 
\end{example}


\begin{exercise}
A bonus observation: show that \eqref{eg:unitary:matrix:diagonal} can be understood as a Taylor expansion with respect to a matrix $T$:
\begin{align}
    U(\theta) &= \sum_{n=0}^\infty \frac{\theta^n}{n!} (iT)^n
    &
    iT=
    \begin{pmatrix}
        i & 0\\
        0 & -i
    \end{pmatrix} \ .
\end{align}
The matrix $T$ is called the \textbf{generator} of the isometry $U$. 
\end{exercise}


\begin{exercise}
Show that the generator of rotations is
\begin{align}
    R(\theta) &= \sum_{n=0}^\infty \frac{\theta^n}{n!} (iT)^n
    &
    iT=
    \begin{pmatrix}
        0 & -1\\
        1 & \pp 0
    \end{pmatrix} \ .
\end{align}
The curious choice of defining $iT$ is so that the generator $T$ is Hermitian, $T^\dag = T$, which is a common convention in physics.
\end{exercise}


% \begin{table}
%     \renewcommand{\arraystretch}{1.3} % spacing between rows
%     \centering
%     \begin{tabular}{ @{} llllll @{} } \toprule % @{} removes space
%         Space
%         & Numbers
%         & Inner Product
%         & Contraction
%         & Matrix
%         & Identity
%         \\ \hline
%         Classic
%             & vector
%             & row-vector
%             & 
%             &       
%         \\
%             & $\vec{v} = v^i \bas{e}_i$
%             & $\row{w} = w_i \rbas{e}^i$
%             & $\rbas{e}^i\bas{e}_j = \delta^i_j$
%             & $M=M\aij{i}{j} \bas{e}_i\otimes \rbas{e}^j$
%             & $\one$
%         \\ \hline
%         Quantum 
%             & ket
%             & bra
%             &       
%             &
%         \\
%             & $\ket{v} = v^i\ket{i}$
%             & $\bra{w} = w_i \bra{i}$
%             & $\langle i | j\rangle = \delta^i_j$
%             & $M=M\aij{i}{j} \ket{i}\bra{j}$
%             & $\ket{i}\bra{i}$
%         \\ \bottomrule
%     \end{tabular}
%     \caption{
%         Dictionary between `classic' notation and `quantum' (bra--ket) notation.
%         \label{tab:classic:bra:ket:dictionary}
%     }
% \end{table}

\section{Self-Adjoint Matrices}

\textbf{Self-adjoint matrices} are those for which
\begin{align}
M^\dag = M \ .     
\label{eq:self:adjoint}
\end{align}
These are known equivalently as \textbf{Hermitian matrices}\index{Hermitian matrix}.
These are what we mean by `nice' matrices. They generalize the notion of symmetric matrices. 
\begin{example}
The \textbf{Pauli matrices} are a set of three $2\times 2$ Hermitian matrices:
\begin{align}
    \sigma^1
    &=
    \begin{pmatrix}
        0 & 1 \\
        1 & 0
    \end{pmatrix}
    &
    \sigma^2
    &=
    \begin{pmatrix}
        0 & -i \\
        i & \pp 0
    \end{pmatrix}
    &
    \sigma^3
    &=
    \begin{pmatrix}
        1 & \pp 0 \\
        0 & -1
    \end{pmatrix} \ .
\end{align}
\end{example}
\begin{exercise}
Show that any \emph{real} linear combination of the Pauli matrices is a Hermitian matrix:
\begin{align}
    a \sigma^1 + b \sigma^2 + c\sigma^3 \ .
\end{align}
Here the numbers $a$, $b$, and $c$ are real. Explain what goes wrong if these numbers are complex---in that case the matrix is no longer Hermitian.

Argue that \emph{all} $2\times 2$ Hermitian matrices take the above form. Yes, this means that the Pauli matrices are a basis for the space of $2\times 2$ Hermitian matrices, which is itself a \emph{real} vector space because the coefficients $a$, $b$, and $c$ are real.
\end{exercise}

% \subsection{Eigen-bra}

% If $\ket{\xi}$ is a normalized eigenvector of the matrix $M$ with eigenvalues $\lambda$, $M\ket{\xi} = \lambda\ket{\xi}$, then 
% \begin{align}
%     \la v, M \xi \ra = \lambda \la v, \xi \ra 
%     \ .
% \end{align}
% Similarly, if we use the definition of the adjoint:
% \begin{align}
%     \la v, M \xi \ra = \la M^\dag\xi, v \ra 
%     = \bar\lambda^* \la \xi, v \ra
%     = \bar\lambda^* \la v, \xi \ra^*
%     \ .
% \end{align}
% Here $\bar\lambda$ is the eigenvalue of $M^\dag$ on $\ket{\xi}$. We have used the antilinearity of the inner product with respect to its first component. 




\subsection{Eigenvalues}

Here is the first really neat thing about nice matrices:
\begin{theorem}[Real eigenvalues]
The eigenvalues of self-adjoint matrices are real.
\end{theorem}
\begin{proof}
Let $M=M^\dag$ be a self-adjoint matrix. Let $\ket{\xi}$ be  an eigenvector eigenvalue $\lambda$:
\begin{align}
    M\ket{\xi} &= \lambda\ket{\xi} 
    \ .
\end{align}
Then consider:
\begin{align}
    \la M \xi, \xi \ra 
    = 
    \la \lambda \xi, \xi \ra
    = 
    \lambda^* \la \xi, \xi \ra \ .
\end{align}
Using the adjoint, we may also write this as
\begin{align}
    \la M \xi, \xi \ra 
    = 
    \la \xi, M^\dag \xi \ra
    =
    \la \xi, M \xi \ra
    = 
    \lambda \la \xi, \xi \ra \ .
\end{align}
Taking the difference of these two expressions gives
\begin{align}
    \left(\lambda^*-\lambda\right) \la \xi, \xi \ra  = 0 \ .
\end{align}
We know that $ \la \xi, \xi \ra\neq 0$, which means that $\lambda^* = \lambda$. This means that $\lambda$ is pure real. 
\end{proof}
Because real vector spaces are complex vector spaces with  a restriction to real coefficients, this theorem holsd for real vector spaces as well. 

\subsection{Eigenvectors}

Here is the second really neat thing about nice matrices:
\begin{theorem}[Orthogonal eigenvectors]
The eigenvectors of self-adjoint matrices are orthogonal.
\end{theorem}
\begin{proof}
Let $M=M^\dag$ be a self-adjoint matrix. Let $\xi_a$ and $\xi_b$ be two distinct eigenvectors ($b\neq a$) with distinct eigenvalues $\lambda_a \neq \lambda_b$:
\begin{align}
    M\ket{\xi_a} &= \lambda_a\ket{\xi_a} 
    &
    M\ket{\xi_b} &= \lambda_b\ket{\xi_b} 
    \ .
\end{align}
Then we may consider
\begin{align}
    \la M \xi_a, \xi_b \ra 
    &= 
    \la \xi_a, M^\dag \xi_b \ra 
    =
    \la \xi_a, M \xi_b \ra
    .
\end{align}
We may use $\la M \xi_a, \xi_b \ra = \la \xi_b, M \xi_a \ra^* = \lambda_b^*\la \xi_a, \xi_b \ra$ to then write
\begin{align}
    (\lambda_a - \lambda_b)\la \xi_a,\xi_b\ra = 0 \ .
\end{align}
We have used $\lambda_a^* = \lambda^a$ since $M$ is self-adjoint. Because $\lambda_a\neq \lambda_b$, we must have that $\la \xi_a, \xi_b\ra = 0$. This means that the two eigenvectors are orthogonal. 
\end{proof}
Given a set of orthogonal eigenvectors, one may normalize them by dividing by their lengths to produce an orthonormal eigenbasis. 

\subsection{Physical observables and nice bases}

What is the significance of all this? 

\paragraph{Real eigenvalues}
You would be forgiven if your first year of university physics left you believing that complex numbers are merely a mathematical convenience. Perhaps you first met complex numbers in physics when learning about waves and the complex exponential $e^{i\theta}$ was a shorthand for writing $\cos\theta$. Sure, you would have complex numbers floating around in your equations, but you were told that what we \emph{really} mean is we always take the real part. Perhaps you learned about complex impedance in a circuit lab, where again you learned that the complex phase is simply a shortcut for trigonometric manipulations that are ultimately about \emph{real} objects. 

Soon in your education you will find that this reassurance appears to be false. What we know of quantum mechanics tells us that nature appears to be \emph{intrinsically} complex. Attempts to construct a manifestly real theory of quantum mechanics end up re-inventing complex numbers.

So this brings up to a puzzle: if nature is so intrinsically complex, why is everything that we measure \emph{real}? The reason turns out to be that all measurements are associated with self-adjoint (Hermitian) operators (matrices). The value of a measurement is an allowed eigenvalue of that operator. Because the operator is self-adjoint, these eigenvalues are real. 

\paragraph{Orthogonal eigenvectors} The states that have well-defined measurable properties are the eigenvectors of those observables. That these eigenvectors are orthogonal means that we can form an eigenbasis with respect to an observable. We also know that one may choose a common eigenbasis for any matrices that commute. If $[A,B]=0$, then the matrices $A$ and $B$ have the same eigenbasis $\ket{\xi_a}$. This means that those eigenstates can each have well-defined measurable properties under two different observables. 

\include{Zchap_LA_10_QM}
\include{Zchap_LA_11_fourier}
\include{Zchap_LA_12_function}
%!TEX root = Physics231.tex
\chapter{Spring Theory}
\label{ch:spring:theory}

\begin{quote}
\emph{Everything is a harmonic oscillator.}
\\ -- every physics student, eventually.\sidenote{At least approximately.}
\end{quote}


\section{Newton}

You know from your earliest physics memories that a harmonic oscillator is a system with a canonical kinetic term and a quadratic potential. The first time you saw it may have been in the context of, say, a spring\sidenote{Or a pendulum.} with some spring constant $k$. This means that the potential is $V(x)=(k/2)x^2$ where $x$ is the displacement from the equilibrium configuration that minimizes potential energy. 

The variable $x$ describes the state of a system. This state evolves in time, so it is a function $x(t)$. Physics tells us how this state evolves in time.\sidenote{This is too arrogant. What I really mean by `physics' is that our mathematical description of physics. Perhaps there is a different mathematical description wherein the evolution rule is not infinitesimal.} It does this by telling us how the state $x(t_0)$ at some time changes with an \emph{infinitesimal} change in time, $t_0 \to t_0 + \D{t}$. This is called the equation of motion. 

In Newtonian physics, the equation of motion is $F=m\ddot{x}$, where $F = -V'(x)$. This gives
\begin{align}
    \frac{ \D[2]{x(t)} }{ \D{t}^2 } 
    = 
    - \frac{ k }{ m } x(t) \ .
\end{align}
The rule for the infinitesimal change is \emph{second order}. Those who are well versed in differential equations will recall that this means you need first order initial data, like the initial velocity.\sidenote{The fancy way of saying this is Picard's theorem for the existence and uniqueness of ordinary differential equations. When I was an undergrad, I took the math-for-mathematicians first year calculus sequence. While my friends taking engineering math learned how to solve a handful of common differential equations, we spent 10 weeks learning Picard's construction to argue that solutions exist. If you wonder why I often make jokes about mathematicians and engineers, this experience is probably why.}


\section{Lagrange, traditionally}
\label{sec:lagrange:spring:theory}

You meet the harmonic oscillator again when upgrading from Newtonian to Lagrangian dynamics.\sidenote{For our purposes, we treat Lagrangian and Hamiltonian dynamics as the same up to a Legendre transform, see Appendix~\ref{ch:Lagrange:Legendre}.} 
% 
One of the fundamental principles in physics is the \textbf{principle of least action}\index{principle of least action}. This is the statement that classically, the trajectory of a system is governed by the minimum of a functional called the action, $S$. A functional is a ``function of functions.'' The position of a particle $\vec{x}(t)$ is a function of time. The action is a function of $\vec{x}(t)$. You may be more familiar with the action as the time integral of the Lagrangian,
\begin{align}
    S = \int \D{t}\, L[\vec{x}, \dot{\vec{x}}] \ ,
\end{align}
where we use the standard convention where we treat the position and velocity as independent variables in phase space. In quantum mechanics, this same object appears as a weight for possible trajectories.\footnote{I recommend looking up Feynman's phasor description of the sum over histories to see how the action in quantum mechanics leads to the principle of least action in classical mechanics.}


In analytical mechanics we derive the equation of motion by varying the Lagrangian. This is called the \textbf{variational priciple}\index{variational principle}. The standard treatment goes something like this. The Lagrangian of a classical particle is the kinetic minus the potential energy,
\begin{align}
    L = \frac{m}{2} \dot{\vec{x}}^2 - V[\vec{x}] \ .
\end{align}
When we treat $\vec{x}$ and $\dot{\vec{x}}$ as independent variables, the infinitesimal change (variation) of the action $S$ with with respect to infinitesimal perturbations $\delta\vec{x}$ and $\delta\dot{\vec{x}}$ is
\begin{align}
    \delta S =
    \int \D{t} \,
    \left[
    \frac{\partial L}{\partial \vec{x}} \delta\vec{x}
    +
    \frac{\partial L}{\partial \dot{\vec{x}}}\delta\dot{\vec{x}}
    \right] \ .
\end{align}
Now we `remember' that $\dot{\vec{x}}$ is actually the time derivative of $\vec{x}(t)$. This means we can integrate the second term by parts to remove the time derivative from $\delta\dot{\vec{x}}$. This is simply the statement that there is \emph{one} variation we're doing: the function $\delta\vec{x}(t)$. So if the path changes by $\vec{x}\to\vec{x}+\delta\vec{x}$, then the time derivative of the path changes by
\begin{align}
    \dot{\vec{x}} 
    \to \dot{\vec{x}} 
    + \frac{\D{}}{\D{t}}(\delta\vec{x}) 
    \equiv \dot{\vec{x}} + \delta\dot{\vec{x}} \ . 
\end{align}
% 
% 
We assume some boundary conditions---the position of the particle at some initial and some final times. This means we can integrate by parts because the boundary terms in the integration by parts vanishes.  We are thus left with (suppressing the limits of integration)
\begin{align}
    \delta S = \int dt\, 
    \left[
    \frac{\partial L}{\partial \vec{x}}
    - 
    \frac{\D{}}{\D{t}}
    \frac{\partial L}{\partial \dot{\vec{x}}} 
    \right]\delta\vec{x} \ .
\end{align}
Because this is true for \emph{any} variation $\delta\vec{x}$ about a path $\vec{x}(t)$ that minimizes $S$, we find that the equation of motion is 
\begin{align}
    \frac{\partial L}{\partial \vec{x}}
    - 
    \frac{\D{}}{\D{t}}
    \frac{\partial L}{\partial \dot{\vec{x}}}  = 0 \ .
\end{align}
For the single particle Lagrangian above, this indeed gives
\begin{align}
    m\ddot{\vec{x}} = V'[\vec{x}] \ .
\end{align}


\section{Linear Algebra of Variations}

If you are like me, you remember the first time you learned the variational principle. It is so unusual compared to Newtonian mechanics that you may even remember the second and third time you learned the variational principle---because the first time didn't sink in. Why would it? Lagrangian mechanics is the first time you meet a new mathematical technique in physics: the \emph{calculus of variations}.

In case the usual treatment carries too much of an air of mystery, we can strip it down to its simplest componentnts by working in discretized time. In contrast to the calculus of variations, let us call this the \emph{linear algebra of variations}.\sidenote{There is no such phrase, nor is this actually different from the traditional treatment. We simply want to understand the calculus of variations from the perspective of linear algebra.} In this section, we first give a lightning review of the calculus of variations, then rephrase it in the more prosaic language of linear algebra.

\paragraph{Review of the calculus of variations}
In ordinary calculus, we vary some independent variable $x$ which may take any value in $\RR$. We search for extrema of the dependent variable, say $f(x)$, by looking at stationary points where $\delta f(x) = f(x+ \delta{x}) - f(x) = 0$. In the {calculus of variations}, we vary a separate independent variable $x(t)$ for \emph{each value of $t$}. That means that $x(0~\text{sec})$ and $x(1.2~\text{sec})$. There are a \emph{continuum} of independent variables, $x(t)$, ``indexed'' by the continuum paramter $t$. Of course, this is simply the path of the particle: for each time $t$, $x(t)$ tells you the position of the particle. The action functional $S[x(t)]$ is a number associated to each path.

In this calculus of variations, the analog of derivatives\sidenote{Math purists, forgive the abuse of the differential here. Just imagine $\D{x} \defeq \Delta x$ in the $\Delta x \to 0$ limit.} 
\begin{align}
    \D{f}(x) = f(x+\D{x}) - f(x) 
\end{align}
is a variation,
\begin{align}
    \delta S[x(t)] = 
     S[x(t)+\delta x(t)] - S[x(t)] 
     \ .
\end{align}
Ordinarly in calculus we write $\D{f}/\D{x}$, but for now I do not want to dig into what the `denominator' of $\delta S[x(t)]/\delta x(t)$ means: unlike $\D{x}$ which physicists like to think of as the infinitesimal limit of $\Delta x \to 0$, $\delta x(t)$ is a function. What does it mean to ``divide by a function?''\sidenote{The answer is that there is an implicit sum over indices. To see this, first examine the case where $f$ depends on two variables.} For what it is worth, the idea that $\D{f}/\D{x}$ is a fraction is already something that we do that causes mathematicians anxiety.\sidenote{This is formally true in an extension of standard analysis called \emph{nonstandard analysis} wherein infinitsimal limits are replaced by \emph{hyperreal numbers} that encode the meaning of `infinitseimal.' I do not know of any example where thinking of $\D{f}/\D{x}$ as a fraction leads to a logical inconsistency---but you can find some insightful discussion online.}

There is also an analog of integrals. In ordinary calculus, the integral is a sum over infinitesimal changes with respect to successive changes in the independent variable. For example, if $h(x) = f'(x)$,
\begin{align}
    \int_{x_0}^{x_1} \D{x} \, h(x)
    =
    \int_{x_0}^{x_1} \D{x} \, \frac{ \D{f} }{ \D{x} } 
    = f(x_1) - f(x_0)
    \ .
\end{align}
Observe that we seem to be able to `cancel out' the infinitesimals $\D{x}$. This is the fundamental theorem of calculus. In the calculus of variations, the analog of an integral is a \textbf{path integral} which is a sum over successive changes in the paths $\delta x(t)$:
\begin{align}
    \int \mathcal{D}x\, Z[x(t)] \ .
\end{align}
We denote the integral measure over all paths by $\mathcal{D}x$ to distinguish it from the single variable measure $\D{x}$---the path integral measure encodes an \emph{infinite} number of single variable measures.\sidenote{Observe that the integrand $Z[x(t)]$ is \emph{not} the action. The name $Z$ should make you think of the partition function.} Path integrals do not usually show up in classical mechanics, but they are a key part of the Lagrangian formulation of quantum mechanics.

With all that being said, let us convert the above machinery into something much more mundane: the linear algebra of variations. The key idea is this: rather than dealing with an \emph{infinite continuum} of independent variables $x(t)$, let us deal with a discrete set of variables that we may arrange as vectors $\ket{x}$. This is analogous to our ``histogram basis'' of function space. The vector space may still be infinite dimensional,\sidenote{Or you can always take a finite limit, say with periodic boundary conditions.} but the discrete spacing means that our differential operators are now naturally interpreted as large matrices rather.




\paragraph{Discrete Time} To avoid doing calculus, we simply replace the continuous time variable $t$ with a discrete lattice of specific times $t_i$ where $t_{i+1} = t_i + \Delta t$ for some finite\sidenote{I have a colleague who finds the word `finite' here rather odd. Usually `finite' is in contrast to infinite, but here we use `finite' in contrast to \emph{infinitesimal}. I refer to Jim Holt's essay on infinitesimals in \emph{When Einstein Walked with G\"odel} for a bit of the history of this idea.} and constant $\Delta t$.

This means that the trajectory of a particle through time, $x(t)$ is now reduced to a discrete list of numbers $x^i$ corresponding to $x(t_i)$. This list of numbers may be infinite, but they are countably infinite. Because we have discretized, there is nothing `in between' $t_i$ and $t_{i+1}$. Instead, we say that physics tells us how $x^{i}$ and $x^{i+1}$ are related to each other. In Newtonian mechanics, we would say that you simply have to know $x^i$ and some notion of the velocity, say $\Delta t\inv(x^i-x^{i-1})$, and you could then deduce $x^{i+1}$ using the discrete version of $F=m\ddot{x}$.
\begin{exercise}
If you have not done this before, it is a worthwhile exercise to write pseudocode for how a computer would determine the trajectory of an asteroid under the influence of a gravitational force over small discrete time steps. You can pretend that you are an engineer in \acro{NASA} in the 1960s. Even as someone who does not often write `honest' code, I highly recommend the classic \emph{Numerical Recipes 3rd Edition: The Art of Scientific Computing}~\autocite{10.5555/1403886} to go deeper into this topic.
\end{exercise}

The principle of least action is rather different from Newtonian mechanics. The action is a functional of the state vector, $S[x^i]$.\sidenote{In fact, you can think of $S[x^i] = \la \vec{x}, \vec{x}\ra$ for an appropriately defined inner product.} The equation of motion comes from minimizing $S$ for all possible values of $x^i$. Because $S$ is an integral over the Lagrangian, and because the Lagrangian's kinetic term relates the values of $x^i$ to that of $x^{i+1}$, we see that minimizing $S$ ends up relating the position of the particle at consecutive times.
\begin{align}
    S_\text{kin} = 
    \sum_{i} \Delta t \,
    \frac{1}{2} 
    \left(
        \frac{ x^{i+1} - x^i }{ \Delta t }
    \right)^2 \ .
\end{align}
The continuum limit $\sum_i \Delta t \to \int \D{t}$ should be clear. The variational principle is the statement that the classical path is the one for which $S$ is minimized over all possible positions at all possible (discrete) times. This variation is simply:
\begin{align}
    \delta S = \frac{\partial{S}}{\partial{x^i}}\delta x^i
    = 
    \frac{\partial{S}}{\partial{x^1}}\delta x^1
    +
    \frac{\partial{S}}{\partial{x^2}}\delta x^2
    + 
    \cdots \ .
    \label{eq:S:partial}
\end{align}
When $S$ is extremized, $\delta S = 0$ for \emph{any} set of variations $\delta x^i$.  Observe that we use the partial derivative notation here: $S$ is no longer strictly a \emph{functional}, but just a function of a large number of independent variables $x^i$. 
\begin{exercise}
Show that for a theory with no potential energy, the $\delta S_\text{kin}=0$ for arbitrary $\delta x^i$ implies that nearest [temporal] neighbor states are equal $x^i = x^{i+1}$. 
\end{exercise}

Let us be clear that in this section the particle moves in a single dimension. The index $i$ on $x^i$ refers to a time stamp, not a spatial direction.\sidenote{We'll get there, don't worry. For now, you can appreciate why it makes sense that in relativity, space and time are treated on similar footing.} The vector 
\begin{align}
    \ket{x} &= x^i \ket{i} 
\end{align}
is a description of all positions of $x(t_i)$ over the temoral history of the particle. The coefficients $x^i$ can take any conitnuous value, but the time slices $t_i$ are discrete.\footnote{The upper index on $x^i$ is really a vector index in that it is understood to contract with, say, basis kets. The lower index on $t_i$ is simply a discrete label.} The basis vectors $\ket{i}$ refer to a unit displacement of the particle in the positive direction at time $t_i$.
\flip{Insert graphic here. Make a discrete time diagram.}




\paragraph{Discretized Action Principle}
We would like to determine the function, $\ket{x}=x(t)$ that minimizes $S[x(t)]$. We may write $S$ as a sum over a Lagrangian $L$. 
\begin{align}
    S &= \sum_i \Delta t \, L[x^i, \dot x^i]  \ .
    % &
    % \ket{x} &= x^i \ket{i}
    % &
    % \ket{\dot{x}} &= \dot{x}^i \ket{i} \ ,
\end{align}
The Lagrangian itself is a function of $t$---that is, it depends on the discrete time index $i$.
Ah, but wait a moment: what do we make of the $\dot x^i$? The dot is a time derivative, and our goal here was to \emph{not do calculus}. We need to pass to the discretized form of the derivative---recall our `histogram basis' introduction to function spaces. We would like to write $\dot x$ as matrix representation of the time derivative acting on the $x$ vector:
\begin{align}
    \ket{\dot{x}} = D\ket{x} \ .
\end{align}
But $D$ itself is ambiguous: we need to decide whetehr $D$ is the \emph{forward} or the \emph{backward} derivative,
\begin{align}
    D_+ \ket{f} 
    &\equiv \sum_i \frac{ f^{i+1} - f^i }{\Delta t} \ket{i}
    \\
    D_- \ket{f} 
    &\equiv \sum_i \frac{ f^{i} - f^{i-1} }{\Delta t} \ket{i} \ .
\end{align}
For our purposes, let us choose $D=D_+$. We drop the explicit index for now when the difference is innoucous, but we restore it when we have to make use of the non-Hermiticity of the derivative, $D_\pm^\dag = - D_\mp$. The definitions above are subject to boundary conditions which we leave implicit.\sidenote{Assume periodic boundary conditions if ambiguity makes you anxious.} 
\begin{exercise}
Confirm that all of our discussion in the rest of this section holds equivalently if we had chosen $D=D_-$.
\end{exercise}

Armed with a definition for the discretized derivative, we may explicitly write out the Lagrangian for the kinetic term:
\begin{align}
    L[x^i]_\text{kin} &= \frac{1}{2} , (Dx)^i (Dx)^i 
    &
    \text{no sum over }i \ .
\end{align}
Eh, this is something of a cumbersome notation. The vector components $(Dx)^i$ are the coefficients of $D\ket{x} = (Dx)^i\ket{i}$, but here we have repeated upper indices\sidenote{This is already a sin!} and a reminder that there is no sum. It is much cleaner to write the action directly:
\begin{align}
    S_\text{kin} = \frac{m}{2} g_{ij} (Dx)^i (Dx)^j \ ,
    \label{eq:S:kin:discrete}
\end{align}
where now we have a completely scalar quantity with indices contracted. In fact, we see that $S_\text{kin} = (m/2)\la Dx, Dx \ra$. We take the metric to be Euclidean so that $g_{ij} = 1$ if $i=j$ and zero otherwise. For the sake of preserving the index structure, we keep writing it explicitly as $g_{ij}$.\sidenote{It is dangerous to write some ``Kronecker'' $\delta_{ij}$ because the object does not generically transform like a (0,2)-tensor.} 
% 
Observe that unlike $L_\text{kin}$, $S_\text{kin}$ already has a sum over all time slices. 


\paragraph{Equation of motion for the free particle, linear algebra version} A free particle is one which has no potential.\sidenote{There is something poetic here along the lines of \emph{The Big Lebowski}.} This means the action is $S=S_\text{kin}$. You already know that the equation of motion in the continuum case is $\ddot{x}(t)=0$. Let us see how we get this in the discretized case.  The key insight is that even though the time slices $t_i$ are discrete, the components $x^i$ are continuous. Thus we may ``do calculus'' on these objects as usual. In this way, the central objects are the \emph{partial derivatives} in \eqref{eq:S:partial}.

% There is a trick to make the steps conceptually clearer, even though it makes everything a bit clunky. Since standard textbooks do everything in a slick way, let's do the clunky derivation. 

Let us start by rewriting $S$ in the form $\frac{1}{2}\la x | A x \ra$, where $A$ is some self-adjoint\sidenote{$A^\dag = A$} matrix. In fact, $A$ is real so that $A^\trans = A$. Why would we want to do this? 
\begin{theorem}\label{thm:quadratic:action}
If $S=\frac{1}{2} \la x | A x \ra$, then the variation of $S$ is
\begin{align}
    \frac{\partial S}{\partial x^i}
    = 
    g_{ij}A\aij{j}{k}x^k \ .
    \label{eq:variation:of:quadratic}
\end{align}
\end{theorem}
\begin{proof}
The action is
\begin{align}
    S= \frac{1}{2}  g_{\ell j}x^\ell A\aij{j}{k}x^k \ .
\end{align}
The product rule from calculus gives
\begin{align}
    2 \frac{\partial S}{\partial x^i}
    &= 
    g_{\ell j}
    \frac{\partial x^\ell}{\partial x^i} A\aij{j}{k}x^k
    + 
    g_{\ell j} x^\ell A\aij{j}{k}
    \frac{\partial x^k}{\partial x^i}
    \\
    &=
    g_{\ell j}
    \delta^\ell_i A\aij{j}{k}x^k
    + 
    g_{\ell j} x^\ell A\aij{j}{k}
    \delta^k_i 
    \label{eq:discrete:2:partial:s:partial:xi:int}
    \ .
\end{align}
The first term on the right-hand side is 
$g_{i j} A\aij{j}{k}x^k$. This is half of the final answer \eqref{eq:variation:of:quadratic}. The second term had better be the same as the first term. To see this, tecall the definition of the adjoint for real matrices, \eqref{eq:adjoint:def}: $(A^\dag)\aij{i}{j} = g_{jk}A\aij{k}{\ell}g^{\ell i}$. Then the self-adjoitness of $A$ gives
\begin{align}
    A\aij{j}{k} = (A^\trans)\aij{j}{k}
    = g_{km}A\aij{m}{n}g^{nj} \ .
\end{align}
Plugging this into the second term of \eqref{eq:discrete:2:partial:s:partial:xi:int} gives
\begin{align}
    g_{\ell j} x^\ell A\aij{j}{k}
    \delta^k_i 
    &=
    g_{\ell j} x^\ell 
    g_{km}A\aij{m}{n}g^{nj}
    \delta^k_i 
    =
    g_{im}A\aij{m}{n}
    g^{nj}
    g_{\ell j} 
    x^\ell 
    =
    g_{im}A\aij{m}{\ell}
    x^\ell \ .
\end{align}
Up to the names of dummy indices, this indeed matches the first term of \eqref{eq:discrete:2:partial:s:partial:xi:int}. Changing these dummy indeices gives the result, \eqref{eq:variation:of:quadratic}.
\end{proof}

What remains, then, is to show how to write \eqref{eq:S:kin:discrete} in the form of Theorem \ref{thm:quadratic:action}. This, in turn, is straightforward from the definitino of the adjoint. Recall that everywhere we wrote $D$ actin on a vector we really meant the \emph{forward} derivative, $D_+$. Then thea ction is
\begin{align}
    S_\text{kin} = \frac{m}{2} \la D_+x , D_+x \ra  \ .
\end{align}
We also know that the adjoint of the forward derivative is the backward derivative, $D_+^\dag =- D_-^\dag$. By the definitin of adjoint, this means that 
\begin{align}
    \la D_+x , D_+x \ra = 
    - \la x , D_- D_+x \ra \ .
\end{align}
Nowe we recognize that $D^2 \equiv D_- D_+$ \emph{is} a self-adjoint operator. We may thus identify $A =- \frac{m}{2}D^2$ in Theorem \ref{thm:quadratic:action} \ . By the result of that theorem, we find that
\begin{align}
    \frac{\partial S_\text{kin}}{\partial x^i}
    = 
    - \frac{m}{2} g_{ij} (D^2)\aij{j}{k}x^k \ .
\end{align}
This is the equation of motion for the free particle. The continuum version is
\begin{align}
    \frac{\delta S_\text{kin}}{\delta x(t)}
    = - \frac{m}{2}\ddot x(t) \ .
\end{align}
\begin{exercise}
Walk through this same derivation in the continuum limit. It may help to write the action as:
\begin{align}
    S_\text{kin} = 
    \int \D{t}\,
    \frac{ m }{ 2 }
    \left[ 
        \frac{ \D{y} }{ \D{t} }
        \frac{ \D{z} }{ \D{t} }
    \right]_{y(t), z(t) =x(t)} \ .
\end{align}
We do this to help keep track of terms in the product rule:
\begin{align}
    \frac{ \delta S_\text{kin} }{ \delta x(t') }
    &=
    \int \D{t}\,
    \frac{ m }{ 2 }
    \left[ 
        \frac{\delta}{\delta y(t')}
        \left(\frac{ \D{y} }{ \D{t} }\right)
        \frac{ \D{z} }{ \D{t} }
        +
        \frac{ \D{y} }{ \D{t} }
        \frac{\delta}{\delta z(t')}
        \left(\frac{ \D{z} }{ \D{t} }\right)
    \right]_{y(t), z(t) =x(t)} \ .
\end{align}
Here the $\delta/\delta x(t')$ is a variation at some fixed $t'$ in the same way that $\partial/\partial x^i$ is a varation of the path at some fixed $t_i$. Now the challenge is: how do we vary a term like $\D{y}/\D{t}$ with respect to $\delta y(t')$? We can integrate by parts so that
\begin{align}
    \int \D{t}\,
    \frac{\delta}{\delta y(t')}
        \left(\frac{ \D{y} }{ \D{t} }\right)
        \frac{ \D{z} }{ \D{t} }
    =
    -
    \int \D{t}\,
    \frac{\delta}{\delta y(t')}
        y(t)
        \frac{ \D[2]{z} }{ \D{t}^2 } 
        \ .
\end{align}
Then use the fact that $\delta y(t)/\delta y(t') = \delta(t-t')$ in the same way that $\partial x^i/\partial x^j = \delta^i_j$:
\begin{align}
    \int \D{t}\,
    \left[\frac{\delta}{\delta y(t')}
                \left(\frac{ \D{y} }{ \D{t} }\right)
                \frac{ \D{z} }{ \D{t} }\right]_{y,z = x}
    =
    - \frac{ \D[2]{x} }{ \D{t'}^2 } \ .
\end{align}
Fill in the remaining steps as needed.
\end{exercise}
The `linear algebra' approach to solving the quadratic is \emph{not} the one you see in analytical mechanics, but it is more convenient in field theory.





\paragraph{Equation of motion for the free particle, calculus version}
Let us redo the free particle with discretized time, but this time let us do it the way that an analytical mechanics textbook might do it. We start from \eqref{eq:S:kin:discrete} and rewrite it in terms of an intermediate variable $y^i$,
\begin{align}
    S_\text{kin} &= \frac{m}{2}g_{ij} y^i y^j
    &
    y^i\defeq D\aij{i}{k} x^k \ .
\end{align}
This variable $y(t)$ is what we typically call the velocity $\dot x(t)$.\sidenote{Remember we start by pretending that the velocity is independent of the path, then we remember that it is not.} By the chain rule, we have
\begin{align}
    \frac{ \partial S_\text{kin} }{ \partial x^i }
    =
    \frac{ \partial S_\text{kin} }{ \partial y^j }
    \frac{ \partial y^j }{ \partial x^i } \ .
\end{align}
The factors are simple:\sidenote{Fill in the steps.}
\begin{align}
    \frac{ \partial S_\text{kin} }{ \partial y^j }
    &= m g_{jk} y^k
    &
    \frac{ \partial y^j }{ \partial x^i }
    &=
    D\aij{j}{i}
    \ .
\end{align}
The second equation makes sense in the discrete case where $D$ is just some matrix, but it does not seem to make any sence in the continuum case where $D\to \D{}/\D{t}$ is a differential operator---what is it acting on? In order to make our expression sensible in the continuum limit, we should make sure $D$ is acting on something. Fortunately, index notation means we can move factors around:
\begin{align}
    \frac{ \partial S_\text{kin} }{ \partial x^i }
    &=
    m g_{jk} y^k
    D\aij{j}{i}
    =
    % m
    % D\aij{j}{i} g_{jk} y^k
    % =
    m 
    \delta^\ell_i
    D\aij{j}{\ell} g_{jk} y^k
    = 
    m g_{im}\, g^{m\ell} D\aij{j}{\ell} g_{jk} \, y^k \ .
\end{align}
where we use $\delta^\ell_i = g_{im} g^{m\ell}$. Now we recognize the transpose: 
\begin{align}
    g^{m\ell} D\aij{j}{\ell} g_{jk} \equiv (D^\trans)\aij{m}{k}
    \ .
\end{align}
Because we have assumed $D=D_+$, we  have $D^\trans = - D_-$. As a result, we have
\begin{align}
    \frac{ \partial S_\text{kin} }{ \partial x^i }
    = - m g_{im} (D_-)\aij{m}{k} (D_+)\aij{k}{j} x^j \ .
\end{align}
As before, we know that the product $D_-D_+ = D^2$ is the self-adjoint Laplacian operator so that we recover the usual $-m\ddot{x} = 0$ equation of motion. 

\begin{exercise}
Show that the continuum limit of this procedure is exactly what we do in Section~\ref{sec:lagrange:spring:theory}
\end{exercise}



\section{Discrete Time Harmonic Oscillator}

The harmonic oscillator action has a potential term that is quadratic in the position variable. In the discrete time limit,
\begin{align}
    S = 
        \frac{m}{2} g_{ij}(Dx)^i(Dx)^j 
        - \frac{k}{2} g_{ij} x^i x^j \ .
\end{align}
The relative minus sign ensures that there is a stable minimum energy configuration.
\begin{exercise}
Using complete sentences, explain why this is the most generic form of an action that is quadratic in $x$. Use homogeneity and isotropy in time to argue why there can be no tensor structures other htan $g_{ij}$.
\end{exercise}

Appreciate that the \emph{quadratic} part of the action gives a linear equation of motion. This is obvious because varying a quadratic action with respect to $x^i$ gives terms that are linear in the position variable. This means that higher-order terms, e.g.\ trilinear terms like $(x^i)^3$, must be treated separately.\footnote{Often we treat these as a small correction to the quadratic action. This amounts to saying that everything is a perturbation on a harmonic oscillator. The perturbations can be treated systematically using a technique called Feynman diagrams.}
% \begin{exercise}
% Confirm that \eqref{eq:HO:action:discretized} is the discretization of the action for a simple harmonic oscillator. For example, show that in the continuum limit you recover the simple harmonic oscillator action. 
% \end{exercise}
% 
% 

% \flip{This next exercise is probably wrong: I think the point is that $D_+^\dag = D_-$ so that we can define the forward derivative for vectors and the backward derivative for forms. The action takes the form $(D_+q)^2$ so that the equation of motion is $D_-D_+ q=0$.}
% \begin{exercise} \label{ex:integrate:kinetic:term:by:parts}
% If $x(t)$ is the position of a particle of mass $m$, then its kinetic term is $m\dot x(t)/2$. One way to derive the equation of motion for a theory is to integrate this term by parts, see Section~\ref{eq:variational}. Let us see how this works. Rather than $\dot x^2$, let us consider $(D_+f)(D_-g)$. We write $f$ and $g$ to make it easier to keep track of terms, but for the specific case of the kinetic term we know $f=g=x(t)$. 

% The integral over time corresponds to a sum over the time slice (histogram) index $i$ of $f$ and $g$:
% \begin{align}
%     \int dt\, \dot x(t)^2 \to \sum_i \Delta x \, (D_+f)^i(D_-g)^i \ .
% \end{align}
% Show that this sum takes the form
% \begin{align}
%     \sum_i \Delta x \, (D_+f)^i(D_-g)^i 
%     = -\sum_i f^i (D^2 g)^i 
%     + f^{N+1}(g^N-g^{N-1})
%     - f^1(g^1 - g^0) \ .
% \end{align}
% I was sloppy with the limits of the sum above. What values does $i$ range over in that sum?
% % HINT:
% \acro{Partial answer:}
% Here's the pattern:
% \begin{align}
%     D_+fD_-g &= 
%     f^2g^1 - f^2g^0 - f^1 g^1 + f^1 g^0\\
%     &+ f^3 g^2 - f^3 g^1 - f^2 g^2 + f^2 g^1 \\
%     &+ f^4 g^3 - f^4 g^2 - f^3 g^3 + f^3 g^2 \\
%     &+ \cdots \\
%     & + f^{N+1} g^N - f^{N+1} g^{N-1} - f^N g^N + f^N g^{N-1}  \ .
% \end{align}
% \end{exercise}
% 
% 

% Recalling Exercise~\ref{ex:integrate:kinetic:term:by:parts}, we may integrate the derivative term by parts to obtain

From the previous section, we are now confident that we can integrate the first term by parts---in fact, we saw how to interpret this in the discrete limit two ways---so that
\begin{align}
    S = 
    -
    \frac{m}{2} g_{ij}x^i (D^2x)^j 
    - \frac{k}{2} g_{ij} x^ix^i
    \equiv
    -\frac{1}{2}x^i Q_{ij} x^j
    \ ,
    \label{eq:HO:action:discretized:integrated}
\end{align}
where we have defined the `quadratic term' operator
\begin{align}
    Q_{ij} = 
    m\, g_{ik}(D^2)\aij{k}{j} 
    + 
    k\, g_{ij} 
    = Q_{ji} \ .
    \label{eq:Qij}
\end{align}
\begin{exercise}
Confirm that $Q_{ij}$ is symmetric, $Q_{ij} = Q_{ji}$. 
\end{exercise}
\begin{exercise}
What is the \emph{support} of \eqref{eq:Qij}? That is, what is the condition on $i$ and $j$ to find a non-zero $Q_{ij}$?
\end{exercise}

The principle of least action posits that the classical trajectory $x(t)$---encoded in the discretized $x^i$ coefficients---is the configuration that extremizes $S$. To do this, we vary each $x^i$ independently. This is what the functional variation $\delta x(t)$ means in the continuous limit. The extremization corresponds to the $N$ equations
\begin{align}
    \frac{\partial{S}}{\partial{x^1}} &= 0
    &
    \frac{\partial{S}}{\partial{x^2}} &= 0
    &
    \cdots&
    &
    \frac{\partial{S}}{\partial{x^N}} &= 0 \ .
\end{align}
where relevant, these are subject to the boundary conditions. 
% 
Of course, the $x^i$ are not independent: we know that the kinetic term relates the displacement at time $i$ to the neighboring displacements at time $i\pm 1$.
% 
Rather than solving each of these one by one, we can solve the equation for a generic $x^i$:
\begin{align}
    -\frac{ \partial{S} }{ \partial{x^i} } &= 
    \frac{ \partial }{ \partial{x^i} }
    \left[
    \frac{1}{2} Q_{ii}(x^i)^2 + \sum_{j\neq i} \frac{1}{2}\left(Q_{ij} x^ix^j + Q_{ji}x^jx^i + \right) 
    \right] \\
    &= 
    Q_{ii}x^i+ \sum_{j\neq i} Q_{ij}x^j \ .
\end{align}
where there is no sum over repeated $i$ indices and we have explicitly written the sum over $j\neq i$. In the last term we have used $Q_{ij} = Q_{ji}$. Expanding this expression gives
\begin{align}
    0&=
    -\frac{\partial{S}}{\partial{x^i}}
    \\
    &=
    -2\frac{m}{\Delta t^2} x^i + k x^i 
    + \frac{m}{\Delta t^2}x^{i+1} 
    + \frac{m}{\Delta t^2}x^{i-1}
    \\
    &= 
    g_{ij}\left[m(D^2x)^j + k x^j\right] \ .
\end{align}
This gives us the discretized equation of motion:
\begin{align}
    (D^2x)^j = - k x^j \ ,
\end{align}
which we recognize is the discretized Hooke's law $m\ddot{x} = -k x$. 

What is significant about this approach is that we never invoked the standard Euler--Lagrange trick of first treating $x$ and $\dot{x}$ as independent variables and then `remembering' that they are related by a time derivative when it was convenient. We took a more direct path where we independently varied each piece of the path in time, $x^i$, while treating a differential operator $Q_{ij}$ as a matrix. We noticed that the $x^i$ are \emph{coupled} by the kinetic term, which tells us that wiggles in $x^i$ will affect the variation of $x^{i+\pm 1}$. In quantum mechanics, this approach is known as the path integral formulation.

\begin{example}
Having established the variational principle approach to the equation of motion ,we may repeat the derivation of Hooke's law by starting with the continuum limit. The action is
\begin{align}
    S&= \int_{t_0}^{t_1} dt\, \frac{m}{2}\dot x(t)^2 - \frac{k}{2}x(t)^2
    \\&
    = -\int_{t_0}^{t_1} dt\, \frac{1}{2}x(t)
    \left[m\left(\frac{d}{dt}\right)^2 + k\right]x(t)
    \\&    
    \equiv -\int_{t_0}^{t_1} dt\, \frac{1}{2} x(t) \mathcal O x(t) \ .
\end{align}
The classical equation of motion is thus
\begin{align}
    \mathcal Ox(t) = \left[m\left(\frac{d}{dt}\right)^2 + k\right]x(t) = 0 \ m \ ,
\end{align}
or in other words: $\ddot x(t) = -kx(t)$.
\end{example}
We have implicitly assumed that the boundary terms vanish.

On physical principles, the kinetic term is typically proportional to a term with two time derivatives. You can argue this based on spacetime symmetry. This means that one must perform an integration by parts on this term to write the equation of motion operator $\mathcal O$. This may give some insight on the meaning of the relative minus sign between the kinetic and potential energy terms in the Lagrangian. 




\section{Coupled Harmonic Oscillators}


Let's take stock of what we have accomplished. We considered a particle that can move in one direction of space subject to a harmonic oscillator potential. We discretized time to dissect the meaning of the variational principle. There is a trivial extension if the particle can move in multiple spatial directions: simply replace $x(t) \to \vec{x}(t)$ where $\vec{x}$ is a \emph{spatial} vector.\sidenote{Do not confuse this with $\ket{x}$ which is a vector with respect to the lattice of time slices.} 
\begin{exercise}
Work out the equation of motion for a particle in three dimensions of space subject to a harmonic oscillator potential. The spring constant $k$ is upgraded to a symmetric matrix $k\aij{i}{j}$, which you may take to be $k\delta^i_j$ for simplicity. Be sure to distinguish between $\partial/\partial x^i$, $\partial/\partial y^i$, and $\partial/\partial z^i$.
\end{exercise}
\begin{exercise}
What happens when $k\aij{i}{j}$ is not proportional to $\delta^i_j$?
\end{exercise}

We now want to move from single particles to fields. 
A particle has a single position at each time. Actually, we need to be more precise with our language. Rather than talking about `positions' of a particle, we should talk about a \emph{displacement from equilibrium.} So rather than talking about $x(t)$, we write $q(t)$.\sidenote{For example, we may not be tracking an actual position of a particle, but the temperature at a single point.} In contrast, a \textbf{field}\index{field} is an object that has some displacement at each time \emph{and at each position}. 

In the discretized limit, rather than just having a lattice of time slices, we have a lattice of space and time. For a given $x_a$ and $t_i$. The basis vectors are $\ket{a,i} = \ket{a}\otimes \ket{i}$: a unit displacement at position $x_a$ and time slice $t_i$. A state is $\ket{q} = q^{ai}\ket{a,i}$. 
% 
Rather than going through everything with both discrete space and discrete time, let's be lazy and return to continuum time. 

Consider an array of $N$ harmonic oscillators. The displacement at time $t$ of each oscillator $i$ with respect some baseline value to be $q_i(t)$. If each harmonic oscillator is independent, then each one would have a separate equation of motion---this is essentially $N$ copies of the single harmonic oscillator. Let us instead consider the case where
\begin{enumerate}
    \item The point $q_i(t)=0$ is not special and does not represent an equilibrium value. In fact, let us throw out the $V[q_i]=\frac{1}{2}kq_i^2$ potential energy term and replace it with something else.
    \item Let us couple the harmonic oscillator to the positions of its \emph{nearest neighbors}. The nearest neighbors of the $q_i$ oscillator are the $q_{i\pm 1}$ oscillators. 
\end{enumerate}
The potential energy of the system looks like
\begin{align}
    V[q_1,\cdots q_N] &=
    \cdots
    \frac{1}{2}k(q_{i+1}-q_i)^2 +
    \frac{1}{2}k(q_{i}-q_{i-1})^2 +
    \cdots \ .
\end{align}
This simply means that the locations of the $i^\text{th}$ oscillator's neighbors will pull the $i^{\text{th}}$ oscillator in one direction or the other. There is \emph{no} potential that pulls $q_i$ to its origin: $(k/2)q_i^2$. 


The Lagrangian  for this system is
\begin{align}
    L[q_1,\cdots, q_N]
    &= 
    \sum_i \frac{1}{2} m \dot q_i^2 - \frac{1}{2}k (q_{i+1}-q_i)^2 \ .
\end{align}
You can work out the equation of motion simply by varying the action, $S=\int dt\,L$. For example, the equation of motion for the $q_i(t)$ oscillator is
\begin{align}
    m\ddot{q}_i(t) = -k(q_{i}-q_{i-1}) -k(q_{i}-q_{i+1}) \ .
\end{align}
We observe that the spring constant terms push to decrease the difference between $q_i$ and its nearest neighbors. The force from the neighboring springs are in the direction to reduce the difference. In the continuum limit,
\begin{align}
    m\ddot{q}(t,x) = k\frac{\partial^2{q(t,x)}}{\partial{x}^2} \ .
\end{align}
We can write this as the d'Alembertian,
\begin{align}
    0 = m \partial^2 q \equiv m\left(\partial_t^2 - c^2 \partial_x^2\right) q(t,x) \ ,
\end{align}
where we recognize the speed of propagation $c^2 = k/m$.





\section{Feynman}
% Classical to quantum
% \flip{Classical to quantum, not just minimum, but the entire ensemble. }

The \textbf{path integral}\index{path integral} is something that shows up in the `sum over histories' approach to quantum mechanics. It is used extensively in quantum field theory because it gives a powerful way to impose constraints on a quantum system using Lagrange multipliers. What is not always obvious is that the path integral can be a more intuitive interpretation of the standard Euler--Lagrange approach. 

\paragraph{Sum Over Histories}
% 
\flip{To be filled in. Please see the example in class of the $N$-slit experiment and the proposal that each path contributes $\e^{iS}$.}

\begin{newrule}[Sum over histories]
Given some observed initial state $\ket{\text{in}}$ and some observed final state $\ket{\text{out}}$, the amplitude to go from $\ket{\text{in}}$ to $\ket{\text{out}}$ is
\begin{align}
    \mathcal M(\text{in}\to\text{out})= 
    \sum_\textnormal{path} \e^{iS_\textnormal{path}} \ ,
\end{align}
where the sum is over all `paths' (histories) that permit the initial state to evolve into the final state. $S_\textnormal{path}$ is the action for a given path---the same action that you are used to from classical mechanics, the time integral over the Lagrangian, $S= \int \D{t}\, L$.
\end{newrule}

\paragraph{Path Integral}
Suppose a quantum state has initial and final positions\sidenote{We deliberate use the notation $q$ for position instead of $x$. This is to aid our transition to quantum \emph{field theory}. Old textbooks sometimes refer to \emph{second quantization}; it is precisely to distinguish the role of $q$, a field displacement, and $x$ a continuous field index indicating spacetime position.}
\begin{align}
    \ket{\textnormal{in}} &= \ket{q_0}
    &
    \ket{\textnormal{out}} &= \ket{q_T} \ .
\end{align}
These are labeled so that we observe $\ket{q_0}$ at time $t=0$ and we observe $\ket{q_T}$ at time $t=T$. The amplitude to go from $\ket{q_0} \to \ket{q_T}$ over the time interval $T$ is
\begin{align}
    \bra{q_T} \hat U(T) \ket{q_0} &= \bra{q_T} \e^{-iT\hat H} \ket{q_0} \ .
\end{align}
Here $\hat U(t)$ is the operator that evolves states over time $t$. As you know from quantum mechanics, the Hamiltonian $\hat H$, is the operator that enacts infinitesimal evolution in time. This means that we enact finite time translations by exponentiating the Hamiltonian.\sidenote{Please make sure you are familiar with this idea: Hermitian operators \emph{generate} infinitesimal changes to a system. The exponential of these Hermitian operators are unitary operators that enact finite transformations. The idea shows up over and over again in this class.}

There are two tricks to deploy at this step:
\begin{enumerate}
    \item Break the finite time evolution operator into a product of $(N+1)$ time evolutions:
    \begin{align}
        \hat U(T) = \hat U(\Delta t)\hat U(\Delta t)\cdots \hat U(\Delta t) \ .
    \end{align}
    It is \emph{obvious} that $\hat U(t)\hat U(s) = \hat U(t+s)$.
    \item Insert the identity as a sum over a complete set of intermediate states:
    \begin{align}
        \one &= \int \D{q_i}\, \ket{q_i}\!\bra{q_i} \ .
    \end{align}
    We insert a copy of the identity in-between each of the $U(\Delta t)$s. Thus the label $i$ runs from $1$ to $N$.
\end{enumerate}
\begin{exercise}
Confirm that $\hat U(t)\hat U(s) = \hat U(t+s)$. \textsc{Hint}: this is trivial, but convince yourself that it is trivial.
\end{exercise}

Applying these two tricks to our amplitude\sidenote{$\D[N]{\vec{q}} \defeq \D{q_1}\cdots\D{q_N}$} 
\begin{align}
    \bra{q_T} \hat U(T) \ket{q_0} 
    &= \int \D[N]{\vec{q}} \,
    \bra{q_T} \hat U(\Delta T) \ket{q_N}
    \bra{q_N}\hat U(\Delta T) \ket{q_{N-1}}\cdots 
    \\
    &=\int \D[N]{\vec{q}} \,
    \prod_i^{N+1} K_i \ , \label{eq:path:integral:time:evolution:product:of:Ks}
\end{align}
where the factor $K_i$ is a complex number:
\begin{align}
    K_i &\defeq \bra{q_{i}} \hat U(\Delta T) \ket{q_{i-1}} \ ,
\end{align}
where we define $\bra{q_{N+1}} = \bra{q_T}$. We have reduced the amplitude to go from our initial state to our final state to a product of $K_i$. The form of $K_i$ is simply that of the amplitude to go from $\ket{q_{i-1}}$ to $\ket{q_{i}}$ over a small time $\Delta t$. The integral over intermediate states $\D[N]{\vec{q}}$ tells us that we are summing over all possible configurations of those unobserved intermediate states. 

In order to evaluate $K_i$, let us remind ourselves of the form of the Hamiltonian in $U(\Delta T) = \exp(i\Delta T\hat H)$,
\begin{align}
    \hat H = \frac{\hat P^2}{2m} + V(\hat Q) \ ,
\end{align}
where we use the non-relativistic kinetic term $\frac{1}{2}m\dot q^2 = p^2/2m$ and an arbitrary potential $V(q)$ that only depends on the $q$ displacements, but not the momenta $p$. Because $\hat H$ is an operator, we write $\hat P$ for the momentum operator and $\hat Q$ for the position\sidenote{This `position' is not necessarily spatial position. In general it is a displacement of a field value.} operator. All Hamiltonians in ordinary field theory have this type of separation: there is a \emph{kinetic term} that depends on momentum and a \emph{potential term} that depends on displacements. When we pass to relativistic theories, the kinetic term converts into a form that includes spatial derivatives in a way that it is Lorentz invariant. 

Once again there is another silly trick. 
\begin{enumerate}
    \setcounter{enumi}{2}
    \item Insert the momentum-space representation of the identity,
    \begin{align}
        \one &= \int \Dbar{p}\, \ket{p}\!\bra{p} 
        &
        \Dbar{p}\defeq \frac{\D{p}}{2\pi} \ .
    \end{align}
    Observe that this has a relative factor of $2\pi$ compared to the resolution of the `position space' identity. This is the famous---and seemingly inscrutable---factor that shows up in the Fourier transform, see Appendix~\ref{app:Fourier}.\sidenote{If you spend time thinking about this, you may wonder why we do not split the $2\pi$ evenly between positions and momenta. This would certainly be symmetric, but this choice ends up being fairly convenient in physics. It means, for example, that the conjugate variable for time is essentially the angular momentum, $E=\omega$.}\footnote{See e.g.\ \url{https://physics.stackexchange.com/a/141737/166736}} 
\end{enumerate}
Because $\hat H$ contains a term that evaluates easily in momentum space and a term that evaluates easily in position space, we insert the complete set of momentum states into $K_i$ to allow the $\hat P$-dependent terms to hit a $\bra{p}$ while the $Q$-dependent terms hit a $\ket{q_i}$:
\begin{align}
    K_i &=
    \int \Dbar{p} \la {q_{i}} | p \ra \bra{p} \e^{-i\Delta t\frac{ \hat P^2}{2m}} \e^{-i\Delta tV(\hat Q)} \ket{q_{i-1}}
    \label{eq:path:integral:Ki:operator}
    \\ &=
    \int \Dbar{p} \, \e^{-i\Delta t\frac{ p^2}{2m}} \e^{-i\Delta tV(q_{i-1})}\,  \la {q_{i}} | p \ra \la p | q_{i-1} \ra  \ .
    \label{eq:path:integral:Ki:eigen}
\end{align}
The factors on the right are simply the Fourier components between the position and momentum eigenstates:
\begin{align}
    \la {q_{i}} | p \ra &= \e^{-ipq_i}
    &
    \la p | q_{i-1} \ra &= \e^{+ipq_{i-1}}
    \ .
\end{align}
Plugging this into $K_i$ gives
\begin{align}
    K_i &=
    \e^{-i\Delta tV(q_{i-1})} \int \Dbar{p} \, \e^{-i\Delta t\frac{ p^2}{2m} - i p(q_i - q_{i-1}) } \ .
\end{align}
If you do not worry\sidenote{%
    Will not worry. On the other hand, you may want to worry about whether the assumptions we made to solve these Gaussian integrals is still valid when there are complex coefficients. This scratches the surface of why some mathematicians are skeptical about whether anything we do in quantum field theory is well posed.} 
about the factors of $i$, then this integral is precisely in the form of a Gaussian with a source term. This has a known solution:
\begin{align}
    G_{a,J} \equiv \int \D{x}\; e^{-\frac{a}{2}x^2 + Jx} = \sqrt{\frac{2\pi}{a}}  e^{\frac{J^2}{2a}} \ .
    % \tag{\ref{eq:Gaussian:with:source}}
\end{align}
We identify
\begin{align}
    a &= \frac{i\Delta t}{m}
    &
    J&= -i(q_i - q_{i-1}) \ .
\end{align}
Plugging in this result gives
\begin{align}
    K_i &=
    \e^{-i\Delta tV(q_{i-1})} 
    \frac{1}{2\pi}
    \sqrt{\frac{2\pi m}{i \Delta t}}
    \e^{\frac{-(q_i-q_{i-1})^2}{2i\Delta t/m}}
    \\
    &=
    \sqrt{\frac{m}{i 2\pi \Delta t}}
    \e^{\frac{i}{2}m\frac{(q_i-q_{i-1})^2}{\Delta t^2}\Delta t}
    \e^{-iV(q_{i-1})\Delta t} 
    \\
    &= \mathcal N \e^{i\Delta t\, L[q_i]}
    \label{eq:path:integral:infinitesmimal:kernel}
     \ .
\end{align}
In the last line we have identified $(q_i-q_{i-1})/\Delta t = \dot q_i$ and observed that the argument of the combined exponential is simply $i\Delta t$ times the \emph{Lagrangian} at point $i$. The ugly square root is some complex constant factor that \emph{does not matter}. We simply write it as a normalization $\mathcal N$. 

Now we insert this into the amplitude to go from our initial state to our final state, \eqref{eq:path:integral:time:evolution:product:of:Ks}:
\begin{align}
    \bra{q_T} \hat U(T) \ket{q_0} 
    &= \mathcal N^{N+1} \int \D[N]{\vec{q}} \,
    \e^{i \sum_i \Delta t\, L[q_i]} \ .
\end{align}
Recalling that $i$ encodes `time slices,' we recognize that the argument of the exponential is simply an \emph{integral} of the Lagrangian,
\begin{align}
    i \sum_i \Delta t\, L[q_i] \to i\int \D{t}\, L[q(t)] \equiv iS[q(t)] \ ,
\end{align}
where we have recovered the action. In the continuum limit, it is conventional to replace the notation $\D[N]{\vec{q}}$ with $\mathcal{D}q(t)$ so that our final expression for the $\Delta t \to 0$ limit is
\begin{align}
    \bra{q_T} \hat U(T) \ket{q_0} 
    &= \mathcal N' \int \mathcal{D}q(t)\;
    \e^{i S[q(t)]} \ .
    \label{eq:path:integral:formalism}
\end{align}
Here $\mathcal N'$ is a re-packaging of the powers of $\mathcal N$. The integration over all trajectories $\mathcal {D}q(t)$ is called a \textbf{path integral}\index{path integral}.

\begin{exercise}
Why was it necessary for us to split $U(T)$ into several small steps? \textsc{Hint}: did we make any assumptions to go from \eqref{eq:path:integral:Ki:operator} to \eqref{eq:path:integral:Ki:eigen}? 
\end{exercise}

\paragraph{Recovering Lagrange}
% \flip{To be filled in. Discussion of phasor diagrams, steepest descent, and why the classical path dominates. Motivates the principle of least action as the $\hbar \to 0$ limit.}
The path integral formulation limits to the classical equation of motion in the limit $\hbar \to 0$. 
\begin{exercise}
Show that the $\hbar \to 0$ limit reduces to the standard Euler--Lagrange equation of motion for a classical system. You can argue heristically using a tool called \textbf{phasor diagrams}, though a more technical derivation will take the leading order terms in an expansion in $\hbar$.
\end{exercise}


\paragraph{Recovering Schr\"odinger}
This section is outside the primary scope of the class. However, I encourage you to at least skim it to see an alternative explanation of where the Schr\"odinger equations comes from. It should convince you that the sum over histories (path integral) construction gives you the same dynamics as before. 
% 
% 
We show that the Schr\"odinger equation appears as the single-particle, non-relativistic limit over the path integral formalism.\sidenote{This section draws from the excellent pedagogical article by David Derbes in the \emph{American Journal of Physics} (one of my favorite publications) in 1996. This artcile was selected as one of the Fermat's Library journal club articles.\footnotemark}\footnotetext{\cite{10.1119/1.18114}; \url{https://fermatslibrary.com/s/feynmans-derivation-of-the-schrodinger-equation}} The power of the path integral formalism is that it generalizes beyond non-relativistic quantum mechanics. Like with all new theoretical formalisms, it behooves one to confirm that the new formalism is mathematically equivalent to the established formalism in the appropriate limit.

Suppose we interpret $q(t)$ as the position of a non-relativistic particle at time $t$. Then in this basis, we can think of the amplitude $\la q_t | \hat U(t) | q_0 \ra$ as the wavefunction $\psi(q, t)$ for the particle to be in position $q$ at time $t$ given some initial state $\ket{q_0}$ at $t=0$. We can rewrite \eqref{eq:path:integral:time:evolution:product:of:Ks} as follows:
\begin{align}
    \psi(q,t) &= \la q_t | \hat U(t) | q_0 \ra 
    \\&= \int \D{q_N} \, K_{\!N} \int \D[N-1]{\vec{q}} \,
    \prod_i^{N} K_i 
    \\
    &= \int \D{q'} \, K_{\!N}\, \psi(q',t-\Delta t) \ .
\end{align}
We have relabeled $q_N \to q'$. It is the last displacement variable that we integrate over. In fact, for convenience,\sidenote{You do not have to do this step.} let us shift all of our time coordinates by $\Delta t$: 
\begin{align}
    \psi(q,t+\Delta t) 
    &= \int \D{q'} \, K_{\!N}\, \psi(q',t) \ .
\end{align}
This equation says that to go from the $(N-1)^\textnormal{th}$ step to the $N^\textnormal{th}$ step, you integrate the wavefunction at last position over all possible intermediate steps, weighted by the \emph{kernel},\sidenote{Kernels are just fancy names for distributions, which are in turn fancy names for functions that are integrated over.} $K_N$. We can now use the form of the kernel, \eqref{eq:path:integral:infinitesmimal:kernel}:
\begin{align}
    K_{\!N} &= \mathcal N \e^{i\Delta t\, L[q_i]}
    &
    \mathcal N = \sqrt{\frac{m}{2\pi i \Delta t}} \ .
    \label{eq:path:integral:normalization}
\end{align}
Let us now be careful with the scaling with the small time slice $\Delta t$:
\begin{align}
    i\Delta t L[q_i] = \frac{im}{2}\frac{(q_i-q_{i-1})^2}{\Delta t} - i\Delta t V(\bar q) \ .
\end{align}
Note that the kinetic term goes like $\Delta t^{-1}$, which is \emph{large}. The potential term scales linearly with $\Delta t$, so it is small. We write $V(\bar q)$ to mean the average value of $V$ in the interval $(q_{i-1}, q_i)$. For our purposes, $\bar q = q$ assuming a sufficiently smooth potential. This means we can Taylor expand the potential term, but not the kinetic term:
\begin{align}
    \psi(q,t+\Delta t) &= \mathcal N\int \D{q'}\; 
    \e^{\frac{im}{2}\frac{(q-q')^2}{\Delta t}} 
    \left[1 - i\Delta t \, V(\bar q)\right]\,
    \psi(q',t) \ .
\end{align}
We drop higher order terms in $\Delta t$.
Now we assume that $q= q' + \xi$ and make the plausible case that $\xi$ is small.\sidenote{To be clear, we integrate over even large values of $\xi$. The sin that I am sweeping under the rug is the idea that the rapidly varying oscillation in the $\e^{\textnormal{kinetic}}$ term will cause integrants with large values of $\xi$ to cancel. This is the steepest descent method in the method of phasors.} This means we can also replace $\bar q \to q$ to leading order.
We can then expand $\psi(q',t)$, where we shall suppress the $t$ argument so that spatial derivatives are clear:
\begin{align}
    \psi(q') = \psi(q) -\xi \psi'(q) + \frac{1}{2}\xi^2 \psi''(q) + \cdots \ .
\end{align}
To be clear: we now write $\psi'(q,t)$ to mean $\frac{\D{}}{\D{q}}\psi(q,t)$.
Plugging this in:
\begin{wide}
\begin{align}
    \psi(q,t+\Delta t) &= 
    \mathcal N\int \D{\xi}\; 
    \e^{\frac{im}{2}\frac{\xi^2}{\Delta t}} 
    \left[1 - i\Delta t \, V(q)\right]\,
    % \\
    % &\phantom{=
    % \mathcal N\int \D{\xi}\; 
    % }\times 
    \left[\psi(q,t) -\xi \psi'(q,t) + \frac{1}{2}\xi^2 \psi''(q,t)\right] 
    \\
    % &= 
    % \mathcal N\int \D{\xi}\; 
    % \e^{\frac{im}{2}\frac{\xi^2}{\Delta t}} 
    % \left[1 - i\Delta t \, V(\bar q)\right] \psi(q,t)
    % \\
    % &\phantom{=}+
    % \mathcal N\int \D{\xi}\; 
    % \e^{\frac{im}{2}\frac{\xi^2}{\Delta t}} 
    % \frac{1}{2}\xi^2 \psi''(q,t) \ .
    &= 
    \mathcal N\int \D{\xi}\; 
    \e^{\frac{im}{2}\frac{\xi^2}{\Delta t}} 
    \left[\psi(q,t) - i\Delta t \, V(q)\,\psi(q,t)
    + \frac{1}{2}\xi^2 \psi''(q,t)
    \right]
    \ .
 % \ . \nonumber
 \label{eq:psi:path:integral:schrodinger:intermediate}
\end{align}
\end{wide}
In the last line we removed the term linear in $\xi$ because it cancels against the rest of the integral, which is even in $\xi$. We also remove terms of order $\mathcal O(\Delta t\,\xi^2)$ since these are small. Observe that the $\xi$-independent terms in the bracket are simple Gaussian integrals:
\begin{align}
    \int \D{\xi}\; 
    \e^{\frac{im}{2}\frac{\xi^2}{\Delta t}} 
    = 
    \sqrt{\frac{2\pi i\Delta t}{m}} = \mathcal{N}^{-1} \ .
    \label{eq:path:integral:Gaussian:trick:1}
\end{align}
Observe that this factor exactly cancels the annoying normalization in \eqref{eq:path:integral:normalization}---but only for the $\xi$-independent terms in the bracket. The remaining term is a moment of the Gaussian integral. There happens to be a convenient trick for this by taking a derivative of 
\begin{align}
    G_{a,J} \equiv \int \D{x}\; e^{-\frac{a}{2}x^2 + Jx} = \sqrt{\frac{2\pi}{a}}  e^{\frac{J^2}{2a}} 
    \tag{\ref{eq:Gaussian:with:source}} \ ,
\end{align}
where we see that
\begin{align}
    -2\frac{\D{G_{a,0}}}{\D{a}} &= \int \D{x}\; x^2 e^{-\frac{a}{2}x^2} 
    = \frac{1}{a}\sqrt{\frac{2\pi}{a}} \ .
    \label{eq:Gaussian:trick:diff:quad}
\end{align}
\begin{exercise}
While the method above reconstructs even moments of the Gaussian integral, one can reconstruct all moments by taking derivatives with respect to the source, $J$. Show that you can reconstruct the second moment of the Gaussian integral, \eqref{eq:Gaussian:trick:diff:quad}, by taking two derivatives with respect to $J$. 
\end{exercise}
This gives us 
\begin{align}
    \int \D{\xi}\; \xi^2
    \e^{\frac{im}{2}\frac{\xi^2}{\Delta t}} 
    = 
    \frac{1}{-im/\Delta t}
    \sqrt{\frac{2\pi i\Delta t}{m}} 
    =\frac{i\Delta t}{m} \mathcal{N}^{-1} \ .
    \label{eq:path:integral:Gaussian:trick:2}
\end{align}

We can insert \eqref{eq:path:integral:Gaussian:trick:1} and \eqref{eq:path:integral:Gaussian:trick:2} into the expression for $\psi(q,t+\Delta t)$, \eqref{eq:psi:path:integral:schrodinger:intermediate}. The ugly factor of $\mathcal N$ conveniently cancels. We are left with:
\begin{align}
    \psi(q,t+\Delta t) &= \psi(q,t) - i \Delta t V(q)\, \psi(q,t) + \frac{i\Delta t}{2m} \psi''(q,t)
\end{align}
We can rearrange terms and recognize the time derivative:
\begin{align}
    \frac{\D{\psi(q,t)}}{\D{t}}
    &= i\left[\frac{1}{2m}\left(\frac{\textnormal{d}}{\D{q}}\right)^2 - V(q)\right]\psi(q,t) \ ,
\end{align}
which we recognize as the Schr\"odinger equation. You can debate whether or not this counts as a `derivation' of the Schr\"odinger equation since we implicitly assumed the Schr\"odinger equation when writing the time evolution operator. However, if you instead took as a starting point the observation that each path is weighted by $\e^{iS[q(t)]}$, then one would recover the Schr\"odinger equation as a consequence. 








\begin{exercise}
Derive the same results by taking a time derivative of the continuous expression \eqref{eq:path:integral:formalism}. The steps are analogous, but this will test if you understand what the continuous expression means.
\end{exercise}


\include{Zchap_LA_20_manifolds}

\include{Zchap_LA_30_GreensFunctions}

%%% for Physics 17
%%% \include{Zchap_LA_ZZ_looseends}

\part{Complex Analysis}

Even though the measurable universe is described by real numbers, one must face the realization that the mathematical framework by which we describe physics appears to beg to be understood with respect to complex numbers. 

\include{Zchap_CA_00_intro}
\include{Zchap_CA_20_kramerskronig}



\part{Harmonic Oscillator}

The harmonic oscillator is the building block of perturbative physics. It may not be strictly true that \emph{everything in nature is a harmonic oscillator}, but it is indeed the case that our models for nature are naturally perturbations on exactly solvable systems whose actions are quadratic in the dynamical variable.



%!TEX root = Physics231.tex


\chapter{Preliminaries}

Before we jump in, you may want to review the classical harmonic oscillator in Chapter~\ref{ch:spring:theory}. In what follows, we attack the harmonic oscillator problem using Green's functions. Following our discussion in that chapter, we will then move to the wave equation as the extension of the harmonic oscillator to spacetime.

\chapter{A Harmonic Green's Function}


\section{The most familiar problem in physics}
% \lecdate{lec~12} 

We return to the problem of solving for Green's functions. Our favorite example---arguably, the \emph{only} example\footnote{Upon generalizing to higher dimensions, curvilinear coordinates... most everything is a harmonic oscillator. Physicists have Harmonic Oscillators in different area codes.}---is the harmonic oscillator. The differential operator is
\begin{align}
  \mathcal O = \left(\frac{d}{dt}\right)^2 + \omega_0^2 \ .
  \label{eq:O:HO}
\end{align}
\begin{example}
The equation of motion for a harmonic oscillator with spring constant $k$ and displacement $f(t)$ from equilibrium is simply the equation $F=m\ddot{f} = -k f(x)$. We have defined $\omega_0^2 = k/m$.
\end{example}
\begin{exercise}
Compare our definition of $\omega_0^2 = k/m$ to what you remember the resonant frequency of an ideal spring to be. What---you do not remember? That is okay, physicists do not have to remember these details: we can re-derive them. What---you do not have time to rederive this? Okay, use dimensional analysis. 
\end{exercise}
\begin{exercise}
One may argue that the harmonic oscillator operator should be defined with an overall minus sign. Life is too short to quibble over signs in these notes. Please confirm that if  $\mathcal O_\text{minus} = -\mathcal O$, then the Green's function is simply $G_\text{minus} = - G$. Relate this to the physical interpretation of the Green's function equation. {Hint:} if you pluck a guitar string in the opposite direction, you still get a `unit' wave.
\end{exercise}
The Green's function equation tells us that the Green's function $G(t,t_0)$ is the spring `response' at time $t$ to a `unit displacement' at $t_0$:
\begin{align}
  G''(t,t_0) + \omega_0^2 G(t,t_0) = \delta(t-t_0) \ .
  \label{eq:HO:Greens:eqn}
\end{align}
Recall that the arguments $t$ and $t_0$ are analogous to the indices of a finite-dimensional matrix. We are primarily concerned about the $t$-dependence of $G(t,t_0)$.  For notational convenience, we will set $t_0=0$ and not list it explicitly. What this means is that we are choosing a `zero' for how we measure time: $t=0$ corresponds to the moment of the `unit displacement' source. Positive times correspond to times after the unit displacement. Negative times correspond to times before the unit displacement.

\begin{exercise}
Our system and the laws of physics in this model are time translation invariant. It does not matter where we pick the zero of time. This tells us that we can always make the replacement $t\to t-t_0$ to recover the source time. In other words, as long as the system is invariant along translations of the coordinates $x$, we have $G(x,y) = G(x-y)$. Convince yourself that this is true. Consider what kinds of scenarios would break this relation.
\end{exercise}

\section{Fourier Transform}

The first thing to do is to write $G(t,t_0)$ as a Fourier transform with respect to $t$. Please refer to Appendix~\ref{app:Fourier} for our set of Fourier transform conventions. We can write $G(t)$ as an integral over Fourier modes with frequency $\omega$ and weight (Fourier transform) $\tilde G(\omega)$:
\begin{align}
  G(t) &= \int_{-\infty}^\infty\dbar \omega \, e^{-i\omega t} \tilde G(\omega) 
  &
  \dbar = \frac{d}{2\pi}
  \ .
  \label{eq:HO:Greens:Fourier}
\end{align}
We say that $\tilde G(\omega)$ is the Fourier transform of $G(t)$. The key point is that the $t$-dependence of $G(t)$ has been sequestered into the $e^{-i\omega t}$ plane waves. This is convenient since these plane waves are eigenfunctions of the derivative operator:
\begin{align}
  \frac{d}{dt} e^{-i\omega t} &= -i\omega e^{-i\omega t} \ .
\end{align}
The left-hand side of the Green's function equation \eqref{eq:HO:Greens:eqn} is
\begin{align}
  \mathcal O_t G(t,t_0) 
  &= 
  -
  \int_{-\infty}^\infty \dbar \omega \, 
  \left(\omega^2-\omega_0^2\right) e^{-i\omega t} \tilde G(\omega,t_0) \ .
\end{align}
The right-hand side is simply the Fourier transform of $\delta(t-t_0)$:
\begin{align}
  \delta(t-t_0)
  &=
  \int_{-\infty}^\infty \dbar \omega \, e^{-i\omega (t-t_0)} \ .
  \label{eq:delta:fourier}
\end{align}
\begin{exercise}
Use our conventions for the Fourier transform \eqref{eq:HO:Greens:Fourier} (see also Appendix~\ref{app:Fourier}) to confirm the Fourier representation of $\delta(t-t_0)$ in \eqref{eq:delta:fourier}. In our notation, the Fourier coefficients $\tilde f(\omega)$ of a function $f(t)$ is
\begin{align}
  \tilde f(\omega) &= 
  % \frac{1}{2\pi}
  \int_{-\infty}^\infty d t\, e^{i\omega t} f(t) \ .
\end{align}
\end{exercise}
So the Green's function equation for the 1D harmonic oscillator, \eqref{eq:HO:Greens:eqn}, tells us
\begin{align}
  -
  \int_{-\infty}^\infty \dbar \omega \, 
  \left(\omega^2-\omega_0^2\right) e^{-i\omega t} \tilde G(\omega)
  &=
  \int_{-\infty}^\infty \dbar \omega \, e^{-i\omega t}
  \ .
  \label{eq:G:HO:Fourier:equation:integrals}
\end{align}
For simplicity we have set $t_0=0$ and don't write it explicitly. There's a rather unscrupulous\footnote{I don't think this is rigorously valid, but the result is true. The ends don't justify the means, but let's take this morally ambiguous shortcut to make the big picture clear. I encourage you to live the rest of your lives with virtue.} way to solve this equation for $\tilde G(\omega)$. Since the two sides of this expression are equal, they have the same Fourier expansion. This implies that the Fourier coefficients are equal. Since both sides are already written as Fourier expansions, we can just match the coefficients of the basis functions, $e^{-i\omega t}$. This gives us:
\begin{align}
  \tilde G(\omega) &= \frac{-1}{\omega^2-\omega_0^2}
  \label{eq:G:HO:Fourier:term}
\end{align}
\begin{exercise}
Prove \eqref{eq:G:HO:Fourier:term} honestly. {Hint}: Start with \eqref{eq:G:HO:Fourier:equation:integrals} and project out the Fourier coefficients. Recall that you do this by taking the inner product with one of the basis functions and then using the orthogonality of the eigenbasis. You may need to be careful with the normalization.
\end{exercise}
\begin{exercise}
In \eqref{eq:G:HO:Fourier:equation:integrals} we had already set $t_0 = 0$. What is the expression for $\tilde G(\omega)$ if we kept $t_0$ explicit?
\label{eq:ex:Gtilde:with:t0:explicit}
\end{exercise}
That was the critical step: we have successfully solved for the Green's function Fourier coefficient. This means that we have a closed form expression for the Green's function by plugging $\tilde G(\omega)$ into \eqref{eq:HO:Greens:Fourier}:
\begin{align}
  G(t) &=  \int_{-\infty}^\infty \dbar \omega
  \, 
  \frac{-e^{-i\omega t}}{\omega^2-\omega_0^2} \ .
  \label{eq:G:HO:Fourier:Rep:t0:0}
\end{align}
All that's left is for us to actually \emph{do} this integral. Fortunately, this integral should look very similar. It seems to beg for us to solve using the residue theorem. 
\begin{exercise}
What are the poles of the integrand in \eqref{eq:G:HO:Fourier:Rep:t0:0}? What are their associated residues? What is the residue if $t_0\neq 0$?
\end{exercise}

\subsection{Contour Integral}

The integral \eqref{eq:G:HO:Fourier:Rep:t0:0} looks like it's perfect for contour integration. There's an exponential factor on top that will determine the convergence, and the denominator can be factored to see where the poles are. Except we notice something troubling:
\begin{align}
  \frac{-e^{-i\omega t}}{\omega^2-\omega_0^2} 
  &=
  \frac{-e^{-i\omega t}}{(\omega - \omega_0)(\omega + \omega_0)} \ .
  \label{eq:HO:integrand:real:pole}
\end{align}
The poles are located at $\omega = \pm \omega_0$. These are \emph{on the real axis}, precisely along the integration contour! How annoying!

Now you may want to have an existential moment. Remind yourself that all we're doing is solving for the behavior of the one-dimensional \emph{harmonic oscillator}. This is an eminently \emph{physical} system. We could have written this system as $\mathcal O f(t) = s(t)$ where $f(t)$ is the displacement of a harmonic oscillator and $s(t)$ is some driving function. The Green's function, $G(t,t_0)$ gives the response of the system to a `unit' driving function, $s(t) = \delta(t-t_0)$. The response of the system should be perfectly physical. And yet---\emph{and yet}---we now face an integral \eqref{eq:G:HO:Fourier:Rep:t0:0} that seems to run right into not only one, but \emph{two} singularities along the integration contour!
\begin{center}
\includegraphics[width=.4\textwidth]{figures/Lec_2017_16_HO_Poles.pdf}
\end{center}
Please do take a brief moment to ponder at this existential crisis. This is an apparent failure of our theory---our calculation hits a pole, you get an infinity, and it's not clear what's going on. Rather than anguish, though, we should see if our theory is trying to tell us something. 

In order to figure out what's going on, let's do something wild. The poles on the real axis are in the way. Let's just move them out of the way by giving them small imaginary parts. Each pole can either be pushed above or below the real line:
\begin{center}
\includegraphics[width=.4\textwidth]{figures/Lec_2017_16_HO_PoleShift.png}
\end{center}
To be clear: are \emph{changing} the integral.\footnote{I have altered the poles. Pray that I do not alter them further.} We are exploring what happens when the expression \eqref{eq:G:HO:Fourier:Rep:t0:0} is changed, as if the poles were slightly complex the whole time. In other words, we are \emph{changing} the original differential operator, \eqref{eq:O:HO}. We're not yet motivating \emph{why} we want to do this; we'll justify it post-facto. For now, let's just be happy that our contours no longer bump into the poles. 

\section{Going under the poles?}
Let's try going under the poles, that is: let's push the poles into the upper half-plane with small imaginary parts:
\begin{center}
\includegraphics[width=.6\textwidth]{figures/Lec_2017_16_underpole.png}
\end{center}
This means that we have deformed the integrand \eqref{eq:HO:integrand:real:pole}:
\begin{align}
\frac{-e^{-i\omega t}}{(\omega - \omega_0)(\omega + \omega_0)}
&\to
\frac{-e^{-i\omega t}}{(\omega - \omega_0 - i\varepsilon)(\omega + \omega_0 - i\varepsilon)} \ .
\label{eq:HO:integrand:upper:poles}
\end{align}
To make it clear that we are treating $\omega$ as a complexified variable, let's write it as $z$:
\begin{align}
  \omega \to z \ .
 \end{align}
Now let's go through the exercise of performing the integral along the real line by using the residue theorem. To do this have to choose a contour that includes the real line. We have two options, $C_+$ and $\bar C_-$, where the bar reminds us of the orientation:
\begin{center}
\includegraphics[width=.6\textwidth]{figures/Lec_2017_16_pokeup.pdf}
\end{center}
Which contour do we pick? Does it matter? (Yes!) The contours differ by whether the `large arc' has positive or negative real part. We can parameterize these arcs with respect to their real and imaginary parts:
\begin{align}
  z(\theta) &= R\cos\theta + i R\sin\theta \ ,
\end{align}
where for $C_+$ the arc is given by $\theta \in [0,\pi]$ while for $\bar C_-$ the arc is $\theta\in[0,-\pi]$. Note that the \emph{orientation} matters here---the $\bar C_-$ arc is moving in the clockwise direction.\footnote{Do people reading this even know why it is called `clockwise'? Analog clocks had mechanical arms that rotate.} The component in the imaginary axis is determined by $R\sin\theta$. Because $R$ is real and greater than zero, we have
\begin{align}
  \text{for $C_+$:}\, \sin\theta &> 0
  &
  \text{for $\bar C_-$:}\, \sin\theta &<0 \ .
  \label{eq:HO:which:contour}
\end{align}
%
How do we pick a contour? The criteria for completing the real line into a contour is that the integral along the contour vanishes, that way
\begin{align}
  \oint_C dz\, f(z) 
  = \int_{-\infty}^\infty dx\, f(x) + \lim_{R\to \infty} R \int_\text{arc} d\theta f(Re^{i\theta}) 
  = \int_{-\infty}^\infty dx\, f(x) \ .
\end{align}
And hence one can use the residue theorem to solve the contour integral on the left-hand side and relate it to the real integral on the right-hand side\footnote{I say `real integral' but of course the integrand \eqref{eq:HO:integrand:upper:poles} is now a complex quantity. The point is that we  only integrating along the real line.}. 
%
So which contour satisfies
\begin{align}
  \lim_{R\to \infty} R \int_\text{arc} d\theta 
  % f(Re^{i\theta})  
  \frac{-e^{-i(R\cos\theta + iR\sin\theta) t}}{(Re^{i\theta} - \omega_0 - i\varepsilon)(Re^{i\theta} + \omega_0 - i\varepsilon)}
  \to 0 \ ?
\end{align}
I was a little slick and chose to write $z = R\cos\theta + iR\sin\theta$ in the numerator while writing the more compact $z=Re^{i\theta}$ in the denominator. The convergence\footnote{I tend to say convergence, what I really mean is `vanishes as $R\to \infty$.'} of the integrand depends on the exponential factor in the numerator:
\begin{align}
  e^{-i(R\cos\theta + iR\sin\theta) t} = 
  e^{-iRt\cos\theta} e^{Rt\sin\theta} \ .
\end{align}
The first factor is purely oscillatory and is $\mathcal O(1)$ as $R\to\infty$. The second factor is either exponentially decaying or exponentially growing with large $R$. Convergence depends specifically on this second factor, $e^{-Rt\sin\theta}$. Since $R>0$, this depends on two things: the sign of $\sin\theta$ and the sign of $t$. The integral along the arc vanishes in the $R\to\infty$ limit when:
\begin{align}
  t > 0 \quad\text{and}\quad \sin\theta < 0 &
  &\text{or}&&
  t < 0 \quad\text{and}\quad \sin\theta > 0 \ .
  \label{eq:t:gtr:less:upper:poles}
\end{align}
Or in other words:
\begin{align}
  t>0 &\to \bar C_{-} 
  &
  t<0 &\to C_+ \ .
  \label{eq:t:sign:contour:HO}
\end{align}
% \begin{itemize}
%   \item $t>0$ and $\sin\theta > 0$ or
%   \item $t<0$ and $\sin\theta < 0$ \ .
% \end{itemize}
Recall that we have set $t_0 = 0$; if you did Exercise~\ref{eq:ex:Gtilde:with:t0:explicit}, you would know that restoring $t_0$ corresponds to $t\to t-t_0$ in the conditions above.

The Green's function equation \eqref{eq:Greens:func:as:inverse} is simply that $$\mathcal O_t G(t,t_0) = \delta(t-t_-)\ .$$ This tells us that the Green's function $G(t,t_0)$ is the \emph{response} measured at time $t$ to a $\delta$-function source at $t_0$. Thus the sign of $t-t_0$ corresponds to whether we are measuring the response \emph{before} or \emph{after} the cause. If you are like me, your eyes just lit up with joy. Something very exciting has happened: in this odd discussion about pushing around the poles of $\tilde G$, we have found that we were secretly talking about \emph{causality}.  

Before we progress, let's be clear: we expect that effects happen \emph{after} causes. This means that physical dynamics obey $t-t_0 >0$. A nice theory would tell us that you cannot have an effect before a cause. At the beginning of this course, we talked about locality as one of the motivating factors for why physical dynamics are described by differential operators. This was closely tied to the idea of causality, since non-local effects with spacelike separation can be non-causal in some reference frame. Perhaps amusingly---but beyond the scope of this course or my expertise---the notion of causality and locality in fundamental theories of nature continues to be a theme in cutting edge theoretical research. 

\paragraph{Causal propagation.}
For the specific case of $t>0$ (or $t-t_0>0$ if we restore $t_0$), we observe from \eqref{eq:t:gtr:less:upper:poles} that the integral \eqref{eq:HO:integrand:upper:poles} converges when $\sin\theta <0$. In other words, for times \emph{after} the initial cause at $t_0=0$, the appropriate contour is $\bar C_-$. 
\begin{exercise}
What is the value of the contour integral 
\begin{align}
  \oint_{\bar{C}_-} dz\, \frac{-e^{-i z t}}{(z - \omega_0 - i\varepsilon)(z + \omega_0 - i\varepsilon)} \ ?
\end{align}
Hint: you shouldn't have to do any work.
\end{exercise}
% The contour integral above vanishes along the arc when $t>0$, so it is the correct contour when doing the Fourier transform to get $G(t)$ from $\tilde G(\omega)$. However, 
This contour encloses no poles. No poles means no residues. This means that the contour integral vanishes and the Fourier transform vanishes. In other words, 
\begin{align}
  G(t>0) &= 0 \ .
  \label{eq:HO:G:t:gtr:0:zero}
\end{align}

\begin{example}
The statement that $G(t>0)=0$ should make you sad and possibly upset. This means that there is \emph{no} effect that happens after some cause at $t_0=0$.
\end{example}

\paragraph{Acausal propagation.} By now you see the writing on the wall. Something perverse has happened: after a $\delta$-function disturbance, information is not traveling forward in time because $G(t)=0$ for $t>0$. This, in turn, was self-evident in because for $t>0$, we were required to take the $\bar C_-$ contour in order that the contribution of the arc vanishes. At this point, you must feel the dread of what's coming next. If we take the $t<0$ case in \eqref{eq:t:gtr:less:upper:poles}, we are forced to take the $C_+$ contour. This contour indeed picks up poles and so we expect it to be non-zero. Thus we find that $G(t<0) \neq 0$, which is awful because this $t<0$ corresponds to an \emph{effect that happens before the cause}. 

\begin{exercise}
Taking $t_0=0$ for simplicity, show that the exact result for the acausal Green's function is
\begin{align}
  G(t<0) &= \frac{-1}{\omega_0}\sin\omega_0t \ .
  \label{eq:HO:G:t:less:0:nonzero}
\end{align}
Hint: the relevant residues are $\pm e^{\pm i\omega_0t}/2\omega_0$ \ .
\end{exercise}
You should \emph{do} the above exercise as practice for the main contour integral of this entire course. If you are stuck, you can keep reading these notes.\footnote{By now you should know that learning physics requires \emph{doing} physics. Reading/watching someone else do physics is one of the least efficient ways of improving yourself.} 

To help you along, let me even set up the integral for you. We start with $G(t)$ in \eqref{eq:HO:integrand:real:pole}, where we rewrote the integrand in \eqref{eq:G:HO:Fourier:Rep:t0:0} to make the pole structure clear. We then said that we wanted to push the poles off the real line. This \emph{changes} the physical problem, but we'll worry about what that means later. In this case, we move the poles into the upper half-plane so that each pole has a small imaginary part. The poles are then located at $\pm \omega_0 + i \varepsilon$. The resulting integral is:
\begin{align}
  G(t) &= 
  \int_{-\infty}^\infty \frac{d\omega}{2\pi} 
  \frac{-e^{-i\omega t}}{(\omega - \omega_0-i\varepsilon)(\omega + \omega_0-i\varepsilon)} 
  \ .
\end{align}
Complete the real integral into a contour integral; we'll re-label our variables to remind us that we're in the complex plane, $\omega\to z$. The contour that we select, $C_+$, came from the $e^{-i\omega t}$ factor. The result is:
\begin{align}
  \oint_{C_+}
  \frac{dz}{2\pi} 
  \frac{-e^{-iz t}}{(z - \omega_0-i\varepsilon)(z + \omega_0-i\varepsilon)}
  =&\phantom{+} 
  \int_{-\infty}^\infty \frac{d\omega}{2\pi} 
  \frac{-e^{-i\omega t}}{(\omega - \omega_0-i\varepsilon)(\omega + \omega_0-i\varepsilon)} 
  \nonumber
  \\
  &
  + 
  \cancel{
  \lim_{R\to\infty} \int_{0}^\pi \frac{Rd\theta}{2\pi} 
  \frac{-e^{-i(R\cos\theta + i R\sin\theta)t}}{(Re^{i\theta} - \omega_0-i\varepsilon)(Re^{i\theta} + \omega_0-i\varepsilon)} 
  }
  \label{eq:HO:G:integral:adv:1}
  \\
  &=
  2\pi i \sum_{z_\pm=\pm\omega_0+i\varepsilon} \text{Res}_f(z) \ ,
  \label{eq:HO:G:integral:adv:2}
\end{align}
where $f$ is understood to be the integrand. 
%
From here you should be able to read off the result. Confirm that we can take the $\varepsilon\to 0$ limit; this tiny shift was only there to tell us that the poles live in the upper half-plane.

\paragraph{Advanced propagator.}
The causal ($t>0$) and acausal ($t<0$) cases in \eqref{eq:HO:G:t:less:0:nonzero} and \eqref{eq:HO:G:t:gtr:0:zero} respectively may be combined into a single expression:
\begin{align}
  G_\text{adv}(t) &= \frac{-1}{\omega_0}\sin(\omega_0t) \Theta(-t) \ ,
  \label{eq:HO:G:adv:theta}
\end{align}
where we use the Heaviside step function
\begin{align}
  \Theta(x) = 
  \begin{cases}
  1 & \text{if}\; x>0\\
  0 & \text{if}\; x<0 \ .
  \end{cases}
  \label{eq:Heaviside:Theta}
\end{align}
\begin{exercise}
Check the overall sign. I keep screwing this up. You are henceforth responsible for the overall sign and I absolve myself of any mistakes due to this sign.
\end{exercise}
This is the \textbf{advanced} propagator. Even though `advanced' is often a positive adjective when applied to your progress\footnote{I once read a story about a particle physicist who started a conversation with a stranger on the plane. The stranger explains, ``My son is really quite brilliant. He's studying physics at an Ivy league university.'' The physicist says, ``Oh, how nice. You know, I'm a physicist.''  The stranger then asks what kind of physics, to which the physicist says ``I study elementary particles.'' The stranger responds with, ``Ah yes. Well, my son studies advanced particles.'' I no longer remember where I read this story, but I can attest to meeting people like this.} In this context, however, `advanced' is a \emph{bad} word. It means \emph{acausal}: the cause comes \emph{after} the effect. This does not make any sense, and we reject this as non-sense. 

After an appropriate amount of indignant disdain, we can reflect on how the heck we got here\footnote{If your response is \emph{Christmas! Christmas eve last year}, then I applaud your taste in musical theater.}. We make two observations:
\begin{enumerate}
\item After checking all of our work, the main conclusion from this non-sense result is that our approach of pushing the poles \emph{up} into the upper half-plane must have been wrong. We see that the position of the poles tells us about the causality of the system. By pushing the poles off the real axis, we have imposed a sense of causality in time. We should see what happens when we push the poles \emph{down} into the lower half-plane. 

\item Even though the \emph{physics} of advanced propagator is utter rubbish because it breaks causality, we also appreciate that it is a perfectly reasonable \emph{mathematical} solution. The harmonic oscillator dynamics encoded in our differential operator $\mathcal O_t$ is symmetric with respect to time-reversal symmetry\footnote{The behavior of our models under $t\to -t$ is a powerful diagnostic for understanding our physical theories. One of my friends was in a pub the night before lecture and some of the pub regulars, who knew he was a physicists, asked him what he was lecturing on the next day. He said he was going to talk about reversing the direction of time, to which they cautiously responded: `\emph{... we can do that??}'}. The `physical' interpretation is analogous to seeing a bunch of ripples starting at the edges of a pond that move inward until they reach a singularity ($\delta$-function!?) at which point the pond spits out a small pebble. 
\end{enumerate}


\begin{exercise}
At this point, you should be able to solve the case of the causal Green's function without any additional guidance. In fact, without doing any work, you should see that by pushing the poles down into the lower half plane, the propagator will be causal. With a little bit of thought, you may be able to write the analog of \eqref{eq:HO:G:adv:theta}, the retarded Green's function,
%\footnote{In the same way that `advanced' means something different colloquially versus for physics, we use this word very carefully. One of my friends in graduate school was an international student who liked to make physics puns. At some point she playfully admonished late arriving undergraduates for being `retarded,' which required quite a bit of care to clarify the situation to everyone involved.}, 
without doing any work. You should still proceed to calculate $G_\text{ret}(t)$ the `honest way.' We walk  through it together in the next subsection. 
\label{ex:retarded:G:HO}
\end{exercise}

\section{Going over the poles}

Assuming that you have completed Exercise~\ref{ex:retarded:G:HO}, you have figured out the main result of this course\footnote{Do the exercises. I've come to appreciate how annoying it is to come up with meaningful exercises, rather than rote calculations. At some point there's a phase transition in education where exercises are no longer just `practice' but are pedagogical tools to make sure your brain knows how to \emph{think} physics rather than just \emph{read} physics. To say it differently, as my undergraduate advisor once told me: you should do every single exercise problem you can. If you're short on time, just focus on the ones that you cannot do.}. However, because this example is so important, let's go ahead and walk through it together. We push the poles down into the lower half-plane by giving them a small imaginary part, $-i\varepsilon$.

\begin{center}
\includegraphics[width=.4\textwidth]{figures/Lec_2017_16_pokedown.pdf}
\end{center}

The integral is exactly the same, except the sign of $\varepsilon$ has changed. 
\begin{exercise}
Pop quiz! Unlike the advanced Green's function where we gave the poles a small positive imaginary part, we have now given the poles a small \textbf{negative} imaginary part. How do we assign the contours $C_+$ and $\bar C_-$ to the cases $t>0$ and $t<0$?
\end{exercise}
Did I fool you? The poles may have moved, but the way that we pick contours does not change! The appropriate contour is an unambiguous choice determined by the $e^{-i\omega t}$ factor in the integrand. Thus we \emph{still} have the \emph{same}contour choices in \eqref{eq:t:sign:contour:HO}:
\begin{align}
  t>0 &\to \bar C_{-} 
  &
  t<0 &\to C_+ \ .
\end{align}
What has changed is that now the $\bar C_{-}$ contour encloses poles, whereas the $C_+$ contour encloses no poles. This clearly tells us that the integral over $\omega$ (or its complexificaion $z$) is zero for $t<0$, which means we do not have acausal propagation: 
\begin{align}
  G(t<0) &= 0 \ .
  \label{eq:HO:G:ret:acausal}
\end{align}
Hooray! Signals do not arrive before they are sent.

What about the causal propagation? Remember that $\bar C_{-}$ has a negative orientation. To keep our signs clear, let's write the integral as negative a contour $C_-$  which has the positive orientation: 
\begin{align}
  \oint_{\bar C_-} = -\oint_{C_-} \ .
\end{align}
We can go ahead and do the integral over $C_-$ using the residue theorem. We're simply re-doing \eqref{eq:HO:G:integral:adv:1} and \eqref{eq:HO:G:integral:adv:2} with $\varepsilon\to -\varepsilon$. Writing $f(z)$ to be the integrand, we have:
\begin{align}
  G(t>0) =& 
  -\frac{1}{2\pi} \oint_{C_-}dz\, f(z)
  \\
  =&
  2\pi i \times -\frac{1}{2\pi} 
  \left(\text{Res}_f(\omega_0) + \text{Res}_f(-\omega_0)\right)
  \\
  =& \frac{1}{\omega_0}\sin(\omega_0 t) \ .
\end{align}
Combining this with \eqref{eq:HO:G:ret:acausal} using the Heaviside $\Theta$-function \eqref{eq:Heaviside:Theta} gives the \textbf{retarded Green's function}:
\begin{align}
  G_\text{ret}(t) &= 
  \frac{1}{\omega_0}
  \sin(\omega_0 t)
  \Theta(t) \ .
  \label{eq:HO:Gret:sin:theta}
\end{align}
\begin{exercise}
Restore $t_0\neq 0$ in the above equation for the retarded Green's function.
\end{exercise}
By the way, we will often use the word \textbf{propagator} in place of Green's function. This is because the Green's function propagates information from the source to the observation point\footnote{This is especially clear when solving linear differential equations for which the Green's function represents an exact solution. For non-linear problems, one typically resorts to a perturbation expansion of linear propagators with order-by-order corrections in non-linearity. %This is what Feynman diagrams are: lines represent propagators, which are Green's functions for the linear part of the dynamical equation. The vertices represent non-linearities which must be treated as a Taylor series about a partition function.
}. 

\begin{exercise}
Check the overall sign of the retarded Green's function.  Confirm that the nonzero part of the retarded Green's function ($t>0$) has the opposite sign as the nonzero part of the advanced Green's function ($t<0$). Confirm that this sign difference is due to the orientation of the contour $C_+$ versus $\bar{C}_-$.
\end{exercise}

\section{Other choices? The Feynman Propagator}

We introduced the advanced $G_\text{adv}$ and retarded $G_\text{ret}$ Green's functions. These are both Green's functions for the harmonic oscillator, but for technically different versions of the harmonic oscillator where the differential operator $\mathcal O_t$ has been tweaked to make it slightly complex. We always take the complex part to zero at the end, but we have nudged the poles of the Fourier integrand over a bit and discovered that these have implications on the causality of the Green's function. Clearly the Green's function that we want for most classical applications is the retarded Green's function; this one reflects an effect that happens after the cause.

However, we may as well also explore alternative pole-pushing options. Consider the following mixed prescription:
\begin{center}
\includegraphics[width=.6\textwidth]{figures/Lec_2017_16_Feynman.png}
\end{center}
Here one pole is pushed up and the other is pushed down. We'll only focus on the case of $+\omega_0$ being pushed into the lower half-plane and $-\omega_0$ being pushed into the upper half-plane. I'll leave it to you to figure out the implications of the converse choice.
\begin{example}
Even before we do any of the work, you should identify why this is strange. The exponential factor of the Fourier integrand determines which contour we select. This, in turn, is based on the sign of $t$ and connects to our notion of causality. By splitting the poles like this, we see that there is a non-zero contribution to \emph{both} causal and acausal propagation. What's even more curious, we no longer get the combination of two complex exponentials that gives a (real) sine. Thus the Green's function is manifestly complex. Curious!
\end{example}
We call this the Feynman propagator, $G_F$, which should make you suspect that this choice is not just some mathematical curiosity but actually has physical significance. The result is:
\begin{align}
  G_F(t>0) &= \frac{i}{2\omega_0} e^{-i\omega_0 t}
  &
  G_F(t<0) &= \frac{-i}{2\omega_0} e^{i\omega_0 t} \ .
  \label{eq:HO:G:Feynman}
\end{align}
\begin{exercise}
Prove \eqref{eq:HO:G:Feynman} using the techniques demonstrated in this section.
\end{exercise}
This Green's function turns out to be the correct choice in quantum field theory. Even if you don't know what quantum field theory is, you can guess from the name that this makes sense:
\begin{enumerate}
\item The `quantum' implies that maybe complex numbers aren't so odd. In fact, the $e^{\pm i\omega_0 t}$ looks a lot like $e^{i\hat H t}$, the time-evolution operator. Indeed, frequencies and energies have the same dimension once you throw in the appropriate factors of $\hbar = 1$.
\item The word `field' implies something like a continuum. Continua are represented by a bunch of nearest-neighbor connected springs---in other words, they are naturally represented as coupled harmonic oscillators\footnote{This fact should be better presented in the standard undergraduate curriculum. I recommend Howard Georgi's lectures on waves.}.
\end{enumerate}
What the Feynman propagator/Green's function represents two things, depending on the sign of $t$:
\begin{enumerate}
\item $t>0$: a plane wave moving forward in time with positive energy, $G_F \sim e^{-i\omega t}$. This seems totally fine. The fact that $G_F$ is imaginary is not a big deal, since it now represents something like a wavefunction in quantum mechanics. 
\item $t<0$: a plane wave moving \emph{backward in time} with  with \emph{negative energy}, $G_F \sim e^{-i(-\omega) t}$. This is the suspicious case: note the `double negative' of moving backward in time and having negative energy.
\end{enumerate}
You may know the resolution of the $t<0$ solution since it's a popular story in physics. Dirac found that the solution to the dynamical equation for relativistic electrons contained negative energy electrons moving backward in time. These corresponded to currents of negative charge moving backward in time, which you can hand-wave into understanding as a \emph{positive} charge moving \emph{forward} in time. In other words, rather than describing negative energy electrons moving backward in time, these states correspond to positive-energy positrons moving forward in time. Thus we have a theoretical invitation to the existence of anti-particles.

\begin{flipcomment}
For those interested in particle or nuclear physics, you may also want to notice that the particle and the anti-particle propagators differ by complex conjugation. This is related to the idea that charges reflect some complex nature of the particles. The geometric description of electrodynamics is a $U(1)$ fibration of spacetime. What this means is that every field is allowed to have a complex phase: the $U(1)$ refers to rephasing symmetry, $z\to e^{i\phi} z$. The other familiar forces of fundamental physics are described by the generalization to what is called $SU(N)$ where an $N$-vector of particles $\Psi^i$ transforms under the $N\times N$ special unitary matrix $U^i_{\phantom i j}=\exp(i\phi^a T^a)^i_{\phantom i j}$ so that $\Psi^i \to U^i_{\phantom i j} \Psi^j$. The matrices $T^a$ are called the `generators' of these $N$-dimensional rephasings. For the case $SU(2)$, the $T^a$ are simply the three Pauli matrices, $\sigma^a$ and the $N$-vectors are essentially two-component spinors. The anti-particles are Hermitian conjugates of the particles $\Psi$.
%\label{footnote:U1:SUN}. 
\end{flipcomment}

\section{Feynman Green's Function and the Convergence of the Path Integral}

This section is purely `for culture'. We present a qualitative discussion of how one could have anticipated the Feynman propagator from the perspective of the path integral formulation of quantum mechanics. While path integrals are outside the scope of our course, the overall story here is helpful to see how a theory predicts an equation of motion. The Fourier components of the Feynman Green's function are:
\begin{align}
  \tilde G(k) &=
  \frac{-1}{
  \left[k-(\omega_0 -i\varepsilon)\right]
  \left[k-(-\omega_0 + i\varepsilon)\right]
  }
  =
  \frac{-1}{
  k^2 - \omega_0^2 + 2i\varepsilon \omega_0
  }
  \ .
  \label{eq:Feynman:G:tilde}
\end{align}
Define the small imaginary part to be $2i\epsilon \omega_0 \equiv 2i\delta$. This is just a small, dimensionful quantity.\footnote{A `small dimensionful quantity' implicitly makes reference to another dimensionful quantity against which `smallness' is measured. In this case, it is the resonant frequency $\omega$. Ultimately, the `small' thing is the dimensionless $\varepsilon$.}  This is what we would expect from a Lagrangian:
\begin{align}
  L = \frac{1}{2}\dot q^2 - \frac{1}{2}\left(\omega_0^2 - 2i \delta\right)q^2 \ .
\end{align}
Note: this replacement should be `obvious': we just take $\omega_0^2 \to \omega_0^2 - 2i\delta$.
\begin{exercise}
Check that the equation of motion from this Lagrangian indeed gives a differential operator whose Green's function has Fourier components \eqref{eq:Feynman:G:tilde}.
\end{exercise}
Since the the Lagrangian is related to the action, $S$, by $S=\int dt\, \mathcal L$, one can integrate the $\dot x^2$ term by parts to write:
\begin{align}
  S = -\int dt \, \frac{1}{2}
  q 
  \left[\left(\frac{d}{dt}\right)^2 + \omega_0^2 - 2i\delta \right]
  q 
  =
  -\int dt \, \frac{1}{2}
  q 
  \left[\mathcal O_\text{HO} - 2i\delta \right]
  q  \ ,
  \label{eq:HO:S:Feynman:quad}
\end{align}
where $\mathcal O_\text{HO}$ is the original harmonic oscillator differential operator that we started with in \eqref{eq:O:HO}. This shows us that the Feynman prescription for pushing poles around corresponds to adding a small imaginary part to the harmonic oscillator Lagrangian. This is a little weird since the resulting operator is clearly not Hermitian.

If you're familiar with statistical mechanics or quantum field theory, you may know that the fundamental object is neither the Lagrangian nor the action\footnote{... and is definitely not the Hamiltonian; sorry condensed matter folks and Lin-Manuel Miranda fans.}, but actually the path integral, $Z$ which is defined as an integral over all possible functions $x(t)$:
\begin{align}
  Z &= 
  \int \mathcal Dq(t) \, e^{iS[q]}
  = 
  \int \mathcal Dq(t) \, e^{iS_\text{HO}[q] - \delta q^2} \ .
\end{align}
Here $\mathcal Dx$ is just the measure of the path integral. For discretized space, it corresponds to $\mathcal Dx\sim dq(t_0)\,dq(t_1)\, dq(t_2)\cdots$. The harmonic oscillator action $S_\text{HO}$ is real and thus $e^{iS_\text{HO}}$ oscillates and the path integral doesn't converge. By adding the $- \frac{\delta}{2} q^2$ term to the action we are imposing a small perturbation that forces the integrand to converge for very large values of $x(t)$. With this factor the path integral is well defined. If we approached the harmonic oscillator assuming that path integral is sufficiently convergent\footnote{Physically this is saying that these extreme path configurations should not contribute since the path that minimizes the action should be far away from these cases; so we're just helping the theory focus on the part of the path integral that matters.} and plugging in some factor to force this convergence, then we can run this argument backward and `deduce' the Feynman prescription for the Green's function.

If everything in this subsection is a bit boring, then here's what you should know. This discussion gives us a hint of where all of our physical differential equations come from. The equations of motion that define a theory come from a variational principle with respect to some action. 
\begin{exercise}
[optional] What is the action (Lagrangian, Hamiltonian, whatever...) that yields Maxwell's equations? Note that curiously, you only get half of Maxwell's equations from an action principle. The other half comes from ideas in Section~\ref{sec:stokes:theorem:aside}. For a discussion by a well-known \acro{UCR} mathematical physicist, see \emph{Gauge Fields, Knots And Gravity} by Prof.~John Baez.
\end{exercise}
We then recall that our actions (Lagrangians, Hamiltonians, ...) tend to be \emph{local} function(al)s, so it makes sense that they can be written as a derivative expansion. The actions depend on particles/fields, say\footnote{I'm deliberately changing notation from $x(t)\to q(t)$. This is important because we will eventually want to treat $x$ and $t$ as spacetime coordinates rather than $x(t)$ as `the position of \emph{the} particle at time $t$.' This unfortunately notational confusion tends to be one of the conceptual hurdles in field theory.} $q(t)$. If we write out the leading terms of the action with respect to $q(t)$ you get something like
\begin{align}
  S[q] &= \int dt \, (\text{const.}) + \mathcal O_1 q(t) + \mathcal O_2 q^2(t) + \mathcal O_3 q^2(t) + \cdots
\end{align}
where the $\mathcal O_i$ are differential operators. We can ignore the constant. Usually the $\mathcal O_1$ term doesn't show up, so let's ignore that too; you can remove it by redefining the particle\footnote{The linear term tells you that you're probably not defining the `zero' of your field sensibly. Amusingly, this is closely related to the Higgs phenomenon that gives known fundamental particles most of their mass.}. The $\mathcal O_2 q^2$ term can be integrated by parts. The details of the integration by parts depend on the details of the derivatives in $\mathcal O_2$, but one can write typically write this term as
\begin{align}
  \mathcal O_2 q^2(t) &\equiv - q(t) \mathcal O_\text{quad} q(t) \ ,
\end{align}
which is in the form that we used in \eqref{eq:HO:S:Feynman:quad}. When we derive the equation of motion by varying the action, we end up with
\begin{align}
  -\mathcal O_\text{quad} q(t) = 0 \ ,
\end{align}
where I am probably assuming that $\mathcal O_\text{quad}$ is Hermitian if we want to be technical. The main observation here is that this is a \emph{linear} differential equation. Linear differential equations in physics come from the quadratic pieces of actions. Since we argued that the linear piece is usually not there and we know that the constant piece corresponds to a shift in the vacuum energy (that's \emph{usually} not physical\footnote{Damn you, dark energy.}), then this quadratic piece of the action is the \emph{lowest order} contribution to the action in an expansion of the particle $q(t)$ with respect to the vacuum ($q(t)=0$). The higher-order terms correspond to differential equations that are \emph{not} linear since the variation of that term contains $q^2(t)$. So yes, even though the mathematical structures here are a bit more fancy, let us state out loud that all of our main differential equations in physics basically boil down to taking the leading term in a Taylor expansion. Typical for physicists, eh?

\section{Example: The Damped Harmonic Oscillator}
 
 At this point we are happy that the physical solution to the classical harmonic oscillator is the retarded Green's function, $G_\text{ret}(t)$. There is an unphysical solution, the advanced Green's function, $G_\text{adv}(t)$ that reflects a mathematical symmetry of the theory (time-reversal invariance) but is otherwise not of physical interest. Finally, there is a solution that makes sense for quantum physics, the Feynman Green's function, $G_\text{F}(t)$, which is complex but reflects the existence of anti-particles. Each of the above cases came from \emph{almost} the same harmonic oscillator differential operator, $\mathcal O$, except the operators were slightly deformed so that its poles were off of the real line in the complex momentum plane. We took the limit where this deformation goes to zero, but not until after we've exploited those deformations to do the Fourier integral using the residue theorem. We came to appreciate that the infinitesimal nudge in the imaginary direction was telling us something about the causality of the harmonic oscillator.

 We can understand the nature of this imaginary push by looking at the \emph{damped harmonic oscillator}. The dynamical equation for the damped harmonic oscillator is
 \begin{align}
  \mathcal O_\text{DHO} x(t) = \ddot{x}(t) + 2\gamma\dot x(t) + \omega_0^2 x(t) = F(t) \ ,
 \end{align}
 where $\gamma$ is the \textbf{damping coefficient} and $F(t)$ is a driving force.
 \begin{exercise}
 What are the dimensions of the damping coefficient, $\gamma$?
 \end{exercise}
 
 \paragraph{Back of the envelope version.}
 Before doing anything quantitative, we can guess the behavior. A non-zero $\gamma$ suppresses the response $x(t)$ to the source $F(t)$. Unlike an idealized instrument, a guitar string has non-zero damping and the vibrations get smaller the longer you wait after plucking it. If we imagine a plane wave\footnote{I keep calling $e^{-i\omega t}$ a plane wave because I'm thinking of the generalization to spacetime. For now there are no transverse directions, so of course $e^{-i\omega t}$ is really just a wave.} $x(t)\sim e^{-i\omega t}$ and ask what this wave looks like far away from the source, $F(t)\to 0$, then the differential equation reads:
 \begin{align}
  -\left(\omega^2 +2i\gamma \omega - \omega_0^2\right)e^{-i\omega t} &= 0 \ .
 \end{align}
 The solution for the plane wave's frequency $\omega$ is
 \begin{align}
  \omega = \omega_R - i\gamma \ , \label{eq:DHO:back:envelope:poles}
 \end{align}
 where for $\gamma^2 \ll \omega_0^2$, $\omega_R \approx \omega_0$ is the frequency.  Plugging $\omega$ back into our plane wave gives
 \begin{align}
  x(t) = e^{-i\omega_R t}e^{-\gamma t} \ .
 \end{align}
 The first factor is the usual plane wave. We recognize the $\exp(-\gamma t)$ as an exponential suppression that causes the plane wave to get rapidly smaller for times larger than $1/\gamma$. There are a few lessons:
 \begin{enumerate}
 \item In \eqref{eq:DHO:back:envelope:poles} we have implicitly but decisively bought into the idea of a complex frequency. The interpretation of this complex frequency is that the real part is the usual oscillation and the imaginary part corresponds to damping.

 \item The solution \eqref{eq:DHO:back:envelope:poles} corresponds to finding the poles of the Fourier integrand; recall, for example, \eqref{eq:G:HO:Fourier:equation:integrals}. So damping has something to do with the imaginary part of the poles in the complex plane.

 \item Because the position of the pole in the imaginary direction is related to the \emph{causality} of the theory, we have now connected `damping' to causality in a round-about way.

 \item When we derived the retarded Green's function $G_\text{ret}(t)$ that describes the classical, ideal harmonic oscillator, we took $i\varepsilon\to 0$. For the damped harmonic oscillator the corresponding value is $i\gamma$, which we keep finite\footnote{Here's a funny bit of linguistics in our discipline. Any physical quantity is naturally not-infinite, so it seems silly when physicists make a point to call something `finite.' You would think that this is in contrast to that quantity being infinite. In fact, when physicists say \emph{finite} what they really mean is \emph{not zero}, which is quite different from `not infinite.'}\footnote{Actually, there's a second aside I'd like to make in the footnotes. Jim Holt has an excellent essay, ``The Dangerous Idea of the Infinitesimal'' about how mathematicians struggled with the idea of infinitesimal quantities. The essay is reproduced in his volume \emph{When Einstein Walked with G\"odel}.}; it has some physical significance after all.

 \item The exponential notation for the Fourier transform made all of this very easy: we just let $\omega$ be complex in $e^{-i\omega t}$ and we are able to capture the physics of oscillations and damping. If we had used trigonometric functions then we'd have a bunch of $\sinh$ or $\cosh$ floating around and the convergence would not be so clear.
 \end{enumerate}
 
 \paragraph{More carefully.} Let's go through the steps to derive the Green's function for the damped harmonic oscillator properly. The Green's function equation is
 \begin{align}
  \left[
  \left(\frac{d}{dt}\right)^2
  + 2\gamma \frac{d}{dt}
  + \omega_0^2
  \right]
  G(t)
  &=
  \delta(t) \ ,
 \end{align}
 where we've again set $t_0 = 0$. It is easy at this stage to restore\footnote{There is a theorem that $G(t,t')$ only depends on $|t-t'|$ under reasonable assumptions. You should be able to argue this on physical grounds.} $t_0$ by taking $t\to (t-t_0)$. Plugging in the Fourier representations for the $t$-dependence,
 \begin{align}
  G(t) &= \int \dbar \omega \, e^{-i\omega t}\tilde G(\omega)
  &
  \delta(t) &= \int \dbar \omega e^{-i\omega t} \ ,
 \end{align}
 then gives:
 \begin{align}
  \tilde G(\omega) &=
  \frac{-1}{\omega^2+2i\gamma\omega -\omega_0^2} \ .
 \end{align}
 The poles occur at 
 \begin{align}
  \omega_\pm &= \frac{-2i\gamma \pm \sqrt{-4\gamma + 4\omega_{0^2}}}{2}
  = \pm \sqrt{\omega_0^2 - \gamma^2} - i \gamma \equiv \pm \omega_R - i\gamma\ .
  \label{eq:DHO:poles}
 \end{align}
 \begin{example}
 The poles in \eqref{eq:DHO:poles} have a real part as long as $\omega_0^2 > \gamma^2$. What does it mean when $\omega_0^2 < \gamma^2$? This is the observation that the damping is stronger than the characteristic oscillation frequency, $\omega_0$. Thus the system is properly describe by an exponentially suppression because it never has a chance to oscillate.\footnote{This is one of the magical things about $B$ mesons. $B$ mesons are really nice systems whose characteristic oscillation frequency (between $B$ and anti-$B$ mesons) and characteristic decay width are roughly of the same order of magnitude. When we produce $B$ mesons at, say, the LHCb experiment, we get to observe the decays (imaginary part) of these mesons after $\sim 1$ oscillation which gives a powerful tool for studying matter--anti-matter asymmetry.\label{footnote:B:meson}}
 \end{example}
 Assuming $\omega_0^2 > \gamma^2$, then we may identify $\omega_R = +\sqrt{\omega_0^2 - \gamma^2}$ and the poles may be written as $\pm \omega_R - i\gamma$. This looks a \emph{just} like the pole prescription for the retarded Green's function of the ideal harmonic oscillator with $\omega_0\to \omega_R$ and $i\varepsilon\to i\gamma$. In fact, we can just follow all of the same steps that we did for the retarded Green's function:
 \begin{enumerate}
 \item The poles are in the lower half-plane of $\omega$. This means that the propagator is automatically causal (it's the retarded, not the advanced Green's function).
 \item The exponential factor $e^{-i\omega t}$ in the Fourier integral determines the convergence of the integrand and hence how we complete the contour. For $t>0$ we close the contour on the lower half-plane and pick up these poles in the residue theorem. For $t<0$ we close the contour on the upper half-plane; since there are no poles here, the Fourier integral vanishes.
 \item We end up with our own version of \eqref{eq:HO:Gret:sin:theta} with $\omega_0\to\omega_R$ except we need to account for the finite $\gamma$. 
 \end{enumerate}
 The only substantial difference between the damped harmonic oscillator and the ordinary harmonic oscillator is that $\omega_\pm$ contains a non-negligible imaginary part. The contour integral takes the form:
 \begin{align}
  \oint_C dz \,
  \frac{-e^{-izt}}{
  \left(z-\omega_+\right)
  \left(z-\omega_-\right)
  }
  =&
  \phantom+ 
  2\pi i\left(\text{Res}_f\omega_+ + \text{Res}_f\omega_-\right)
  \\
  =&
  -2\pi i\left[
  \frac{e^{i\omega_+t}}{\omega_+-\omega_-}
  +
  \frac{e^{i\omega_-t}}{\omega_--\omega_+}
  \right]
  \\
  =&
  -2\pi i
  \frac{e^{i\omega_+t} - e^{i\omega_-t}}{\omega_+-\omega_-}
  \\
  =&
  -2\pi i \frac{e^{-i\gamma t}}{2\omega_R}
  \left(e^{i\omega_Rt} - e^{-i\omega_Rt}\right)
  \\
  =&
  \pi \frac{e^{-i\gamma t}}{2\omega_R}
  \sin(\omega t)
   \ .
 \end{align}
 Plugging this into the full solution---which involves keeping track of which contour integral is non-zero and where the factors of $2\pi$ belong---we find:
 \begin{align}
  G(t) = \frac{e^{-\gamma t}}{\sqrt{\omega_0-\gamma^2}} \sin\left(\sqrt{\omega_0^2-\gamma^2} t\right) \Theta(t) \ .
  \label{eq:DHO:G:sin:theta}
 \end{align} 
 \begin{exercise}
Compare and contrast \eqref{eq:DHO:G:sin:theta} to \eqref{eq:HO:Gret:sin:theta}. Why is the general form the same? How do you explain the differences? What is the behavior for $\omega \gg \gamma$, $\omega \ll \gamma$? What does it look like if you restored $t_0$?
 \end{exercise}
 The key observation that we have from this discussion is that the choice of pushing the poles into the lower half plane---which we found was critical for a causal Green's function---is actually physically motivated. Any \emph{physical} oscillating system actually has some small amount of friction (or the appropriate generalization) that causes it to not only oscillate, but also decay. 
 \begin{exercise}
 In quantum mechanics the time evolution operator is $e^{i\hat H t}$. Maybe there's a minus sign in there, I don't remember and I don't really care. It is important that this operator is unitary, or else probabilities are not preserved in time. However, there are plenty of physical systems (say, the $B$ meson in footnote~\ref{footnote:B:meson}) that \emph{are} described my non-Hermitian Hamiltonians, i.e.~non-unitary time evolution operators. Reason tells us that the form of the time-evolution operator should be $e^{i\hat H t}e^{-\gamma t}$, since probabilities less than one make more sense than probabilities bigger than one. What is the physical interpretation of this non-unitary time evolution? Hint: this has to do with the notion that our theory (our Hamiltonian) models something in nature, but it not model \emph{all} of the relevant pieces nature for a given problem.
 \end{exercise}
 At this point in my note I have the following written:
 \begin{quote}
 \emph{I took 5 weeks of physics 231 and all I have is this one stupid Green's function.}
 \end{quote}
 
 
 
\chapter{The Harmonic Oscillator in More Dimensions}



In most graduate courses, we examine simple examples where we can focus on understanding the physics. We then leave it to you to generalize to the horrible real-world scenarios that you will encounter in your other graduate courses and research. There is one generalization that we should take time to do properly together: the shift from one dimension to multiple dimensions. This is the shift from working with functions of time to functions \emph{spacetime}. Equivalently, the harmonic oscillator converts into a wave.
%
Keep in mind that we are not doing any special relativity, even though borrow notation from special relativity\footnote{There was once a professor teaching electrodynamics who used notation informed by relativity. Some of the condensed matter students complained the class is about electricity and magnetism, not relativity. The professor replied: just where do you think magnetism comes from?}. 




\section{The `harmonic oscillator' in spacetime}
The second derivative in higher-dimensional Euclidean space is the Laplacian, $\nabla^2= \partial_x^2 + \partial_y^2 +\cdots$. But when you combine space and time (Minkowski space), there's the famous relative minus sign\footnote{At this level the minus sign is a convenient convention, but we know that the second derivative should really be relativistically invariant: $\partial^2 = \partial_\mu \partial^\mu$.}:
\begin{align}
  \frac{d^2}{dt^2}
  \to 
  \frac{1}{c^2}
  \frac{\partial}{\partial t}^2
  -
  \frac{\partial^2}
  {\partial \vec x^2} \ ,
\end{align}
where $\partial/\partial\vec x = \nabla$ is the gradient. The overall sign is a convention. Against all of my hard-developed instincts for natural units, I have replaced the explicit value of $c$ so that the operator makes sense dimensionally: each term has powers of inverse length squared. In what follows I am likely to make mistakes with the factors of $c$.\footnote{As my adviser used to say: ``If you think there may be a sign error or a factor of two error, then your homework is to fix those errors.'' In this case, if you think there's a missing factor of $c=1$, I am not even sure if I would even acknowledge that it is actually an error.} 

In (1+1)-dimensional of spacetime\footnote{The (1+1) notation means one dimension of space, one dimension of time.} the second derivative is $\partial^2 = c^{-2} \partial_t^2 - \partial_x^2$. That looks familiar, doesn't it? The minus sign tells us that if we move forward in time a little bit, but `look backward' in space, then the function does not change. The description of the state at time $t-\delta t$ and position $x-\delta x$, subject to $\delta x = c\delta t$, is the same as the state at time $t$ and position $x$. 
%
This means that the information was propagated \emph{forward} in space as time also moves forward. This is exactly what we expect from a wave. We recall that the solution to second derivative differential equations are usually trigonometric functions or their exponential counterparts. We further recall that `plane waves' are described with the same funny minus sign:
\begin{align}
  f(t) \sim \sin (\omega t - kx) \ .
  \label{eq:plane:wave}
\end{align}
From this you can read off that the plane wave velocity is $\omega/k$. Of course, you already knew that from dimensional analysis. 

\begin{exercise}
This is a great place to reflect on phase versus group velocity. I should probably write such a discussion into a future version of this notes. \flip{To do.} In the meanwhile, one of my favorite puzzles is the following (perhaps apocryphal) story.\footnote{I think I first heard this story from Tony Zee's lectures at the African School for Theoretical Physics in 2004. He posed this as an \emph{obvious} problem in dimensional analysis. I do not think I properly appreciated the `obvious' aspect until about 15 years later while preparing this course. It was not---and still is not---obvious to me at all.} During the second world war, an admiral wanted to be able to use reconnaissance photographs of ships to determine their velocities. The admiral knew that given data about the ships' positions and velocities, one could extrapolate where those ships are going and when they will arrive. Because ships leave a wake behind them, the admiral figured that a clever physicist should be able to determine a ship's velocity based on the angle of the wake behind the boat.

In deep water, this is called a Kelvin wake. You also see this when watching ducks on the river.\footnote{Not everyone has a chance to watch battleships from a spy plane, but it turns out plenty of universities are situated near rivers where generations of the mathematically inclined could reflect on physical phenomena. One could just sit by the water and think. I suspect this is the reason why Cambridge has such a historically strong fluid dynamics department. Anyway, I used to watch the ducks swimming in the stream near my old Ph.D institution. There were also geese around. I would not watch the geese. I suspect they also leave a Kelvin wake, but what I know with certainty is that geese leave behind gigantic goose shits that make it unpleasant to sit by the water.} The key `aha' moment for understanding the Kelvin wake is that the relevant model is one where the key restoring force comes from gravity---these wakes are formed by \emph{gravity} waves. Do not confuse this with \emph{gravitational waves}, which are quite different. By neglecting surface tension, we find that the phase velocity is $\sqrt{g/k}$, where $k$ is the wave number ($e^{ikx}$). The group velocity (that carries energy) is $d\omega/dk$, where $\omega=\sqrt{gk}$ is the angular frequency. Thus the group velocity is half the phase velocity. 

Using a geometric argument, show that the angle of the Kelvin wake is $\sin^{-1}(1/3)$. The book by Stone \& Goldbart has a nice discussion.\footnote{See also \url{www.itp.uni-hannover.de/fileadmin/itp/emeritus/zawischa/static_html/KWake.html}} Even before jumping into the construction, you should wonder where the ratio 1:3 shows up in the problem. 
\end{exercise}
% A good exercise: phase and group velocity; see appell or cahill


\section{Green's functions for multidimensional spaces}
Our `harmonic oscillator' equation becomes:
\begin{align}
  \left[\frac{1}{c^2}
      \frac{\partial}{\partial t}^2
      -
      \frac{\partial^2}
      {\partial x^2}
      -
      \frac{\partial^2}
      {\partial y^2}
      -
      \frac{\partial^2}
      {\partial z^2}
    \right]
    \varphi(\vec{x},t) = \rho(\vec{x},t) \ .
    \label{eq:phi:wave:eq}
\end{align}
Huh, that looks familiar, doesn't it? We've chosen variables so that this looks just like the wave equation for the scalar potential $\varphi$ subject to a charged source $\rho$ in electrodynamics! In fact, you know how this works. There are three more equations that correspond to the vector potential:
\begin{align}
  \left[\frac{1}{c^2}
      \frac{\partial}{\partial t}^2
      -
      \frac{\partial^2}
      {\partial x^2}
      -
      \frac{\partial^2}
      {\partial y^2}
      -
      \frac{\partial^2}
      {\partial z^2}
    \right]
    \vec A(\vec{x},t) = \vec j(\vec{x},t) \ .
    \label{eq:A:wave:eq}
\end{align}
Now let's go through a series of questions about \eqref{eq:phi:wave:eq} and \eqref{eq:A:wave:eq} to make sure we're on the same page. 
\begin{enumerate}
\item \textbf{Are these differential equations still linear?} Yes. Remember what it means to be linear! $\mathcal O(f+g) = \mathcal Of +\mathcal Og$, for example.
\item \textbf{Are we worried that there are multiple arguments?} Not really. The functions now depend on $t$ and $\vec{x}=(x,y,z)$. You went from an infinite dimensional `vector space' to a still-infinite dimensional vector space\footnote{You can pontificate about whether this infinity is `bigger.' It doesn't really make a difference.}. You can also mumble reassuring words to yourself, perhaps recalling how in prior encounters with the wave equation you've perhaps separated your functions into single-argument factors, $\Psi(t,x) = T(t)X(x)$, or something like that.
\item \textbf{How many equations are there?} There are \emph{four}\footnote{Obligatory TNG reference: \url{https://www.youtube.com/watch?v=wjKQQpPVifY}} equations. There is an equation for $\phi$ and one equation for each component of $\vec A$. In fact, we typically bundle this all up into a four-vector $A_\mu=(\varphi, \vec A)$.\footnote{For now this is just a bundling, so you may argue that it's not obviously a \emph{vector}. The fact that $\varphi$ and $\vec{A}$ transform into each other upon performing a boost is something we learn from physical observations. Alternatively, the identity of $A_\mu$ as a vector (technically, a one-form) is imposed by the mathematical structure---but that mathematical structure was adopted after the empirical progress in the 1800s.}
\item \textbf{How many differential operators are there?} Just one. There is only \emph{one} differential operator, $\partial^2$. It acts on different components of $A_\mu$, but it's the same wave operator acting on each component. 
\item \textbf{... so how many equations are there, really?} All four equations are essentially the same equation with different state functions and different sources. So it's really just one class of differential equation.
\end{enumerate}
Okay, here's the really important one. I'm going to put it in a box to make sure you pay really close attention. Please answer the following question before reading on:
\begin{flipcomment}
\centering
How many Green's functions are there?
\end{flipcomment}
Do we need a Green's function for each component of $A_\mu$? Do we need a Green's function for each dependent variable, $x^\mu = (ct,x,y,z)$? \emph{No}! There is only \emph{one} differential equation, and thus there is only \emph{one} Green's function. 

Let's see how this works. Suppose you have a differential equation in multiple variables,
\begin{align}
  \mathcal O_{\vec{x}} f(\vec{x}) = s(\vec{x}) \ .
\end{align}
Let's say that $\vec{x}=(x_1, \cdots, x_N)$, so that we're working in $N$-dimensional space. Suppose you are given the Green's function $G(\vec{x},\vec{x}')$ for $\mathcal O_{\vec{x}}$. This means that $G$ satisfies the Green's function equation
\begin{align}
  \mathcal O_{\vec{x}} G(\vec{x},\vec{x}') = \delta^{(N)}(\vec{x}-\vec{x}') \ ,
\end{align}
where the $N$-dimensional $\delta$-function is simply the product of one-dimensional $\delta$-functions in the expected way:
\begin{align}
  \delta^{(N)}(\vec{x}-\vec{x}') 
  &=
  \delta(x_1-x_1')\delta(x_2-x_2')\cdots\delta(x_N-x_N') \ .
\end{align}
The solution to the above differential equation is
\begin{align}
  f(\vec{x}) &= \int d^N\vec{x}'\,  G(\vec{x},\vec{x}') s(\vec{x}') \ ,
\end{align}
which you can check by acting on both sides with $\mathcal O_{\vec{x}}$.




\section{Relativistic Notation}

Let's borrow some four-vector notation from special relativity. Nothing that we're doing is based on special relativity, but the fact that it pops out for free is telling. Our indices $\mu$ run over four values which we take to be $\mu = 0,1,2,3$ corresponding to timelike (0) and spacelike (1,2,3) directions. Some objects are naturally defined with an upper index, like $x^\mu = (ct,x,y,z)$. You can lower these indices using the metric, $\eta_{\mu\nu}=\text{diag}(1,-1,-1,-1)$. This means that for some vector $v^\mu$,
\begin{align}
  v^2 \equiv v^\mu v_\mu 
  = c^2(v^0)^2 - (v^1)^2 - (v^2)^2 - (v^3)^2 \ .
 \end{align}
 More importantly, this means that the four-vector of partial derivatives, which is naturally a lower-index object, squares to:
 \begin{align}
   \partial^2\equiv
  \partial_\mu \partial^\mu 
  &= 
  \frac{1}{c^2}
      \frac{\partial}{\partial t}^2
      -
      \frac{\partial^2}
      {\partial x^2}
      -
      \frac{\partial^2}
      {\partial y^2}
      -
      \frac{\partial^2}
      {\partial z^2} \ .
 \end{align}
The placement of the $c^2$ factors should be clear from dimensional analysis. We see that indeed the wave operator is the natural generalization of the harmonic oscillator in spacetime. We won't actually be raising and lowering indices, but we want to be clear that there's a natural `dot product' that includes the relative minus sign between spatial and temporal components.

The states and sources in our system are the (four-)vector potential $A_\mu$ and the current $j_\mu$: 
\begin{align}
  A_\mu &= (\varphi, \vec{A})
  &
  j_\mu &= (\rho, \vec{j}) \ .
\end{align}
The four equations for the vector potential are
\begin{align}
  \partial^2 A_\mu = j_\mu \ .
\end{align}
That's one equation for each value of $\mu$, but each equation has the \emph{same} differential operator $\partial^2$. Our task is to find the Green's function for this operator in spacetime. The Green's function equation is
\begin{align}
  \partial^2 G(x,x') &= \delta^{(4)}(x-x') \ ,
  \label{eq:wave:eq:Greens:eq}
\end{align}
where we write $x$ and $x'$ to mean \emph{four}-vectors with \emph{components} $x^\mu$ and $(x')^\mu$. Given a current $j_\mu(x)$, the vector potential is the solution to the wave equation which integrates over all of the current (source) positions:
\begin{align}
  A_\mu(x) &= \int d^4x' \, G(x,x') j_\mu(x) \ .
\end{align}
Observe that $G$ has two arguments, $x$ and $x'$, but no indices. The two arguments correspond to the observation point $x$ and the source point $x'$. The integral over $d^4x'$ means we're accounting for sources from all positions and all times\footnote{This should remind you of starlight that is emitted a long time ago ($t'$) in a galaxy far, far away ($\vec{x}'$ that we we observe today $(t,x)$.}. We have not specified our choice of coordinates. For now you can assume that we are using Cartesian coordinates, but you may recall that the Green's function for the Poisson equation\footnote{Which we can now call the `harmonic oscillator' in (3+0)-dimensional spacetime.} is most cleanly expressed in spherical coordinates, $G(r)\sim (4\pi r)^{-1}$. You should expect that the natural way to write $G(x,x')$ in (3+1)-dimensional spacetime is some hybrid of spherical and Cartesian coordinates.

\begin{exercise}\label{ex:guess:Greens}
At this point you may want to start going over Homework \#4 where \emph{you} solve the (3+1)- and (2+1)-dimensional Green's functions in detail. I'm quite proud of that homework assignment, there's a lot of neat observations. One observation is that with some physical insight, you could take the Poisson equation's Green function and \emph{guess} the Green's function for the wave equation simply by using the finite speed of light. This is discussed nicely in Tony Zee's latest textbook, \emph{Fly By Night Physics}. 
\end{exercise}


We solve \eqref{eq:wave:eq:Greens:eq} in what is now the `usual' way. We Fourier decompose both sides. Making deliberate choices for our Fourier transform conventions, let us write
\begin{align}
  f(t,\vec{x}) &= \int \dbar^4k \, e^{-ik\cdot x} \tilde f(\omega, \vec{k})
  &
  \dbar^4k &= \frac{dE\, d^3\vec{k}}{(2\pi)^4}
   \ .
\end{align}
\begin{exercise}
Please confirm that with our metric convention, this choice of (3+1)-dimensional Fourier conventions matches the (0+1)-dimensional Fourier conventions we used for the harmonic oscillator.
\end{exercise}
For simplicity we'll just write $f(x)$ and $\tilde f(k)$ to refer to the four-component arguments\footnote{Later on we'll use $k=|\vec{k}|$ to mean the spatial length of the 3-momentum; there should be no ambiguity because the context should be clear. This is an example of `read what I mean, not what I [literally] write.' This is a common `chalkboard physics' peccadillo; usually it's a pain to write \emph{literally} what you mean. Sometimes it's better to write the `easy to interpret' equation, then tell people to interpret it slightly differently. Mathematicians should be incensed. That's okay. This is mathematics for physicists, not physics for mathematicians. The latter course is also quite interesting and theorists should take it if it's ever offered from John Baez.}. 
We have used the Minkowski spacetime inner product,
\begin{align}
  k_\mu x^\mu \equiv k\cdot x = Et - \vec{k}\cdot\vec{x} \ .
\end{align}
From the perspective of the Fourier transform, this is just a sign convention for each Cartesian direction that we are free to make. From the perspective of physics, we know that it is convenient to pick a notation that is manifestly Lorentz invariant\footnote{If we had chosen a different notation, the theory would still be Lorentz invariant---its built into the operator---but it would not be \emph{obviously} Lorentz invariant. This would be like using cylindrical coordinates and losing track about whether or not space is rotationally symmetric.} and automatically looks like a traveling wave. By the way, if you're keeping track of factors of $c$:
\begin{align}
  x^\mu &= (ct,\vec{x})
  &
  k^\mu &= (E/c,\vec{p})
  &
  k_\mu &= (E/c, -\vec{p})
  &
  x\cdot k  &= Et - \vec{p}\cdot\vec{x} \ .
\end{align}
I hate keeping track of factors of $c$. 


Here's what the Fourier transform of the spacetime $\delta$-function looks like:
\begin{align}
  \delta^{(4)}(x-x') &=
  \int \dbar E e^{-iE(t-t')}
  % \times
  \int \dbar k_x e^{+ik_x(x-x')}\
  % \times
  \int \dbar k_y e^{+ik_y(y-y')}\
  % \times
  \int \dbar k_z e^{+ik_z(z-z')} \ .
\end{align}
Again, we have \emph{chosen} the sign of the Fourier transform so that the time-like component has a minus sign, $e^{-iE(t-t')}$, while the space-like component has a plus sign, $e^{+ik_x(x-x')}$. This is convenient since we could package all of these exponential together with respect to $k_\mu = (E, -k_x, -k_y, -k_z)$,
\begin{align}
  \delta^{(4)}(x-x') &=
  \int \dbar^4k e^{-ik\cdot (x-x')} \ .
  \label{eq:4D:delta:Fourier}
\end{align}
Conveniently, this looks just like a---\emph{wait for it}---plane wave:
\begin{align}
  e^{-ik\cdot x} = e^{-i\left(Et - \vec{k}\cdot\vec{x}\right)} \ .
\end{align}
Compare this to \eqref{eq:plane:wave}. By the way, if you're really picky you can bring up factors of $c$ again because you argue that $\delta(t-t')$ should really be $\delta(ct-ct')$. This is worth checking, but I'd rather relegate it to yet another parenthetical footnote\footnote{
  If we keep factors of $c$ explicit, then we recall that $t\to ct$ and $E\to E/c$. The expression $\delta(t-t')$ is replaced by $\delta(ct-ct') = \delta(t-t')/c$. The factor of $1/c$ shows up in the corresponding replacement $dE \to dE/c$. To solve for the Fourier transform $\tilde G$, we compare Fourier coefficients of the $\delta$ function with the Fourier coefficients of the Green's function. In that comparison, these factors of $c$ cancel out. Thus we are free to simply write $\delta(t-t')$ as long as our integration measure is $dE$ and not $dE/c$. If you're ever lost, you should just set $c=1$ and then replace factors of $c$ at the end to make sure the dimensions of $G$ work out.
}.

The left-hand side of the Green's function equation \eqref{eq:wave:eq:Greens:eq} is 
\begin{align}
  \partial^2 \int \dbar^4k \, e^{-ik\cdot x} \tilde G(k,x')
  &=
  \int \dbar^4k \, \left(-\bar{E}^2+\vec{k}^2\right) e^{-ik\cdot x} \tilde G(k,x')
  &
  \bar{E}\equiv E/c
   \ .
  \label{eq:4D:Greens:wave:LHS:Fourier}
\end{align}
We have introduced the $\bar E$ notation to keep track of the factors of $c$ coming from $\partial^2 \sim c^{-2}\partial_t^2 -\cdots$. If you'd had set $c=1$, you could have guessed this because you want to make sure that both terms in the denominator have the same dimensions.
Equating the Fourier coefficients of \eqref{eq:4D:delta:Fourier} with \eqref{eq:4D:Greens:wave:LHS:Fourier} gives us the Fourier modes:
\begin{align}
  \tilde G(k,x') &= \frac{-e^{ikx'}}{\bar{E}^2 - \vec{k}^2} \ .
\end{align}
Thus the wave equation Green's function in (3+1)-dimensions is:
\begin{align}
  G(x,x') &= \int \dbar^4k \, \frac{-e^{ik\cdot (x-x')}}{\bar{E}^2 - \vec{k}^2} \ .
  \label{eq:Greens:function:For:4D:int}
\end{align}
This looks almost trivial---you've integrated something with this \emph{exact} integrand before when we did the harmonic oscillator. The only issue is that now the $\dbar\omega$ has been replaced with a much more\footnote{\emph{The more integrals we come across, the more problems we see}, with apologies to the Notorious \acro{BIG}.} ominous $\dbar^4k$. 

\section{Anticipating the Answer}

We very deliberately choose to motivate this problem from the perspective of electromagnetic waves. It is good training---in fact, this may be the \emph{primary} training---to build an intuition for what you expect the Green's function to be. What you already know?
\begin{enumerate}
  \item You know that the Green's function to the wave equation is the response to a unit source. This source is localized in both space and time: it is a blip of infinitesimal extend at an instant. 
  \item You know that the electromagnetic wave propagates at a finite speed, $c$. 
  \item You know that in the \emph{static} limit, the Green's function is $\Phi(\vec{x}-\vec{x}')= 1/(4\pi |\vec{x}-\vec{x}'|)$. 
\end{enumerate}
The last point is the most intriguing one. A static source is almost like the unit blip for the wave equation, except that it keeps ``blipping'' continuously. Rather than a flash of light, it is a lightbulb that is simply on. In other words, if we knew the Green's function for the wave equation, we could integrate it over a source $s(t,\vec{x}) = \delta^{(3)}(\vec{x}) u(t)$ where the function $u(t) =$ constant. This should give us the Green's function in electrostatics,
\begin{align}
  \int d^{(4)}x'\; G(t-t',\vec{x}-\vec{x}') s(t', \vec{x})' = \frac{1}{4\pi |\vec{x}-\vec{x}'|} \ . 
  \label{eq:guessing:wave:Greens:function}
\end{align}
We have used the observation that $G(x,x') = G(x-x')$. 
From this we can \emph{guess} that the wave equation Green's function is
\begin{align}
  G(t-t',\vec{x}-\vec{x}') &= 
  \frac{\delta(|\vec{x}-\vec{x}'|-ct)}{4\pi |\vec{x}-\vec{x}'|} \ ,
\end{align}
up to factors of $c$ that I tend to be sloppy about. You can check that plugging this into the electrostatic equation \eqref{eq:guessing:wave:Greens:function} with the static source $s(t,\vec{x})$ indeed recovers the solution to the Poisson equation. 

The interpretation of this Green's function is clear. One feels a $1/r$ potential, but information about this potential only propagates at a finite velocity. That means an observer at $\vec{x}$ and time $t$ feels the potential coming from the origin that was emitted at time $t-c|\vec{x}|$. In hindsight, you could argue that the wave equation Green's function \emph{had} to be this.  


 


\section{Sketch of the Fourier integral in (3+1)-dimensions}

Rather than solve this in gory detail, we'll just provide the framework and then we will it to you to dot all the $i$s and cross all the $d$s.\footnote{The aphorism is ``...cross all the $t$s.'' I just made a $\dbar = d/2\pi$ joke. I like to think that footnotes like this contributed to me getting tenure.}  It'll be a little tedious the first time you do do this integral, but don't worry. If you drink alcohol, I recommend a glass of wine\footnote{Or the appropriate substance relevant for you personally. I had a mentor who suggested a glass of red wine before doing a tedious calculation. I knew another physicist who was renowned for doing complicated calculations in quantum field theory who once told me: \emph{every day I drink a glass of red wine... and that is why I will live forever!} I think with enough red wine my calculations would take me forever to complete. There is a great quote from \emph{Calvin \& Hobbes}: ``I was put on this earth to accomplish a certain number of things. Right now I’m so far behind I will never die.'' So there you have it: a model for how it is that red wine may extend one's \sout{life} Ph.D.} before starting.

The solution is most clear if we use \textbf{hypercylindrical} coordinates that correspond to spherical coordinates over the spatial directions and a Cartesian coordinate in the time direction. The integration measure is
\begin{align}
  \dbar^4k &= \frac{1}{(2\pi)^4} dE \, |\vec{k}|^2 d|\vec{k}|\, d\cos\theta \, d\varphi \ .
  \label{eq:4D:momentum:measure}
\end{align}
The polar angle $\theta$ is defined by 
\begin{align}
  (\vec{x}-\vec{x}')\cdot \vec{k} = |\vec{x}-\vec{x}'||\vec {k}| \cos\theta \ .
\end{align}

\paragraph{The azimuthal integral.}
The azimuthal angle $\varphi$ is defined with respect to the azimuthal direction set by $\vec{x}-\vec{x}'$. Observe that the integrand of the Fourier integral \eqref{eq:Greens:function:For:4D:int} is completely independent of $\varphi$ so we might as well do the $d\varphi$ integral first:
\begin{align}
  \int_0^{2\pi} \dbar \varphi &= 2\pi \ .
\end{align}
This conveniently cancels a factor of $(2\pi)$ in the denominator of \eqref{eq:4D:momentum:measure}.

\paragraph{The polar integral.}
We can tidy up our remaining expressions by defining $c=\cos\theta$ and $k = |\vec{k}| = \sqrt{k_x^2 + k_y^2+k_z^2}$ for the radial direction in momentum space; hopefully this doesn't cause any ambiguity with $k = (E/c,k_x,k_y,k_z)$ the four-vector\footnote{It will not cause any ambiguity because you'll notice that all of the factors of $c$ have been removed in the Fourier integrand \eqref{eq:Greens:function:For:4D:int}.}. Similarly, let's write $r$ to be the separation distance between the observer and the source, $r=|\vec{x}-\vec{x}'|$. Thus we have
\begin{align}
  \int_{-1}^1 d\cos\theta \, e^{i\vec{k}\cdot(\vec{x}-\vec{x'})}
  &=
  \int_{-1}^1 dc\, e^{ikrc}
  = 
  \frac{1}{ikr}\left(e^{ikr} - e^{-ikr}\right) \ .
\end{align}
Plugging this into \eqref{eq:Greens:function:For:4D:int} gives
\begin{align}
  G(x,x') &=
  \frac{1}{(2\pi)^3}
  \int dk \, dE \, k^2
  \, 
  \frac{1}{ikr}\left(e^{ikr} - e^{-ikr}\right)
  \frac{- e^{-iE (t-t')}}{\bar{E}^2-k^2}
  \\
  &=
  \frac{1}{8\pi^3 ir}
  \int_0^\infty dk\, 
  k \left(e^{ikr} - e^{-ikr}\right)
  \int_{-\infty}^\infty dE \, 
  \frac{-e^{-iE (t-t')}}{\bar{E}^2-k^2}
  \ .
\end{align}
For simplicity we have renamed $|\vec k| \to k$. Note the limits of the integrals. Because $k$ is a radial variable, it only takes non-negative values. The energy integral, on the other hand, is an honest Fourier transform over all real values.

\paragraph{The energy integral.} Ah! The $dE$ integral is an old friend: you remember it when it used to go by the name $d\omega$ for the harmonic oscillator. You know how to do this one: just follow what we did to derive the retarded Green's function for the harmonic oscillator. Nudge the poles infinitesimally into the lower-half plane\footnote{This is as good a time as any to point out that the rule of nudging the poles into the \emph{lower} half plane for the causal propagator is directly related to our Fourier integral convention. Some references use the opposite convention and thus nudge the poles into the upper half plane.}. The solution is
\begin{align}
   \frac{1}{2\pi}\int_{-\infty}^\infty dE \, 
  \frac{-e^{-iE \Delta t}}{\bar{E}^2-k^2}
  &= \frac{-i}{2k}\left(e^{ikc\Delta t} - e^{-ikc\Delta t}\right)\Theta(\Delta t) \ .
\end{align}
Note that we have been careful to replace the factors of the speed $c$ when integrating since we remember $\bar E \equiv E/c$. This is straightforward from writing $E = c\bar E$ and then changing integration variables $dE = cd\bar{E}$. You could have also guessed this because you know the right-hand side will have exponentials with arguments that go like $k$ and $\Delta t$. In order for these to make sense dimensionally, there must be a factor of $c$ in there\footnote{If you're being strict about units, recall that the `energies' are really frequencies and the `three-momenta' are really wave numbers. Those who like to think in natural units will be confused why this footnote is necessary.}.

Compare this to \eqref{eq:HO:Gret:sin:theta}. We have chosen not to write the difference of exponentials into sines; we will see why shortly. Otherwise, the three-momentum magnitude $k$ plays the role of the frequency $\omega_0$. We went ahead and wrote $\Delta t \equiv t-t'$; in the harmonic oscillator we equivalently said\footnote{Technically we said $t_0=0$ because that subscript made sense for an `initial time,' but now we use a more standard Green's function convention that $x'$ is the spacetime position of the source.} ``set $t'=0$''. 

\paragraph{The [magnitude of the] three-momentum integral.} The remaining integral over $k=|\vec k|$ is
\begin{align}
  G(x,x')&= -
  \frac{\Theta(\Delta t)}{8\pi^2 r} 
  \int^\infty_0
  dk\, 
  \left[
    e^{ik(r+c\Delta t)}
    - e^{ik(r-c\Delta t)}
    - e^{ik(-r+c\Delta t)}
    + e^{ik(-r-c\Delta t)}
  \right] 
  \\
  &= -
  \frac{\Theta(\Delta t)}{8\pi^2 r} 
  \int^\infty_0
  dk\, 
  \left[
    e^{ik(r+c\Delta t)}
    - e^{ik(r-c\Delta t)}
    - e^{i(-k)(r-c\Delta t)}
    + e^{i(-k)(r+c\Delta t)}
  \right] 
  \\
  &= -
  \frac{\Theta(\Delta t)}{8\pi^2 r} 
  \int^\infty_{-\infty}
  dk\, 
  \left[
    e^{ik(r+c\Delta t)}
    - e^{ik(r-c\Delta t)}
  \right] 
  \ .
\end{align}
In the second line we were slick and realized that we could convert the lat two terms into the first two terms if we extended the region of integration to include negative values of $k$. Pat yourself on the back if you caught this simplification. The integrands are now pure exponentials. We know what these are: they're $\delta$-functions!
\begin{align}
  G(x,x')
  &=
  \frac{\Theta(\Delta t)}{4\pi r} 
  \left[ \delta(r-c\Delta t) - \delta(r+c\Delta t) \right]
  \ .
\end{align}
At this point a bit of either physical or mathematical insight is required. The second term is problematic: $\delta(r+c\Delta t)$. Recall that $r$ is a radial coordinate, which means that $r>0$. However, we also have this $\Theta(\Delta t)$ sitting in the function that imposes that $\Delta t > 0$. That means that $\delta(r+c\Delta t) = 0$ for any interesting values of $r$ and $t$ away from the original `$\delta$-function' source\footnote{Such a source is not at all physical and is just a model for a physical source whose details are not resolved by our model. For example: in Newtonian mechanics, when we say that an ideal particle has an ideal elastic collision with an ideal wall, we say that the momentum changes from $\vec{p}$ to $-\vec{p}$ and that this happens instantaneously at the moment the particle hits the wall. This is a discontinuous jump in the momentum, which means that the force, $\vec{F} = \dot{\vec{p}}$ must be infinite. We don't actually believe that there is an infinite force, but the model doesn't care about the details of how the ideal wall is actually a little squishy. The details of that \emph{microphysics} are only a negligible contribution to the model of, say, a basketball hitting the backboard. This is the idea that physical models as effective theories with domains of validity.}. The non-physical term apparently encodes acausal propagation, so we'll go ahead and drop it like its hot. Formally this is because we have the $\Theta(\Delta t)$ forcing the integral over $\delta(r+c\Delta t)$ to be zero.

The result for the retarded Green's function in spacetime is remarkably simple
\begin{align}
  G_\text{ret} = \frac{1}{4\pi r} \Theta(\Delta t) \delta(r-c\Delta t) \ .
\end{align}
Take a moment to appreciate this. The first factor is \emph{exactly} the Green's function for the Poisson operator, $\nabla^2$, in three spatial dimensions. The $\Theta$ function imposes forward-propagation in time (retarded versus advanced). Then there's a $\delta$-function whose sole purpose seems to impose that $r = c\Delta t$. A moment's thought reveals the interpretation. Recall that the Green's function is the solution to an idealized $\delta$-function source at spacetime point $x'$: $\partial^2 G = \delta^4(x-x')$. The $\delta(r-c\Delta t)$ is telling me that the response from my initial idealized $\delta^4$-function source propagates outward (increasing $r$) at speed $c$ along a \emph{wave front}. At time $t$ the wave front reaches radial distance $c(t-t')$. In other words, this is just reminding us about propagation speed in the wave equation. In other words, the speed of electrodynamic waves (light) is constant\footnote{We can be snotty and say ``you mean, the speed of light... \emph{in vacuum},'' but this is really beside the point. We started with Maxwell's equations. We didn't even specify the medium. Maybe it's uniform water. The point is that the \emph{wave equation} told us that the speed of propagation is constant.}. Here's what the propagation of a `blip' of light looks like (this is no surprise to you):
\begin{center}
\includegraphics[width=.9\textwidth]{figures/Lec27_3dblip.png}
\end{center}

With a little bit of thought\footnote{This addresses the point in Exercise~\ref{ex:guess:Greens}.}, you could have \emph{guessed} down this Green's function for the wave equation based on the Green's function from the Poisson equation. All you have to do is to remember that the speed of light is constant and you probably want to mathematically enforce causality. In this way, the Poisson equation, the harmonic oscillator, and the wave equation are all \emph{the same damn thing}. They're the appropriate second derivative in a given (3+0)-, (0+1)-, and (3+1)-dimensional spacetime.

To pontificate further on this: all of the uglier versions of the Poisson equation and the wave equation in cylindrical and spherical coordinates are just cousins of what we've done here. Sure, you may have to replace your Fourier expansion over nice plane waves with spherical harmonics or Bessel functions. But you know the deal: we never cared much about the basis functions, we just cared about using the \emph{orthonormality} of these \emph{nice} basis functions to read off the coefficients ($\tilde G$) so that it's easy to represent the solution. Keep this in mind and you won't get lost when you get to the mathematically more tedious examples in your electrodynamics course---you're doing nothing more and nothing less than the harmonic oscillator in some number of dimensions and in some choice of basis.

\begin{exercise}
Solve for the Green's function of the wave equation in (2+1)-dimensional spacetime. Notice that the behavior is quite different from the one in (3+1)-dimensions. Comment on the physics of waves in \emph{Flatland}.
\emph{Answer}:
\begin{center}
\includegraphics[width=.47\textwidth]{figures/Lec27_2da}
\includegraphics[width=.47\textwidth]{figures/Lec27_2db}
\end{center}
\end{exercise}


\section{Dimensional Reduction}

You can guess that the (2+1)-dimensional Green's function is a different beast. In (3+1) dimensions, we have $\nabla^2 (-1/4\pi r) = \delta^{(3)}(\vec{x})$. On the other hand, in (2+1) dimensions we have $\nabla^2 (\ln r)/(2\pi) = \delta^{(2)}(x)$, where the $\nabla^2$ is the appropriate spatial d'Alembertian in each case.

\begin{exercise}
Do you remember deriving the $\varphi\sim \ln r$ potential somewhere? I think most students must have done this, but perhaps not explicitly in the guise of Green's functions in flatland.
\end{exercise}

\begin{exercise}
What are the dimensions of $G_{(2+1)}$? How do these compare to the dimensions of $G_{(3+1)}$? Keep track of this difference as we go over this section.
\end{exercise}

There is a clever way to deduce the (2+1)-dimensional Green's function. In fact, one could have used this trick to derive the (3+1)-dimensional Green's function from the 3-dimensional Green's function. We propose that
\begin{align}
  G_{(2+1)}(t,x,y) &= \int_{-\infty}^\infty dz\, G_{(3+1)}(t,x,y,z) \ .
\end{align}
That is: we integrate out the extraneous dimension from the higher-dimensional Green's function. To see this, insert the Fourier representation of $G_{(3+1)}$:
\begin{align}
  \int_{-\infty}^\infty dz\, G_{(3+1)}(t,x,y,z)
  &=
  \int dz\,\dbar^4k\, e^{-iEt+i\vec{k}\cdot\vec{x}} \tilde G_{(3+1)}(E,k_x,k_y,k_z) \ .
\end{align}
We can perform the $dz$ integral: $\int dz \,e^{ik_zz} = 2\pi \delta(k_z)$. This gives
\begin{align}
  \int_{-\infty}^\infty dz\, G_{(3+1)}(t,x,y,z)
  &=
  \int \dbar^3k\, e^{-iEt+ik_xx+ik_yy} \tilde G_{(3+1)}(E,k_x,k_y,0) \ .
\end{align}
We can now plug in the form of $\tilde G_{(3+1)}$:
\begin{align}
  \tilde G_{(3+1)}(E,k_x,k_y,0) &=\frac{1}{E^2-k_x^2-k_y^2} = \tilde G_{(2+1)}(E,k_x,k_y) \ .
\end{align}
\flip{check minus signs.} We observe that the (3+1)-dimensional momentum-space Green's function with $k_z=0$ is exactly the (2+1)-dimensional momentum-space Green's function. The difference is how many integrals one performs to go to position space. So we have found that
\begin{align}
  \int dz\, G_{(3+1)} = \int \dbar^3k\, e^{-iEt+ik_xx +ik_y y} \tilde G_{(2+1)}
  \equiv G_{(2+1)} \ .
\end{align}
Thus we have proved our assertion that the (2+1)-dimensional Green's function is an integral of the (3+1)-dimensional Green's function over the `extra' dimension.

We can go ahead and plug in $G_{(3+1)}$ to find
\begin{align}
  G_{(2+1)}(t,x,y) 
  &= \frac{\Theta(t)}{4\pi}\int_{-\infty}^{\infty}dz\,
  \frac{1}{\sqrt{x^2+y^2 + z^2}}
  \delta(\sqrt{x^2+y^2+z^2}-t) \ .
\end{align}
Recall that
\begin{align}
  \delta(f(z)) &= \sum_{z_0} \frac{1}{|f'(z_0)|} \delta(z-z_0) \ ,
\end{align}
where the sum is over zeros of the function, $f(z_0)=0$. For our case we have $f'(z) = z/\sqrt{x^2+y^2+z^2}$ with zeros at $z_\pm = \pm \sqrt{t^2 - x^2+y^2}$. Using the $\Theta(t)$ to remove an unphysical term, we have
\begin{align}
  G_{(2+1)}(t,x,y) &= \frac{\Theta(t)}{2\pi} \frac{1}{\sqrt{t^2-x^2-y^2}} \ .
\end{align}


The following figures give a bit of intuition for this solution relative to an infinite line charge:
\begin{center}
\includegraphics[width=.47\textwidth]{figures/Lec08_XD_2D1}
\includegraphics[width=.47\textwidth]{figures/Lec08_XD_2D2}
\end{center}
The key insight is that $\int dz\, G_{(3+1)}$ is precisely the form of the potential coming from an infinite, uniform line charge extending in the $z$-direction from the origin. The electric field would ordinarily `leak' into the extra dimension, but the infinite line charge cancels the potential that leaks off of a given sheet. Note that by symmetry the non-trivial potential profile is restricted to the plane.



\begin{exercise} Application to biophysics. Derive the position-space two point function in equation (3.17) of \texttt{cond-mat/9601006v2}\footnote{\url{https://arxiv.org/abs/cond-mat/9601006v2}}. One has to perform a two-dimensional (inverse) Fourier transform to show: 
\begin{align}
  \int \dbar^2p 
  \frac{ e^{i \vec{p} \cdot (\vec{x}-\vec{y}) } }{p^4}
  &=
  \frac{1}{16\pi} R^2 \ln R^2 \ ,
\end{align}
where $\vec{R} = \vec{x}-\vec{y}$ and $R = |\vec{R}|$. The dot product in the exponential gives the cosine of an angle. If we move to polar coordinates in momentum space, we get $p\,dp\,d\theta$. Because the two-dimensional measure is $d\theta$ rather than $d\cos\theta$, this integral is somewhat tricky. Derive the above result by performing the analogous integral in \emph{three} dimensions and then dimensionally reducing to two dimensions. 
\emph{Thanks to Yankang Liu (P232 Fall 2022) for bringing up this application.}
\end{exercise}




\include{Zchap_HO_20_QM_path_integral}


\part{Randomness}

One of the unique and valuable features of a physics education is the way that we learn how to deal with randomness. In fact, we see randomness in at least three different ways:
\begin{enumerate}
  \item Noise in experimental data.
  \item The statistical randomness in thermodynamics.
  \item The quantum randomness in quantum mechanics.
\end{enumerate}
The statistical intuition developed from randomness in physics makes our tools especially generalizable---just ask the physicists involved in theoretical work on \acro{AI}.



\include{Zchap_ST_00_Probability}
\chapter{Entropy} 
  % Baez

  You may be familiar with the notion of \textbf{entropy}\index{entropy} from statistical mechanics\sidenote{You probably first heard about entropy in thermodynamics. I encourage you to derive the thermodynamical definition from the statistical definition},
  \begin{align}
    S = -\sum_i p_i \ln p_i \ ,
  \end{align}
  where the sum is over the number of microstate configurations and we have chosen sensible units where $k_\textnormal{B} =1$. 
  \begin{example}
  For the microcanonical ensemble, every allowed microstate is equally probably. That means that if there are $N$ possible microstates consistent with a macroscopic observation like energy conservation, then each microstate $i$ has probability $p_i = 1/N$. The entropy is thus
  \begin{align}
    S = -\frac{1}{N}\sum_i \ln \frac{1}{N} = \ln N \,
  \end{align}
  which is the logarithm of the number of microstates.
  \end{example}
  In statistical physics this is a good working definition: entropy is a measure of the unobserved microstates that are consistent with the observed macroscopic state. The most likely macroscopic state is the one with the most microstates. 

  There is a more poetic way to phrase this idea that turns out to be even more useful. The entropy is a measure of information known about the system. A system with large entropy is one where you know comparatively little: an observed macrostate with large entropy is one where there are many possible microstates and you do not know which microstate the system is in at a given instant. A system with low entropy: say, a chunk of iron in its ferromagnetic phase, is one where you know very well that every microscopic spin is more-or-less aligned in a certain direction. In this way, \emph{entropy is a measure of information}.

  This idea was formalized in what we now refer to as \textbf{information entropy}\index{information entropy} or \textbf{Shannon entropy}\index{Shannon entropy}. It is defined in exactly the same way up to a factor of $\ln 2$:
  \begin{align}
    S = -\sum_i p_i \log_2 p_1 \ ,
  \end{align}
  where the only difference is that the logarithm is now taken with base 2. THis is because in information theory we work with bits of binary information. The notion of information versus `physics' entropy are the same. They are so `the same' that I will use the same symbol for them. 










\section*{Acknowledgments}

\acro{PT}\ thanks the students of Physics 17 (Spring 2022, 2023, and 2024) and Physics 231 (Fall 2016 through 2024) for their feedback and patience.  It was your curiosity, questions, frustrations, and occasional elation that shaped these notes. He also thanks Owen Long for support to create this course and believing in it, Yakov Eliashberg, from whom he first heard the phrase ``\emph{there is no such thing as a position vector},'' and Savas Dimopoulos who showed him the elegance of a well-chosen index convention.
%
% \acro{PT} is supported by the \acro{DOE} grant \acro{DE-SC}/0008541.
\acro{PT} is supported by a \acro{NSF CAREER} award (\#2045333).

%% Appendices
\appendix

\include{Zchap_APP_LagrangeLegendre}
\include{Zchap_APP_Misc}






\include{Zchap_ZZ_ColophonEndnotes.tex}

\printindex

%% Bibliography
%% USING BIBLATEX, SKIP BIBTEX
%% Use inspireHEP bibtex entries when possible
% \bibliographystyle{utcaps} 	% arXiv hyperlinks, preserves caps in title
% \bibliography{bib title without .bib}


\end{document}